% =============================================================================
% APPENDIX: Source Code Listings
% =============================================================================
% Fully self-contained -- works directly on Overleaf.
% Drop this file into your project and % =============================================================================
% APPENDIX: Source Code Listings
% =============================================================================
% Fully self-contained -- works directly on Overleaf.
% Drop this file into your project and % =============================================================================
% APPENDIX: Source Code Listings
% =============================================================================
% Fully self-contained -- works directly on Overleaf.
% Drop this file into your project and % =============================================================================
% APPENDIX: Source Code Listings
% =============================================================================
% Fully self-contained -- works directly on Overleaf.
% Drop this file into your project and \input{appendix_code_listings}
% in your report where you want the appendix to appear.
%
% If you already have these packages in your preamble, comment out the
% \usepackage lines below.
% =============================================================================

% ---- Packages (uncomment any NOT already in your preamble) ------------------
% \usepackage{xcolor}
% \usepackage{listings}
% \usepackage{caption}

% ---- Listing style -----------------------------------------------------------
\definecolor{codebg}{rgb}{0.95,0.95,0.97}
\definecolor{codeframe}{rgb}{0.80,0.80,0.85}
\definecolor{stringgreen}{rgb}{0.15,0.55,0.15}
\definecolor{keywordblue}{rgb}{0.10,0.20,0.70}
\definecolor{commentgray}{rgb}{0.35,0.35,0.35}

\lstdefinestyle{appendix}{
  language        = C++,
  basicstyle      = \footnotesize\ttfamily,
  backgroundcolor = \color{codebg},
  frame           = single,
  rulecolor       = \color{codeframe},
  numbers         = left,
  numberstyle     = \tiny\color{gray},
  stepnumber      = 1,
  numbersep       = 8pt,
  tabsize         = 2,
  showspaces      = false,
  showstringspaces = false,
  breaklines      = true,
  breakatwhitespace = true,
  keywordstyle    = \color{keywordblue}\bfseries,
  commentstyle    = \color{commentgray}\itshape,
  stringstyle     = \color{stringgreen},
  captionpos      = t,
  xleftmargin     = 2em,
  framexleftmargin = 1.5em,
}

\lstdefinestyle{appendix-py}{
  style           = appendix,
  language        = Python,
}

\lstdefinestyle{appendix-ini}{
  style           = appendix,
  language        = {[LaTeX]TeX},
}

% =============================================================================
\chapter{Source Code Listings}  \label{app:code}
% =============================================================================

% -----------------------------------------------------------------------------
\section{Remote Controller Firmware (C++)}
% -----------------------------------------------------------------------------

The handheld remote controller firmware runs on an Adafruit Feather~M0
with an RFM69~433\,MHz transceiver, an SSD1306~128$\times$32~OLED display,
and an analog dual-axis joystick.  PlatformIO is used as the build system.

% ---- main.cpp ---------------------------------------------------------------
\begin{lstlisting}[style=appendix, caption={Remote --- main.cpp}, label={lst:remote-main}]
#include <Arduino.h>
#include "OledDis.h"
#include "RadioTrans.h"
#include "Joystick.h"

#define BUTTON_PIN 10
#define VBATPIN A7
#define PWR_CONTROL 5

volatile uint8_t buttonState = 0;

void handleButton() {
  buttonState = !buttonState;
}

void setup() {
  Serial.begin(115200);
  delay(2000);

  pinMode(PWR_CONTROL, OUTPUT);
  digitalWrite(PWR_CONTROL, HIGH);


  attachInterrupt(digitalPinToInterrupt(BUTTON_PIN), handleButton, FALLING);

  Display_init();

  transiver_init();   

  Serial.println("System ready");
}

void loop() {
  uint8_t x, y;
  int8_t RSSI = 0;
  uint8_t battery = 0;
  char robotName[15] = "No connection";

  readJoystick(&x, &y);


  // Remote battery mesurements
  float measuredvbat = analogRead(VBATPIN);
  measuredvbat *= 2;
  measuredvbat *= 3.3;
  measuredvbat /= 1024.0;



  sendData(x, y, buttonState, &battery, robotName, &RSSI);
  Display_draw(robotName, measuredvbat, battery, RSSI);

  delay(100);
}
\end{lstlisting}

% ---- Joystick.cpp -----------------------------------------------------------
\begin{lstlisting}[style=appendix, caption={Remote --- Joystick.cpp}, label={lst:remote-joystick}]
#include <Arduino.h>


void readJoystick(uint8_t *x, uint8_t *y) {
    *x = -analogRead(A0) >> 2;
    *y = -analogRead(A1) >> 2;
}
\end{lstlisting}

% ---- OledDis.cpp ------------------------------------------------------------
\begin{lstlisting}[style=appendix, caption={Remote --- OledDis.cpp}, label={lst:remote-oled}]
#include <Arduino.h>
#include "OledDis.h"
#include <Wire.h>
#include <Adafruit_SSD1306.h>

#define SCREEN_WIDTH 128
#define SCREEN_HEIGHT 32

Adafruit_SSD1306 display(
  SCREEN_WIDTH,
  SCREEN_HEIGHT,
  &Wire,     
  -1
);

void Display_init() {
  Wire.begin();
  delay(100);
  if (!display.begin(SSD1306_SWITCHCAPVCC, 0x3C)) {
    Serial.println("OLED not found");
    while (1);
  }
  Serial.println("OLED found");

  display.clearDisplay();
  display.setRotation(2);
  display.ssd1306_command(SSD1306_DISPLAYON);
  display.ssd1306_command(SSD1306_NORMALDISPLAY);
  display.display();
}

void Display_test() {
  display.clearDisplay();

  display.drawRect(0, 0, 128, 32, SSD1306_WHITE);

  display.drawLine(0, 0, 127, 31, SSD1306_WHITE);

  display.display();
}

void Display_clear() {
  display.clearDisplay();
  display.display();
}

void Display_draw(char *name, float remoteBattery, uint8_t robotBattery, int8_t RSSI) {
  display.clearDisplay();
  display.setTextSize(1);
  display.setTextColor(SSD1306_WHITE);
  display.setCursor(0, 0);

  display.print("Name:");
  display.println(name);

  display.print("R:");
  display.print(remoteBattery, 1);
  display.print("V");

  display.print("B:");
  display.print(robotBattery / 10.0, 1);
  display.println("V");

  display.print("RSSI:");
  display.println(RSSI);

  display.display();
}
\end{lstlisting}

% ---- RadioTrans.cpp ---------------------------------------------------------
\begin{lstlisting}[style=appendix, caption={Remote --- RadioTrans.cpp}, label={lst:remote-radio}]
#include "RadioTrans.h"
#include <SPI.h>
#include <RFM69.h>
#include <RFM69registers.h>

#define RFM69_CS    0
#define RFM69_RST   1
#define RFM69_INT   A4
#define RFM69_DIO2  A3
#define RFM69_DIO3  A5

#define NETWORKID   42
#define NODEID      2
#define TONODEID    1
#define FREQUENCY   RF69_433MHZ
#define LATENCY_PACKET 0x42
#define PACKET_TEST 0x55

RFM69 radio(RFM69_CS, RFM69_INT);

void transiver_init()
{
    pinMode(RFM69_RST, OUTPUT);
    digitalWrite(RFM69_RST, HIGH);
    delay(50);
    digitalWrite(RFM69_RST, LOW);
    delay(50);

    SPI.begin();

    if (!radio.initialize(FREQUENCY, NODEID, NETWORKID))
    {
        Serial.println("RFM69 init FAILED");
        while (1);
    }

    radio.setHighPower();

    Serial.println("RFM69 init OK");
}

bool sendData(uint8_t x, uint8_t y, uint8_t button, uint8_t *battery, char *robotName, int8_t *RSSI)
{
    uint8_t packet[3];
    packet[0] = x;
    packet[1] = y;
    packet[2] = button;

    if (!radio.sendWithRetry(TONODEID, packet, sizeof(packet), 1, 50))
    {
        Serial.println("No ACK received");
        return false;
    }

    *battery = radio.DATA[0];

    for (int i = 0; i < 14; i++)
    {
        robotName[i] = (char)radio.DATA[i + 1];

        if (robotName[i] == '\0')
            break;
    }

    robotName[14] = '\0';

    *RSSI = radio.RSSI;

    return true;
}
\end{lstlisting}

% ---- Joystick.h -------------------------------------------------------------
\begin{lstlisting}[style=appendix, caption={Remote --- Joystick.h}, label={lst:remote-joystick-h}]
#include <Arduino.h>

void readJoystick(uint8_t *x, uint8_t *y);
\end{lstlisting}

% ---- OledDis.h --------------------------------------------------------------
\begin{lstlisting}[style=appendix, caption={Remote --- OledDis.h}, label={lst:remote-oled-h}]
#pragma once
#include <Arduino.h>

void Display_init();
void Display_test();
void Display_clear();
void Display_draw(char *name, float remoteBattery, uint8_t robotBattery, int8_t RSSI);
\end{lstlisting}

% ---- RadioTrans.h -----------------------------------------------------------
\begin{lstlisting}[style=appendix, caption={Remote --- RadioTrans.h}, label={lst:remote-radio-h}]
#ifndef RFM69_H
#define RFM69_H

#include <Arduino.h>

void transiver_init();

bool sendData(uint8_t x, uint8_t y, uint8_t button, uint8_t *battery, char *robotName, int8_t *RSSI);


#endif
\end{lstlisting}

% ---- platformio.ini ---------------------------------------------------------
\begin{lstlisting}[style=appendix-ini, caption={Remote --- platformio.ini}, label={lst:remote-pio}]
[env:adafruit_feather_m0]
platform = atmelsam
board = adafruit_feather_m0
framework = arduino
monitor_speed = 115200
build_flags = -D USB_SERIAL
lib_deps = 
	adafruit/Adafruit SSD1306@^2.5.16
	lowpowerlab/RFM69@^1.6.0
	lowpowerlab/SPIFlash@^101.1.3
\end{lstlisting}

% -----------------------------------------------------------------------------
\section{Robot-Side Firmware (C++)}
% -----------------------------------------------------------------------------

The robot-side firmware runs on a second Adafruit Feather~M0.  It receives
joystick data over the RFM69 link, reads battery voltage and a human-readable
name from a Raspberry Pi over the serial port, and echoes these back to the
remote inside the radio's ACK payload.

% ---- Robot main.cpp ---------------------------------------------------------
\begin{lstlisting}[style=appendix, caption={Robot --- main.cpp}, label={lst:robot-main}]
#include <Arduino.h>
#include <SPI.h>
#include <RFM69.h>
#include "RemoteTransiver.h"

void setup() {
  Serial.begin(115200);
  while (!Serial) delay(10);

  transiver_init();
}

void loop() {
  transiver_receive();
}
\end{lstlisting}

% ---- RemoteTransiver.cpp ----------------------------------------------------
\begin{lstlisting}[style=appendix, caption={Robot --- RemoteTransiver.cpp}, label={lst:robot-trans}]
#include <SPI.h>
#include <RFM69.h>
#include "RemoteTransiver.h"
#include <RFM69registers.h>

#define RFM69_CS   A5
#define RFM69_INT  6
#define RFM69_RST  A4
#define ESTOP      A3
#define LED_PIN    13


#define NETWORKID  42
#define NODEID     1
#define TONODEID   2
#define FREQUENCY  RF69_433MHZ
#define LATENCY_PACKET 0x42
#define PACKET_TEST 0x55

RFM69 radio(RFM69_CS, RFM69_INT);

float batteryVoltage = 0;       
char robotName[16] = "none/con";     
bool piDataReceived = false;     

unsigned long lastReceiveTime = 0;
unsigned long lastNoConnectionPrint = 0;
const unsigned long CONNECTION_TIMEOUT = 1000;


void transiver_init()
{
    pinMode(LED_PIN, OUTPUT);
    pinMode(RFM69_RST, OUTPUT);
    digitalWrite(RFM69_RST, HIGH);
    delay(100);
    digitalWrite(RFM69_RST, LOW);
    delay(100);

    pinMode(ESTOP, OUTPUT);

    if (!radio.initialize(FREQUENCY, NODEID, NETWORKID))
    {
        Serial.println("RFM69 init FAILED");
        while (1);
    }

    radio.setHighPower();

    Serial.println("RFM69 init OK");
    lastReceiveTime = millis();
}

void readPiSerial()
{
    static char line[64];
    static uint8_t idx = 0;

    while (Serial.available())
    {
        char c = Serial.read();
        if (c == '\r') continue;

        if (c == '\n')
        {
            line[idx] = '\0';
            idx = 0;

            // Find "name:" 
            char* namePtr = strstr(line, "name:");
            char* batPtr  = strstr(line, ",bat:");

            if (namePtr && batPtr)
            {
                namePtr += 5; // skip "name:"

                // Copy name up to the comma
                uint8_t i = 0;
                while (namePtr[i] != ',' && namePtr[i] != '\0' && i < 15)
                {
                    robotName[i] = namePtr[i];
                    i++;
                }
                robotName[i] = '\0';

                // Parse float manually
                batPtr += 5; // skip ",bat:"
                batteryVoltage = atof(batPtr);
                piDataReceived = true;


            }

        }
        else
        {
            if (idx < sizeof(line) - 1)
                line[idx++] = c;
            else
                idx = 0;
        }
    }
}


void handleDataPacket(uint8_t x, uint8_t y, uint8_t button)
{

    Serial.print("x:");
    Serial.print(x);
    Serial.print(", y:");
    Serial.print(y);
    Serial.print(", btn:");
    Serial.println(button);

    if (button == 1)
        digitalWrite(ESTOP, HIGH);
    else
        digitalWrite(ESTOP, LOW);
}


void transiver_receive()
{

    readPiSerial();

    if (radio.receiveDone())
    {
        if (radio.ACKRequested())
        {

            uint8_t ackPayload[16];

            // Battery voltage, e.g. 24.3V -> 243
            ackPayload[0] = (uint8_t)(batteryVoltage * 10.0);

            // Copy robot name into remaining bytes
            uint8_t i = 0;
            for (; i < 14 && robotName[i] != '\0'; i++)
            {
                ackPayload[i + 1] = robotName[i];
            }

            // Zero-fill rest
            for (; i < 14; i++)
            {
                ackPayload[i + 1] = 0;
            }

            radio.sendACK(ackPayload, 15);
        }

        uint8_t x      = radio.DATA[0];
        uint8_t y      = radio.DATA[1];
        uint8_t button = radio.DATA[2];

        handleDataPacket(x, y, button);
        lastReceiveTime = millis();
    }

    if (millis() - lastReceiveTime >= CONNECTION_TIMEOUT) {
        if (millis() - lastNoConnectionPrint >= 1000) {
            Serial.println("no connection"); 
            lastNoConnectionPrint = millis();
        }
    }
}
\end{lstlisting}

% ---- RemoteTransiver.h ------------------------------------------------------
\begin{lstlisting}[style=appendix, caption={Robot --- RemoteTransiver.h}, label={lst:robot-trans-h}]
#ifndef RFM69_H
#define RFM69_H

#include <Arduino.h>

void transiver_init();
void transiver_receive();
void readPiSerial();

#endif
\end{lstlisting}

% ---- Robot platformio.ini ---------------------------------------------------
\begin{lstlisting}[style=appendix-ini, caption={Robot --- platformio.ini}, label={lst:robot-pio}]
[env:adafruit_feather_m0]
platform = atmelsam
board = adafruit_feather_m0
framework = arduino
monitor_speed = 115200
build_flags = -D USB_SERIAL
lib_deps = 
	adafruit/Adafruit SSD1306@^2.5.16
	lowpowerlab/RFM69@^1.6.0
	lowpowerlab/SPIFlash@^101.1.3
\end{lstlisting}

% -----------------------------------------------------------------------------
\section{Test and Measurement Scripts (Python)}
% -----------------------------------------------------------------------------

The following Python scripts were used during the project to characterise
the radio link's latency, packet-loss, and signal-strength (RSSI)
performance.  All measurement scripts communicate with the robot-side
Feather~M0 over a serial port and use \texttt{matplotlib} to produce plots.

% ---- test_feather.py --------------------------------------------------------
\begin{lstlisting}[style=appendix-py, caption={Feather serial-communication test}, label={lst:test-feather}]
#!/usr/bin/env python3
"""
Feather M0 Communication Test Script
============================
Usage: python test_feather.py

This script tests serial communication with the Feather M0 (RobotSide).

Setup:
  - Windows: Uses COM8 (change below if different)
  - Linux: Uses /dev/ttyACM1 (change below if different)

Run:
  python test_feather.py

What it does:
  1. Sends "name:TESTBOT,bat:12.34\n" to Feather
  2. Waits for response
  3. Reports what was received

Output examples:
  - "Response: OK" = Feather received and responded
  - "No response" = Feather not receiving data
"""

import serial
import time
import sys

WINDOWS_PORT = 'COM8'
LINUX_PORT = '/dev/ttyACM1'

if sys.platform.startswith('win'):
    SERIAL_PORT = WINDOWS_PORT
elif sys.platform.startswith('linux'):
    SERIAL_PORT = LINUX_PORT
else:
    SERIAL_PORT = '/dev/ttyUSB0'  # Mac

BAUD_RATE = 115200
TIMEOUT = 2

def test_feather():
    print("=== Feather M0 Communication Test ===")
    print(f"Platform: {sys.platform}")
    print(f"Port: {SERIAL_PORT}")
    print(f"Baud: {BAUD_RATE}")
    print("-" * 40)

    try:
        print("Opening serial port...")
        ser = serial.Serial(SERIAL_PORT, BAUD_RATE, timeout=TIMEOUT)
        time.sleep(2)
        print("Port opened!")

        print("\nSending test data...")
        test_msg = "name:TESTBOT,bat:12.34\n"
        ser.write(test_msg.encode())
        print(f"Sent: {test_msg.strip()}")

        print("\nReading response...")
        received_any = False
        
        for i in range(3):
            time.sleep(0.5)
            if ser.in_waiting > 0:
                data = ser.read(ser.in_waiting)
                print(f"Response {i+1}: {data.decode('utf-8', errors='replace').strip()}")
                received_any = True
            else:
                print(f"Attempt {i+1}: No data available")

        if not received_any:
            print("\nNo response received from Feather!")
            print("This means Feather is not responding to the serial data.")

        ser.close()
        print("\n=== Test Complete ===")

    except serial.SerialException as e:
        print(f"\nERROR: {e}")
        print("\nTroubleshooting:")
        print("  - Make sure Feather is connected")
        print("  - Check correct COM port (Windows: COMx, Linux: /dev/ttyACMx)")
        print("  - Close other programs using the port")
        print("  - Try a different USB port")

if __name__ == "__main__":
    test_feather()
\end{lstlisting}

% ---- BitRateLatencyTest.py --------------------------------------------------
\begin{lstlisting}[style=appendix-py, caption={Latency vs.\ bitrate measurement}, label={lst:bitrate-latency}]
import serial
import time
import matplotlib.pyplot as plt
import csv
import re
from datetime import datetime

SERIAL_PORT = "COM8"
BAUD_RATE = 115200

BITRATE_LABEL = "1.2 kbps"   # Change this manually for each test
MEASURE_TIME = 10

timestamp = datetime.now().strftime("%Y%m%d_%H%M%S")

csv_filename = f"latency_{BITRATE_LABEL.replace(' ', '_')}_{timestamp}.csv"
plot_filename = f"latency_{BITRATE_LABEL.replace(' ', '_')}_{timestamp}.png"

ser = serial.Serial(SERIAL_PORT, BAUD_RATE, timeout=1)
time.sleep(2)

rtt_values = []
oneway_values = []

pattern = r"RTT:\s+(\d+)\s+us\s+\|\s+One-way latency:\s+([\d.]+)\s+ms"

print(f"\nRecording latency data for {BITRATE_LABEL}...")

start = time.time()

while time.time() - start < MEASURE_TIME:
    line = ser.readline().decode(errors='ignore').strip()
    match = re.search(pattern, line)

    if match:
        rtt = int(match.group(1))
        latency = float(match.group(2))
        rtt_values.append(rtt)
        oneway_values.append(latency)
        print(f"RTT: {rtt} us | Latency: {latency} ms")

ser.close()

print(f"\nCollected {len(rtt_values)} samples")

with open(csv_filename, mode='w', newline='') as file:
    writer = csv.writer(file)
    writer.writerow(["Sample", "Bitrate", "RTT_us", "OneWayLatency_ms"])
    for i in range(len(rtt_values)):
        writer.writerow([i + 1, BITRATE_LABEL, rtt_values[i], oneway_values[i]])

print(f"CSV saved as: {csv_filename}")

avg_latency = sum(oneway_values) / len(oneway_values)

samples = list(range(1, len(oneway_values) + 1))

plt.figure(figsize=(8,5))
plt.plot(samples, oneway_values, marker='o', label="One-way Latency")
plt.axhline(avg_latency, linestyle='--', label=f"Average = {avg_latency:.3f} ms")
plt.xlabel("Sample Number")
plt.ylabel("Latency (ms)")
plt.title(f"Latency Test - {BITRATE_LABEL}")
plt.grid(True)
plt.legend()
plt.tight_layout()
plt.savefig(plot_filename, dpi=300)
print(f"Plot saved as: {plot_filename}")
plt.show()
\end{lstlisting}

% ---- PackageSizeLatencyTest.py ----------------------------------------------
\begin{lstlisting}[style=appendix-py, caption={Latency vs.\ packet-size measurement}, label={lst:pkt-size-latency}]
import serial
import time
import matplotlib.pyplot as plt
import csv
import re
from datetime import datetime

SERIAL_PORT = "COM8"
BAUD_RATE = 115200

PACKET_SIZE = 5     # Change for each test
BIT_RATE = 115.2
MEASURE_TIME = 10

timestamp = datetime.now().strftime("%Y%m%d_%H%M%S")

csv_filename = f"packet_size_{PACKET_SIZE}B_{BIT_RATE}kb{timestamp}.csv"
plot_filename = f"packet_size_{PACKET_SIZE}B_{BIT_RATE}kb{timestamp}.png"

ser = serial.Serial(SERIAL_PORT, BAUD_RATE, timeout=1)
time.sleep(2)

rtt_values = []
latency_values = []

pattern = r"Size:\s+(\d+)\s+bytes\s+\|\s+RTT:\s+(\d+)\s+us\s+\|\s+One-way latency:\s+([\d.]+)\s+ms"

print(f"\nRecording latency data for packet size = {PACKET_SIZE} bytes")

start = time.time()

while time.time() - start < MEASURE_TIME:
    line = ser.readline().decode(errors='ignore').strip()
    match = re.search(pattern, line)

    if match:
        size = int(match.group(1))
        rtt = int(match.group(2))
        latency = float(match.group(3))
        rtt_values.append(rtt)
        latency_values.append(latency)
        print(f"Size: {size} B | RTT: {rtt} us | Latency: {latency} ms")

ser.close()

print(f"\nCollected {len(latency_values)} samples")

with open(csv_filename, mode='w', newline='') as file:
    writer = csv.writer(file)
    writer.writerow(["Sample", "PacketSize_Bytes", "RTT_us", "OneWayLatency_ms"])
    for i in range(len(latency_values)):
        writer.writerow([i + 1, PACKET_SIZE, rtt_values[i], latency_values[i]])

print(f"CSV saved as: {csv_filename}")

avg_latency = sum(latency_values) / len(latency_values)

samples = list(range(1, len(latency_values) + 1))

plt.figure(figsize=(8,5))
plt.plot(samples, latency_values, marker='o', label="One-way Latency")
plt.axhline(avg_latency, linestyle='--', label=f"Average = {avg_latency:.3f} ms")
plt.xlabel("Sample Number")
plt.ylabel("Latency (ms)")
plt.title(f"Latency vs Time - Packet Size {PACKET_SIZE} Bytes")
plt.grid(True)
plt.legend()
plt.tight_layout()
plt.savefig(plot_filename, dpi=300)
print(f"Plot saved as: {plot_filename}")
plt.show()
\end{lstlisting}

% ---- packagelossTest.py -----------------------------------------------------
\begin{lstlisting}[style=appendix-py, caption={Packet loss vs.\ distance}, label={lst:packet-loss}]
import serial
import time
import matplotlib.pyplot as plt
import csv
from datetime import datetime

SERIAL_PORT = "COM8"
BAUD_RATE = 115200

DISTANCES = list([20, 40, 50, 70])
EXPECTED_PACKETS = 100

timestamp = datetime.now().strftime("%Y%m%d_%H%M%S")

csv_filename = f"packet_loss_test_{timestamp}.csv"
plot_filename = f"packet_loss_plot_{timestamp}.png"

ser = serial.Serial(SERIAL_PORT, BAUD_RATE, timeout=1)
time.sleep(2)

results = {}

for dist in DISTANCES:
    print(f"\nMove robot to {dist} meters")
    input("Press ENTER when ready...")
    print("Waiting for packets...")

    received_packets = set()
    start_time = time.time()

    while time.time() - start_time < 10:
        line = ser.readline().decode(errors='ignore').strip()

        if "Received packet:" in line:
            try:
                packet_num = int(line.split(":")[1])
                received_packets.add(packet_num)
                print(f"{dist}m -> Packet {packet_num}")
            except:
                pass

    received_count = len(received_packets)
    lost_count = EXPECTED_PACKETS - received_count
    loss_percent = (lost_count / EXPECTED_PACKETS) * 100

    results[dist] = {
        "received": received_count,
        "lost": lost_count,
        "loss_percent": loss_percent
    }

    print("\n========== RESULT ==========")
    print(f"Distance: {dist} m")
    print(f"Received: {received_count}")
    print(f"Lost: {lost_count}")
    print(f"Packet Loss: {loss_percent:.1f}%")
    print("============================")

ser.close()

with open(csv_filename, mode='w', newline='') as file:
    writer = csv.writer(file)
    writer.writerow(["Distance_m", "Received_Packets", "Lost_Packets", "Packet_Loss_Percent"])
    for dist in DISTANCES:
        writer.writerow([dist, results[dist]["received"], results[dist]["lost"], results[dist]["loss_percent"]])

print(f"\nCSV saved as: {csv_filename}")

distances = [dist for dist in DISTANCES]
losses = [results[dist]["loss_percent"] for dist in DISTANCES]

plt.figure(figsize=(8, 5))
plt.plot(distances, losses, marker='o')
plt.xlabel("Distance (m)")
plt.ylabel("Packet Loss (%)")
plt.title("Packet Loss vs Distance")
plt.grid(True)
plt.tight_layout()
plt.savefig(plot_filename, dpi=300)
print(f"Plot saved as: {plot_filename}")
plt.show()
\end{lstlisting}

% ---- OrrientationsTest.py ---------------------------------------------------
\begin{lstlisting}[style=appendix-py, caption={RSSI vs.\ antenna orientation}, label={lst:orientation}]
import serial
import time
import matplotlib.pyplot as plt
import csv
from datetime import datetime

SERIAL_PORT = "COM8"
BAUD_RATE = 115200

ORIENTATIONS = ["Front", "Back", "Left", "Right"]
MEASURE_TIME = 10
ANTENNA = "Horizontal_16.5cm"

timestamp = datetime.now().strftime("%Y%m%d_%H%M%S")

csv_filename = f"rssi_orientation_data_{ANTENNA}_{timestamp}.csv"
plot_filename = f"rssi_orientation_plot_{ANTENNA}_{timestamp}.png"

ser = serial.Serial(SERIAL_PORT, BAUD_RATE, timeout=1)
time.sleep(2)

results = {}

def parse_line(line):
    try:
        parts = line.split(",")
        x = int(parts[0].split(":")[1])
        rssi = int(parts[1].split(":")[1])
        return x, rssi
    except:
        return None, None

for orientation in ORIENTATIONS:
    print(f"\nRotate robot to: {orientation}")
    print("Push joystick to start recording...")

    while True:
        line = ser.readline().decode().strip()
        x, rssi = parse_line(line)
        if x is not None and x > 240:
            print(f"Recording {orientation} orientation...")
            break

    rssi_values = []
    start = time.time()

    while time.time() - start < MEASURE_TIME:
        line = ser.readline().decode().strip()
        x, rssi = parse_line(line)
        if rssi is not None:
            rssi_values.append(rssi)

    results[orientation] = rssi_values
    print(f"Collected {len(rssi_values)} samples")

ser.close()

with open(csv_filename, mode='w', newline='') as file:
    writer = csv.writer(file)
    writer.writerow(["Orientation", "Sample_Number", "RSSI_dBm"])
    for orientation in ORIENTATIONS:
        for i, rssi in enumerate(results[orientation]):
            writer.writerow([orientation, i + 1, rssi])

print(f"\nCSV data saved as: {csv_filename}")

orientations = []
means = []
mins = []
maxs = []

for orientation in ORIENTATIONS:
    vals = results[orientation]
    if len(vals) == 0:
        continue
    orientations.append(orientation)
    means.append(sum(vals) / len(vals))
    mins.append(min(vals))
    maxs.append(max(vals))

plt.figure(figsize=(8, 5))
plt.plot(orientations, means, marker='o', label="Average RSSI")
plt.fill_between(range(len(orientations)), mins, maxs, alpha=0.2, label="Variation")
plt.xlabel("Orientation")
plt.ylabel("RSSI (dBm)")
plt.title("RSSI vs Orientation")
plt.grid(True)
plt.legend()
plt.tight_layout()
plt.savefig(plot_filename, dpi=300)
print(f"Plot saved as: {plot_filename}")
plt.show()
\end{lstlisting}

% ---- RSSItest.py ------------------------------------------------------------
\begin{lstlisting}[style=appendix-py, caption={RSSI vs.\ distance and antenna length}, label={lst:rssi-test}]
import serial
import time
import matplotlib.pyplot as plt
import csv
from datetime import datetime

SERIAL_PORT = "COM8"
BAUD_RATE = 115200

DISTANCES = [1, 5, 10, 20, 40]
MEASURE_TIME = 10
ANTENNA = "Horizontal_16.5cm"

timestamp = datetime.now().strftime("%Y%m%d_%H%M%S")

csv_filename = f"rssi_data_{ANTENNA}_{timestamp}.csv"
plot_filename = f"rssi_plot_{ANTENNA}_{timestamp}.png"

ser = serial.Serial(SERIAL_PORT, BAUD_RATE, timeout=1)
time.sleep(2)

results = {}

def parse_line(line):
    try:
        parts = line.split(",")
        x = int(parts[0].split(":")[1])
        rssi = int(parts[1].split(":")[1])
        return x, rssi
    except:
        return None, None

for dist in DISTANCES:
    print(f"\nMove to {dist} meters. Push joystick to start...")

    while True:
        line = ser.readline().decode().strip()
        x, rssi = parse_line(line)
        if x is not None and x > 240:
            print(f"Recording at {dist} m...")
            break

    rssi_values = []
    start = time.time()

    while time.time() - start < MEASURE_TIME:
        line = ser.readline().decode().strip()
        x, rssi = parse_line(line)
        if rssi is not None:
            rssi_values.append(rssi)

    results[dist] = rssi_values
    print(f"Collected {len(rssi_values)} samples")

ser.close()

with open(csv_filename, mode='w', newline='') as file:
    writer = csv.writer(file)
    writer.writerow(["Distance_m", "Sample_Number", "RSSI_dBm"])
    for dist in DISTANCES:
        for i, rssi in enumerate(results[dist]):
            writer.writerow([dist, i + 1, rssi])

print(f"\nCSV data saved as: {csv_filename}")

distances = []
means = []
mins = []
maxs = []

for d in DISTANCES:
    vals = results[d]
    if len(vals) == 0:
        continue
    distances.append(d)
    means.append(sum(vals) / len(vals))
    mins.append(min(vals))
    maxs.append(max(vals))

plt.figure()
plt.plot(distances, means, marker='o', label="Average RSSI")
plt.fill_between(distances, mins, maxs, alpha=0.2, label="Variation")
plt.xlabel("Distance (m)")
plt.ylabel("RSSI (dBm)")
plt.title("RSSI vs Distance")
plt.grid(True)
plt.legend()
plt.tight_layout()
plt.savefig(plot_filename, dpi=300)
print(f"Plot saved as: {plot_filename}")
plt.show()
\end{lstlisting}

% ---- plot_rssi_antenna.py ---------------------------------------------------
\begin{lstlisting}[style=appendix-py, caption={Multi-file RSSI plotting (antenna comparison)}, label={lst:rssi-plot}]
"""
RSSI vs Distance -- comparison across antenna lengths.

Filename format:  rssi_data_*_<LENGTH>cm_<timestamp>.csv
Columns: Distance_m, Sample_Number, RSSI_dBm
"""

import os
import re
import glob
import pandas as pd
import numpy as np
import matplotlib.pyplot as plt
import matplotlib.ticker as ticker
from pathlib import Path
from collections import defaultdict

DATA_DIR = os.path.dirname(os.path.abspath(__file__))
OUTPUT_DIR = DATA_DIR

LENGTH_ORDER = [16.5, 16.7, 16.9, 17.1, 17.3, 17.5, 17.8]

PALETTE = ["#1f77b4", "#ff7f0e", "#2ca02c", "#d62728",
           "#9467bd", "#8c564b", "#e377c2"]

def parse_length(path: str):
    m = re.search(r'_([\d]+[_.][\d]+)cm[_.]', Path(path).name)
    if m:
        return float(m.group(1).replace("_", "."))
    m2 = re.search(r'_(\d+)cm[_.]', Path(path).name)
    if m2:
        return float(m2.group(1))
    return None

def load_all(data_dir: str) -> pd.DataFrame:
    pattern = os.path.join(data_dir, "rssi_data_*.csv")
    files = glob.glob(pattern)
    if not files:
        raise FileNotFoundError(f"No files matching 'rssi_data_*.csv' in {data_dir}")

    by_length = defaultdict(list)
    for f in sorted(files):
        length = parse_length(f)
        if length is None:
            continue
        df = pd.read_csv(f)
        df.columns = [c.strip() for c in df.columns]
        df = df.rename(columns={c: c.lower().replace(" ", "_") for c in df.columns})
        df = df.rename(columns={"distance_m": "distance", "rssi_dbm": "rssi"})
        df["length_cm"] = length
        by_length[length].append(df)

    frames = []
    for length, dfs in sorted(by_length.items()):
        frames.append(pd.concat(dfs, ignore_index=True))
    return pd.concat(frames, ignore_index=True)

def build_stats(df: pd.DataFrame) -> pd.DataFrame:
    g = df.groupby(["length_cm", "distance"])["rssi"]
    stats = pd.DataFrame({
        "mean": g.mean(), "median": g.median(), "std": g.std(),
        "p5": g.quantile(0.05), "p95": g.quantile(0.95), "n": g.count(),
    }).reset_index()
    return stats

def plot_rssi_vs_distance(stats: pd.DataFrame, out_dir: str):
    lengths = sorted(stats["length_cm"].unique(),
                     key=lambda x: LENGTH_ORDER.index(x) if x in LENGTH_ORDER else x)
    fig, ax = plt.subplots(figsize=(9, 6))
    for i, length in enumerate(lengths):
        sub = stats[stats["length_cm"] == length].sort_values("distance")
        color = PALETTE[i % len(PALETTE)]
        marker = ["o", "s", "^", "D", "v", "P", "X"][i % 7]
        ax.plot(sub["distance"], sub["median"],
                color=color, marker=marker, linewidth=2, markersize=7,
                label=f"{length} cm", zorder=3)
        ax.fill_between(sub["distance"], sub["p5"], sub["p95"],
                        alpha=0.08, color=color)
        ax.fill_between(sub["distance"],
                        sub["median"] - sub["std"], sub["median"] + sub["std"],
                        alpha=0.15, color=color)
    ax.set_xscale("log")
    ax.set_xlabel("Distance (m)", fontsize=11)
    ax.set_ylabel("RSSI (dBm)", fontsize=11)
    ax.set_title("RSSI vs Distance -- Antenna Length Comparison", fontsize=12)
    distances = sorted(stats["distance"].unique())
    ax.set_xticks(distances)
    ax.set_xticklabels([f"{d} m" for d in distances])
    ax.xaxis.set_minor_formatter(ticker.NullFormatter())
    ax.legend(title="Antenna length", loc="lower left")
    ax.grid(True, which="both", linestyle="--", alpha=0.4)
    fig.tight_layout()
    fig.savefig(os.path.join(out_dir, "plot_rssi_vs_distance.png"), dpi=150)
    plt.close(fig)

def plot_bar_per_distance(stats: pd.DataFrame, out_dir: str):
    lengths = sorted(stats["length_cm"].unique(),
                     key=lambda x: LENGTH_ORDER.index(x) if x in LENGTH_ORDER else x)
    distances = sorted(stats["distance"].unique())
    n_len = len(lengths)
    x = np.arange(len(distances))
    width = 0.7 / n_len
    offsets = np.linspace(-(n_len - 1) / 2, (n_len - 1) / 2, n_len) * width
    fig, ax = plt.subplots(figsize=(11, 6))
    for i, length in enumerate(lengths):
        sub = stats[stats["length_cm"] == length].sort_values("distance")
        vals = [sub[sub["distance"] == d]["median"].values[0]
                if len(sub[sub["distance"] == d]) else np.nan for d in distances]
        errs = [sub[sub["distance"] == d]["std"].values[0]
                if len(sub[sub["distance"] == d]) else 0 for d in distances]
        ax.bar(x + offsets[i], vals, width, yerr=errs,
               color=PALETTE[i], alpha=0.8, label=f"{length} cm", capsize=3)
    ax.set_xticks(x)
    ax.set_xticklabels([f"{d} m" for d in distances])
    ax.set_ylabel("Median RSSI (dBm)")
    ax.set_title("Median RSSI per Distance -- Antenna Length Comparison")
    ax.legend(loc="lower left")
    ax.grid(axis="y", linestyle="--", alpha=0.4)
    fig.tight_layout()
    fig.savefig(os.path.join(out_dir, "plot_rssi_bar_per_distance.png"), dpi=150)
    plt.close(fig)

def plot_boxplot_far(df: pd.DataFrame, out_dir: str):
    max_dist = df["distance"].max()
    sub = df[df["distance"] == max_dist]
    lengths = sorted(sub["length_cm"].unique(),
                     key=lambda x: LENGTH_ORDER.index(x) if x in LENGTH_ORDER else x)
    groups = [sub[sub["length_cm"] == l]["rssi"].values for l in lengths]
    labels = [f"{l} cm" for l in lengths]
    colors = PALETTE[:len(lengths)]
    fig, ax = plt.subplots(figsize=(9, 5))
    bp = ax.boxplot(groups, tick_labels=labels, patch_artist=True,
                    medianprops={"color": "black", "linewidth": 2})
    for patch, color in zip(bp["boxes"], colors):
        patch.set_facecolor(color)
        patch.set_alpha(0.75)
    ax.set_title(f"RSSI Distribution at {max_dist} m")
    ax.grid(axis="y", linestyle="--", alpha=0.4)
    fig.tight_layout()
    fig.savefig(os.path.join(out_dir, f"plot_rssi_boxplot_{max_dist}m.png"), dpi=150)
    plt.close(fig)

if __name__ == "__main__":
    df = load_all(DATA_DIR)
    stats = build_stats(df)
    os.makedirs(OUTPUT_DIR, exist_ok=True)
    plot_rssi_vs_distance(stats, OUTPUT_DIR)
    plot_bar_per_distance(stats, OUTPUT_DIR)
    plot_boxplot_far(df, OUTPUT_DIR)
\end{lstlisting}

% ---- plot_latency_vs_bitrate.py ---------------------------------------------
\begin{lstlisting}[style=appendix-py, caption={Latency vs.\ bitrate plotting}, label={lst:latency-bitrate-plot}]
"""
Latency vs Bitrate -- fixed packet size (5 B).

Loads all CSVs matching: latency_*kbps*.csv
Columns: Sample, Bitrate (string), RTT_us, OneWayLatency_ms
"""

import os, re, glob
import pandas as pd
import numpy as np
import matplotlib.pyplot as plt
import matplotlib.ticker as ticker
from pathlib import Path

DATA_DIR = os.path.dirname(os.path.abspath(__file__))
OUTPUT_DIR = DATA_DIR
PACKET_SIZE_B = 5
BITRATE_ORDER = [1.2, 9.6, 57.6, 115.2, 250.0]

def load_all_csvs(data_dir: str) -> pd.DataFrame:
    pattern = os.path.join(data_dir, "latency_*kbps*.csv")
    files = glob.glob(pattern)
    if not files:
        raise FileNotFoundError(f"No CSVs matching 'latency_*kbps*.csv' in {data_dir}")
    frames = []
    for f in sorted(files):
        df = pd.read_csv(f)
        df.columns = [c.strip() for c in df.columns]
        col_map = {c: ["sample", "bitrate_str", "RTT_us", "OWL_ms"][i]
                   for i, c in enumerate(df.columns)}
        df = df.rename(columns=col_map)
        frames.append(df)
    combined = pd.concat(frames, ignore_index=True)
    combined["bitrate_kbps"] = combined["bitrate_str"].str.extract(r"([\d.]+)").astype(float)
    return combined

def remove_outliers(df: pd.DataFrame) -> pd.DataFrame:
    clean = []
    for br, g in df.groupby("bitrate_kbps"):
        med = g["RTT_us"].median()
        clean.append(g[g["RTT_us"] < 5 * med])
    return pd.concat(clean, ignore_index=True)

def build_stats(df: pd.DataFrame) -> pd.DataFrame:
    g = df.groupby("bitrate_kbps")["RTT_us"]
    stats = pd.DataFrame({
        "n": g.count(), "mean_ms": g.mean()/1000, "median_ms": g.median()/1000,
        "std_ms": g.std()/1000, "p5_ms": g.quantile(0.05)/1000,
        "p95_ms": g.quantile(0.95)/1000,
    }).reset_index()
    order = {b: i for i, b in enumerate(BITRATE_ORDER)}
    stats["_order"] = stats["bitrate_kbps"].map(order).fillna(999)
    return stats.sort_values("_order").drop(columns="_order").reset_index(drop=True)

def plot_rtt_line(stats: pd.DataFrame, out_dir: str):
    fig, ax = plt.subplots(figsize=(8, 5))
    ax.plot(stats["bitrate_kbps"], stats["median_ms"],
            color="#2196F3", marker="o", linewidth=2, markersize=7, zorder=3,
            label="Median RTT")
    ax.fill_between(stats["bitrate_kbps"], stats["p5_ms"], stats["p95_ms"],
                    alpha=0.15, color="#2196F3", label="5th-95th percentile")
    ax.fill_between(stats["bitrate_kbps"],
                    stats["median_ms"] - stats["std_ms"],
                    stats["median_ms"] + stats["std_ms"],
                    alpha=0.25, color="#2196F3", label=r"$\pm1\sigma$")
    for _, row in stats.iterrows():
        ax.annotate(f"{row['median_ms']:.1f} ms",
                    xy=(row["bitrate_kbps"], row["median_ms"]),
                    xytext=(0, 10), textcoords="offset points", ha="center", fontsize=8)
    ax.set_xscale("log")
    ax.set_xlabel("Bitrate (kbps)")
    ax.set_ylabel("RTT (ms)")
    ax.set_title(f"Round-Trip Time vs Bitrate  |  Packet size = {PACKET_SIZE_B} B")
    ax.set_xticks(stats["bitrate_kbps"])
    ax.set_xticklabels([f"{b:g}" for b in stats["bitrate_kbps"]])
    ax.xaxis.set_minor_formatter(ticker.NullFormatter())
    ax.legend()
    ax.grid(True, which="both", linestyle="--", alpha=0.4)
    fig.tight_layout()
    fig.savefig(os.path.join(out_dir, "plot_rtt_vs_bitrate.png"), dpi=150)
    plt.close(fig)

def plot_owl_line(df: pd.DataFrame, out_dir: str):
    owl = df.groupby("bitrate_kbps")["OWL_ms"]
    owl_stats = pd.DataFrame({
        "median_ms": owl.median(), "std_ms": owl.std(),
        "p5_ms": owl.quantile(0.05), "p95_ms": owl.quantile(0.95),
    }).reset_index()
    order = {b: i for i, b in enumerate(BITRATE_ORDER)}
    owl_stats["_order"] = owl_stats["bitrate_kbps"].map(order).fillna(999)
    owl_stats = owl_stats.sort_values("_order").drop(columns="_order").reset_index(drop=True)
    fig, ax = plt.subplots(figsize=(8, 5))
    ax.plot(owl_stats["bitrate_kbps"], owl_stats["median_ms"],
            color="#4CAF50", marker="s", linewidth=2, markersize=7, zorder=3,
            label="Median one-way latency")
    ax.fill_between(owl_stats["bitrate_kbps"],
                    owl_stats["median_ms"] - owl_stats["std_ms"],
                    owl_stats["median_ms"] + owl_stats["std_ms"],
                    alpha=0.25, color="#4CAF50", label=r"$\pm1\sigma$")
    for _, row in owl_stats.iterrows():
        ax.annotate(f"{row['median_ms']:.1f} ms",
                    xy=(row["bitrate_kbps"], row["median_ms"]),
                    xytext=(0, 10), textcoords="offset points", ha="center", fontsize=8)
    ax.set_xscale("log")
    ax.set_xlabel("Bitrate (kbps)")
    ax.set_ylabel("One-Way Latency (ms)")
    ax.set_title(f"One-Way Latency vs Bitrate  |  Packet size = {PACKET_SIZE_B} B")
    ax.set_xticks(owl_stats["bitrate_kbps"])
    ax.set_xticklabels([f"{b:g}" for b in owl_stats["bitrate_kbps"]])
    ax.xaxis.set_minor_formatter(ticker.NullFormatter())
    ax.legend()
    ax.grid(True, which="both", linestyle="--", alpha=0.4)
    fig.tight_layout()
    fig.savefig(os.path.join(out_dir, "plot_owl_vs_bitrate.png"), dpi=150)
    plt.close(fig)

def plot_boxplot(df: pd.DataFrame, out_dir: str):
    stats = build_stats(df)
    bitrates = stats["bitrate_kbps"].tolist()
    groups = [df[df["bitrate_kbps"] == br]["RTT_us"].values / 1000 for br in bitrates]
    labels = [f"{b:g} kbps" for b in bitrates]
    fig, ax = plt.subplots(figsize=(9, 5))
    bp = ax.boxplot(groups, tick_labels=labels, patch_artist=True,
                    medianprops={"color": "black", "linewidth": 2})
    colors = ["#2196F3", "#4CAF50", "#FF9800", "#9C27B0", "#F44336"]
    for patch, color in zip(bp["boxes"], colors):
        patch.set_facecolor(color); patch.set_alpha(0.7)
    ax.set_title(f"RTT Distribution per Bitrate  |  Packet size = {PACKET_SIZE_B} B")
    ax.grid(axis="y", linestyle="--", alpha=0.4)
    fig.tight_layout()
    fig.savefig(os.path.join(out_dir, "plot_boxplot_bitrate.png"), dpi=150)
    plt.close(fig)

def plot_bar_comparison(df: pd.DataFrame, out_dir: str):
    stats = build_stats(df)
    bitrates = stats["bitrate_kbps"].tolist()
    x = np.arange(len(bitrates))
    owl_med = [df[df["bitrate_kbps"] == br]["OWL_ms"].median() for br in bitrates]
    owl_std = [df[df["bitrate_kbps"] == br]["OWL_ms"].std() for br in bitrates]
    fig, ax = plt.subplots(figsize=(9, 5))
    width = 0.35
    ax.bar(x - width/2, stats["median_ms"], width, yerr=stats["std_ms"],
           capsize=4, color="#2196F3", alpha=0.8, label="RTT")
    ax.bar(x + width/2, owl_med, width, yerr=owl_std,
           capsize=4, color="#4CAF50", alpha=0.8, label="One-Way Latency")
    ax.set_xticks(x)
    ax.set_xticklabels([f"{b:g} kbps" for b in bitrates])
    ax.set_ylabel("Latency (ms)")
    ax.set_title(f"RTT vs One-Way Latency per Bitrate  |  Packet size = {PACKET_SIZE_B} B")
    ax.legend()
    ax.grid(axis="y", linestyle="--", alpha=0.4)
    fig.tight_layout()
    fig.savefig(os.path.join(out_dir, "plot_bar_rtt_owl.png"), dpi=150)
    plt.close(fig)

if __name__ == "__main__":
    df = load_all_csvs(DATA_DIR)
    df = remove_outliers(df)
    os.makedirs(OUTPUT_DIR, exist_ok=True)
    plot_rtt_line(build_stats(df), OUTPUT_DIR)
    plot_owl_line(df, OUTPUT_DIR)
    plot_boxplot(df, OUTPUT_DIR)
    plot_bar_comparison(df, OUTPUT_DIR)
\end{lstlisting}

% ---- plot_latency.py --------------------------------------------------------
\begin{lstlisting}[style=appendix-py, caption={Latency vs.\ packet-size plotting}, label={lst:latency-pktsize-plot}]
"""
Latency comparison across packet sizes and bitrates.

Filename format: packet_size_<SIZE>B_<BITRATE>kb<timestamp>.csv
Columns: test_num, packet_size_B, RTT_us, one_way_us
"""

import os, re, glob
import pandas as pd
import numpy as np
import matplotlib.pyplot as plt
import matplotlib.ticker as ticker
from pathlib import Path

DATA_DIR = os.path.dirname(os.path.abspath(__file__))
OUTPUT_DIR = DATA_DIR

PACKET_SIZES = [5, 16, 32, 48, 64]
BITRATES = [9.6, 57.6, 115.2]

COLORS = {5: "#2196F3", 16: "#4CAF50", 32: "#FF9800", 48: "#9C27B0", 64: "#F44336"}
MARKERS = {9.6: "o", 57.6: "s", 115.2: "^"}
LINESTYLES = {9.6: "-", 57.6: "--", 115.2: ":"}

def parse_filename(path: str):
    name = Path(path).stem
    m = re.search(r'packet_size_(\d+)B_([\d._]+)kb', name)
    if m:
        return int(m.group(1)), float(m.group(2).replace("_", "."))
    return None

def load_all_csvs(data_dir: str):
    pattern = os.path.join(data_dir, "packet_size_*.csv")
    files = glob.glob(pattern)
    if not files:
        raise FileNotFoundError(f"No CSVs matching 'packet_size_*.csv' in {data_dir}")
    data = {}
    for f in sorted(files):
        key = parse_filename(f)
        if key is None:
            continue
        df = pd.read_csv(f)
        df.columns = ["test_num", "packet_size_B", "RTT_us", "one_way_us"]
        median_rtt = df["RTT_us"].median()
        df = df[df["RTT_us"] < 5 * median_rtt]
        data[key] = df
    return data

def summary_table(data: dict):
    rows = []
    for (ps, br), df in sorted(data.items()):
        rows.append({
            "Packet size (B)": ps, "Bitrate (kbps)": br,
            "RTT mean (ms)": df["RTT_us"].mean()/1000,
            "RTT median (ms)": df["RTT_us"].median()/1000,
            "RTT std (ms)": df["RTT_us"].std()/1000,
            "OWL median (ms)": df["one_way_us"].median()/1000,
            "OWL std (ms)": df["one_way_us"].std()/1000,
        })
    return pd.DataFrame(rows)

def plot_rtt_vs_packet_size(stats: pd.DataFrame, out_dir: str):
    fig, ax = plt.subplots(figsize=(8, 5))
    for br in sorted(stats["Bitrate (kbps)"].unique()):
        sub = stats[stats["Bitrate (kbps)"] == br].sort_values("Packet size (B)")
        ax.plot(sub["Packet size (B)"], sub["RTT median (ms)"],
                marker=MARKERS[br], linestyle=LINESTYLES[br],
                linewidth=1.8, markersize=7, label=f"{br} kbps")
        ax.fill_between(sub["Packet size (B)"],
                        sub["RTT median (ms)"] - sub["RTT std (ms)"],
                        sub["RTT median (ms)"] + sub["RTT std (ms)"], alpha=0.10)
    ax.set_xlabel("Packet Size (bytes)")
    ax.set_ylabel("Median RTT (ms)")
    ax.set_title("Round-Trip Time vs Packet Size")
    ax.set_xticks(PACKET_SIZES)
    ax.legend(title="Bitrate")
    ax.grid(True, linestyle="--", alpha=0.4)
    fig.tight_layout()
    fig.savefig(os.path.join(out_dir, "plot_rtt_vs_packet_size.png"), dpi=150)
    plt.close(fig)

def plot_owl_vs_packet_size(stats: pd.DataFrame, out_dir: str):
    fig, ax = plt.subplots(figsize=(8, 5))
    for br in sorted(stats["Bitrate (kbps)"].unique()):
        sub = stats[stats["Bitrate (kbps)"] == br].sort_values("Packet size (B)")
        ax.plot(sub["Packet size (B)"], sub["OWL median (ms)"],
                marker=MARKERS[br], linestyle=LINESTYLES[br],
                linewidth=1.8, markersize=7, label=f"{br} kbps")
        ax.fill_between(sub["Packet size (B)"],
                        sub["OWL median (ms)"] - sub["OWL std (ms)"],
                        sub["OWL median (ms)"] + sub["OWL std (ms)"], alpha=0.10)
    ax.set_xlabel("Packet Size (bytes)")
    ax.set_ylabel("Median One-Way Latency (ms)")
    ax.set_title("One-Way Latency vs Packet Size")
    ax.set_xticks(PACKET_SIZES)
    ax.legend(title="Bitrate")
    ax.grid(True, linestyle="--", alpha=0.4)
    fig.tight_layout()
    fig.savefig(os.path.join(out_dir, "plot_owl_vs_packet_size.png"), dpi=150)
    plt.close(fig)

def plot_grouped_bar(stats: pd.DataFrame, out_dir: str):
    bitrates = sorted(stats["Bitrate (kbps)"].unique())
    packet_sizes = sorted(stats["Packet size (B)"].unique())
    n_br, n_ps = len(bitrates), len(packet_sizes)
    x = np.arange(n_ps)
    width = 0.7 / n_br
    offsets = np.linspace(-(n_br-1)/2, (n_br-1)/2, n_br) * width
    fig, axes = plt.subplots(1, 2, figsize=(13, 5))
    for ax, metric, ylabel in [
        (axes[0], "RTT median (ms)", "Median RTT (ms)"),
        (axes[1], "OWL median (ms)", "Median One-Way Latency (ms)"),
    ]:
        for i, br in enumerate(bitrates):
            sub = stats[stats["Bitrate (kbps)"] == br].sort_values("Packet size (B)")
            vals = [sub[sub["Packet size (B)"] == ps][metric].values[0]
                    if len(sub[sub["Packet size (B)"] == ps]) else 0 for ps in packet_sizes]
            std_col = metric.replace("median", "std")
            errs = [sub[sub["Packet size (B)"] == ps][std_col].values[0]
                    if len(sub[sub["Packet size (B)"] == ps]) else 0 for ps in packet_sizes]
            ax.bar(x + offsets[i], vals, width, yerr=errs, label=f"{br} kbps", capsize=3)
        ax.set_xlabel("Packet Size (bytes)")
        ax.set_ylabel(ylabel)
        ax.set_xticks(x)
        ax.set_xticklabels([f"{ps} B" for ps in packet_sizes])
        ax.legend(title="Bitrate")
        ax.grid(axis="y", linestyle="--", alpha=0.4)
    fig.suptitle("Latency Comparison -- All Packet Sizes and Bitrates")
    fig.tight_layout()
    fig.savefig(os.path.join(out_dir, "plot_grouped_bar.png"), dpi=150)
    plt.close(fig)

def plot_boxplot(data: dict, out_dir: str):
    bitrates = sorted({k[1] for k in data})
    packet_sizes = sorted({k[0] for k in data})
    fig, axes = plt.subplots(1, len(bitrates), figsize=(5*len(bitrates), 5), sharey=True)
    if len(bitrates) == 1:
        axes = [axes]
    for ax, br in zip(axes, bitrates):
        groups = [data[(ps, br)]["RTT_us"].values/1000 for ps in packet_sizes if (ps, br) in data]
        labels = [f"{ps} B" for ps in packet_sizes if (ps, br) in data]
        bp = ax.boxplot(groups, labels=labels, patch_artist=True,
                        medianprops={"color": "black", "linewidth": 2})
        for patch, ps in zip(bp["boxes"], [ps for ps in packet_sizes if (ps, br) in data]):
            patch.set_facecolor(COLORS.get(ps, "#888")); patch.set_alpha(0.75)
        ax.set_title(f"{br} kbps")
        ax.grid(axis="y", linestyle="--", alpha=0.4)
    axes[0].set_ylabel("RTT (ms)")
    fig.suptitle("RTT Distribution -- Grouped by Bitrate")
    fig.tight_layout()
    fig.savefig(os.path.join(out_dir, "plot_boxplot_rtt.png"), dpi=150)
    plt.close(fig)

def plot_heatmap(stats: pd.DataFrame, out_dir: str):
    import matplotlib.colors as mcolors
    for metric, title, fname in [
        ("RTT median (ms)", "Median RTT (ms)", "plot_heatmap_rtt.png"),
        ("OWL median (ms)", "Median One-Way Latency (ms)", "plot_heatmap_owl.png"),
    ]:
        pivot = stats.pivot(index="Packet size (B)", columns="Bitrate (kbps)", values=metric)
        fig, ax = plt.subplots(figsize=(6, 4))
        im = ax.imshow(pivot.values, aspect="auto", cmap="YlOrRd",
                       norm=mcolors.Normalize(vmin=pivot.values.min(), vmax=pivot.values.max()))
        ax.set_xticks(range(len(pivot.columns)))
        ax.set_xticklabels([f"{c} kbps" for c in pivot.columns])
        ax.set_yticks(range(len(pivot.index)))
        ax.set_yticklabels([f"{r} B" for r in pivot.index])
        ax.set_title(f"Heatmap -- {title}")
        for i in range(len(pivot.index)):
            for j in range(len(pivot.columns)):
                val = pivot.values[i, j]
                if not np.isnan(val):
                    ax.text(j, i, f"{val:.2f}", ha="center", va="center", fontsize=9,
                            color="white" if val > pivot.values.max()*0.65 else "black")
        plt.colorbar(im, ax=ax, label=title)
        fig.tight_layout()
        fig.savefig(os.path.join(out_dir, fname), dpi=150)
        plt.close(fig)

if __name__ == "__main__":
    data = load_all_csvs(DATA_DIR)
    stats = summary_table(data)
    os.makedirs(OUTPUT_DIR, exist_ok=True)
    plot_rtt_vs_packet_size(stats, OUTPUT_DIR)
    plot_owl_vs_packet_size(stats, OUTPUT_DIR)
    plot_grouped_bar(stats, OUTPUT_DIR)
    plot_boxplot(data, OUTPUT_DIR)
    plot_heatmap(stats, OUTPUT_DIR)
\end{lstlisting}

\endinput

% in your report where you want the appendix to appear.
%
% If you already have these packages in your preamble, comment out the
% \usepackage lines below.
% =============================================================================

% ---- Packages (uncomment any NOT already in your preamble) ------------------
% \usepackage{xcolor}
% \usepackage{listings}
% \usepackage{caption}

% ---- Listing style -----------------------------------------------------------
\definecolor{codebg}{rgb}{0.95,0.95,0.97}
\definecolor{codeframe}{rgb}{0.80,0.80,0.85}
\definecolor{stringgreen}{rgb}{0.15,0.55,0.15}
\definecolor{keywordblue}{rgb}{0.10,0.20,0.70}
\definecolor{commentgray}{rgb}{0.35,0.35,0.35}

\lstdefinestyle{appendix}{
  language        = C++,
  basicstyle      = \footnotesize\ttfamily,
  backgroundcolor = \color{codebg},
  frame           = single,
  rulecolor       = \color{codeframe},
  numbers         = left,
  numberstyle     = \tiny\color{gray},
  stepnumber      = 1,
  numbersep       = 8pt,
  tabsize         = 2,
  showspaces      = false,
  showstringspaces = false,
  breaklines      = true,
  breakatwhitespace = true,
  keywordstyle    = \color{keywordblue}\bfseries,
  commentstyle    = \color{commentgray}\itshape,
  stringstyle     = \color{stringgreen},
  captionpos      = t,
  xleftmargin     = 2em,
  framexleftmargin = 1.5em,
}

\lstdefinestyle{appendix-py}{
  style           = appendix,
  language        = Python,
}

\lstdefinestyle{appendix-ini}{
  style           = appendix,
  language        = {[LaTeX]TeX},
}

% =============================================================================
\chapter{Source Code Listings}  \label{app:code}
% =============================================================================

% -----------------------------------------------------------------------------
\section{Remote Controller Firmware (C++)}
% -----------------------------------------------------------------------------

The handheld remote controller firmware runs on an Adafruit Feather~M0
with an RFM69~433\,MHz transceiver, an SSD1306~128$\times$32~OLED display,
and an analog dual-axis joystick.  PlatformIO is used as the build system.

% ---- main.cpp ---------------------------------------------------------------
\begin{lstlisting}[style=appendix, caption={Remote --- main.cpp}, label={lst:remote-main}]
#include <Arduino.h>
#include "OledDis.h"
#include "RadioTrans.h"
#include "Joystick.h"

#define BUTTON_PIN 10
#define VBATPIN A7
#define PWR_CONTROL 5

volatile uint8_t buttonState = 0;

void handleButton() {
  buttonState = !buttonState;
}

void setup() {
  Serial.begin(115200);
  delay(2000);

  pinMode(PWR_CONTROL, OUTPUT);
  digitalWrite(PWR_CONTROL, HIGH);


  attachInterrupt(digitalPinToInterrupt(BUTTON_PIN), handleButton, FALLING);

  Display_init();

  transiver_init();   

  Serial.println("System ready");
}

void loop() {
  uint8_t x, y;
  int8_t RSSI = 0;
  uint8_t battery = 0;
  char robotName[15] = "No connection";

  readJoystick(&x, &y);


  // Remote battery mesurements
  float measuredvbat = analogRead(VBATPIN);
  measuredvbat *= 2;
  measuredvbat *= 3.3;
  measuredvbat /= 1024.0;



  sendData(x, y, buttonState, &battery, robotName, &RSSI);
  Display_draw(robotName, measuredvbat, battery, RSSI);

  delay(100);
}
\end{lstlisting}

% ---- Joystick.cpp -----------------------------------------------------------
\begin{lstlisting}[style=appendix, caption={Remote --- Joystick.cpp}, label={lst:remote-joystick}]
#include <Arduino.h>


void readJoystick(uint8_t *x, uint8_t *y) {
    *x = -analogRead(A0) >> 2;
    *y = -analogRead(A1) >> 2;
}
\end{lstlisting}

% ---- OledDis.cpp ------------------------------------------------------------
\begin{lstlisting}[style=appendix, caption={Remote --- OledDis.cpp}, label={lst:remote-oled}]
#include <Arduino.h>
#include "OledDis.h"
#include <Wire.h>
#include <Adafruit_SSD1306.h>

#define SCREEN_WIDTH 128
#define SCREEN_HEIGHT 32

Adafruit_SSD1306 display(
  SCREEN_WIDTH,
  SCREEN_HEIGHT,
  &Wire,     
  -1
);

void Display_init() {
  Wire.begin();
  delay(100);
  if (!display.begin(SSD1306_SWITCHCAPVCC, 0x3C)) {
    Serial.println("OLED not found");
    while (1);
  }
  Serial.println("OLED found");

  display.clearDisplay();
  display.setRotation(2);
  display.ssd1306_command(SSD1306_DISPLAYON);
  display.ssd1306_command(SSD1306_NORMALDISPLAY);
  display.display();
}

void Display_test() {
  display.clearDisplay();

  display.drawRect(0, 0, 128, 32, SSD1306_WHITE);

  display.drawLine(0, 0, 127, 31, SSD1306_WHITE);

  display.display();
}

void Display_clear() {
  display.clearDisplay();
  display.display();
}

void Display_draw(char *name, float remoteBattery, uint8_t robotBattery, int8_t RSSI) {
  display.clearDisplay();
  display.setTextSize(1);
  display.setTextColor(SSD1306_WHITE);
  display.setCursor(0, 0);

  display.print("Name:");
  display.println(name);

  display.print("R:");
  display.print(remoteBattery, 1);
  display.print("V");

  display.print("B:");
  display.print(robotBattery / 10.0, 1);
  display.println("V");

  display.print("RSSI:");
  display.println(RSSI);

  display.display();
}
\end{lstlisting}

% ---- RadioTrans.cpp ---------------------------------------------------------
\begin{lstlisting}[style=appendix, caption={Remote --- RadioTrans.cpp}, label={lst:remote-radio}]
#include "RadioTrans.h"
#include <SPI.h>
#include <RFM69.h>
#include <RFM69registers.h>

#define RFM69_CS    0
#define RFM69_RST   1
#define RFM69_INT   A4
#define RFM69_DIO2  A3
#define RFM69_DIO3  A5

#define NETWORKID   42
#define NODEID      2
#define TONODEID    1
#define FREQUENCY   RF69_433MHZ
#define LATENCY_PACKET 0x42
#define PACKET_TEST 0x55

RFM69 radio(RFM69_CS, RFM69_INT);

void transiver_init()
{
    pinMode(RFM69_RST, OUTPUT);
    digitalWrite(RFM69_RST, HIGH);
    delay(50);
    digitalWrite(RFM69_RST, LOW);
    delay(50);

    SPI.begin();

    if (!radio.initialize(FREQUENCY, NODEID, NETWORKID))
    {
        Serial.println("RFM69 init FAILED");
        while (1);
    }

    radio.setHighPower();

    Serial.println("RFM69 init OK");
}

bool sendData(uint8_t x, uint8_t y, uint8_t button, uint8_t *battery, char *robotName, int8_t *RSSI)
{
    uint8_t packet[3];
    packet[0] = x;
    packet[1] = y;
    packet[2] = button;

    if (!radio.sendWithRetry(TONODEID, packet, sizeof(packet), 1, 50))
    {
        Serial.println("No ACK received");
        return false;
    }

    *battery = radio.DATA[0];

    for (int i = 0; i < 14; i++)
    {
        robotName[i] = (char)radio.DATA[i + 1];

        if (robotName[i] == '\0')
            break;
    }

    robotName[14] = '\0';

    *RSSI = radio.RSSI;

    return true;
}
\end{lstlisting}

% ---- Joystick.h -------------------------------------------------------------
\begin{lstlisting}[style=appendix, caption={Remote --- Joystick.h}, label={lst:remote-joystick-h}]
#include <Arduino.h>

void readJoystick(uint8_t *x, uint8_t *y);
\end{lstlisting}

% ---- OledDis.h --------------------------------------------------------------
\begin{lstlisting}[style=appendix, caption={Remote --- OledDis.h}, label={lst:remote-oled-h}]
#pragma once
#include <Arduino.h>

void Display_init();
void Display_test();
void Display_clear();
void Display_draw(char *name, float remoteBattery, uint8_t robotBattery, int8_t RSSI);
\end{lstlisting}

% ---- RadioTrans.h -----------------------------------------------------------
\begin{lstlisting}[style=appendix, caption={Remote --- RadioTrans.h}, label={lst:remote-radio-h}]
#ifndef RFM69_H
#define RFM69_H

#include <Arduino.h>

void transiver_init();

bool sendData(uint8_t x, uint8_t y, uint8_t button, uint8_t *battery, char *robotName, int8_t *RSSI);


#endif
\end{lstlisting}

% ---- platformio.ini ---------------------------------------------------------
\begin{lstlisting}[style=appendix-ini, caption={Remote --- platformio.ini}, label={lst:remote-pio}]
[env:adafruit_feather_m0]
platform = atmelsam
board = adafruit_feather_m0
framework = arduino
monitor_speed = 115200
build_flags = -D USB_SERIAL
lib_deps = 
	adafruit/Adafruit SSD1306@^2.5.16
	lowpowerlab/RFM69@^1.6.0
	lowpowerlab/SPIFlash@^101.1.3
\end{lstlisting}

% -----------------------------------------------------------------------------
\section{Robot-Side Firmware (C++)}
% -----------------------------------------------------------------------------

The robot-side firmware runs on a second Adafruit Feather~M0.  It receives
joystick data over the RFM69 link, reads battery voltage and a human-readable
name from a Raspberry Pi over the serial port, and echoes these back to the
remote inside the radio's ACK payload.

% ---- Robot main.cpp ---------------------------------------------------------
\begin{lstlisting}[style=appendix, caption={Robot --- main.cpp}, label={lst:robot-main}]
#include <Arduino.h>
#include <SPI.h>
#include <RFM69.h>
#include "RemoteTransiver.h"

void setup() {
  Serial.begin(115200);
  while (!Serial) delay(10);

  transiver_init();
}

void loop() {
  transiver_receive();
}
\end{lstlisting}

% ---- RemoteTransiver.cpp ----------------------------------------------------
\begin{lstlisting}[style=appendix, caption={Robot --- RemoteTransiver.cpp}, label={lst:robot-trans}]
#include <SPI.h>
#include <RFM69.h>
#include "RemoteTransiver.h"
#include <RFM69registers.h>

#define RFM69_CS   A5
#define RFM69_INT  6
#define RFM69_RST  A4
#define ESTOP      A3
#define LED_PIN    13


#define NETWORKID  42
#define NODEID     1
#define TONODEID   2
#define FREQUENCY  RF69_433MHZ
#define LATENCY_PACKET 0x42
#define PACKET_TEST 0x55

RFM69 radio(RFM69_CS, RFM69_INT);

float batteryVoltage = 0;       
char robotName[16] = "none/con";     
bool piDataReceived = false;     

unsigned long lastReceiveTime = 0;
unsigned long lastNoConnectionPrint = 0;
const unsigned long CONNECTION_TIMEOUT = 1000;


void transiver_init()
{
    pinMode(LED_PIN, OUTPUT);
    pinMode(RFM69_RST, OUTPUT);
    digitalWrite(RFM69_RST, HIGH);
    delay(100);
    digitalWrite(RFM69_RST, LOW);
    delay(100);

    pinMode(ESTOP, OUTPUT);

    if (!radio.initialize(FREQUENCY, NODEID, NETWORKID))
    {
        Serial.println("RFM69 init FAILED");
        while (1);
    }

    radio.setHighPower();

    Serial.println("RFM69 init OK");
    lastReceiveTime = millis();
}

void readPiSerial()
{
    static char line[64];
    static uint8_t idx = 0;

    while (Serial.available())
    {
        char c = Serial.read();
        if (c == '\r') continue;

        if (c == '\n')
        {
            line[idx] = '\0';
            idx = 0;

            // Find "name:" 
            char* namePtr = strstr(line, "name:");
            char* batPtr  = strstr(line, ",bat:");

            if (namePtr && batPtr)
            {
                namePtr += 5; // skip "name:"

                // Copy name up to the comma
                uint8_t i = 0;
                while (namePtr[i] != ',' && namePtr[i] != '\0' && i < 15)
                {
                    robotName[i] = namePtr[i];
                    i++;
                }
                robotName[i] = '\0';

                // Parse float manually
                batPtr += 5; // skip ",bat:"
                batteryVoltage = atof(batPtr);
                piDataReceived = true;


            }

        }
        else
        {
            if (idx < sizeof(line) - 1)
                line[idx++] = c;
            else
                idx = 0;
        }
    }
}


void handleDataPacket(uint8_t x, uint8_t y, uint8_t button)
{

    Serial.print("x:");
    Serial.print(x);
    Serial.print(", y:");
    Serial.print(y);
    Serial.print(", btn:");
    Serial.println(button);

    if (button == 1)
        digitalWrite(ESTOP, HIGH);
    else
        digitalWrite(ESTOP, LOW);
}


void transiver_receive()
{

    readPiSerial();

    if (radio.receiveDone())
    {
        if (radio.ACKRequested())
        {

            uint8_t ackPayload[16];

            // Battery voltage, e.g. 24.3V -> 243
            ackPayload[0] = (uint8_t)(batteryVoltage * 10.0);

            // Copy robot name into remaining bytes
            uint8_t i = 0;
            for (; i < 14 && robotName[i] != '\0'; i++)
            {
                ackPayload[i + 1] = robotName[i];
            }

            // Zero-fill rest
            for (; i < 14; i++)
            {
                ackPayload[i + 1] = 0;
            }

            radio.sendACK(ackPayload, 15);
        }

        uint8_t x      = radio.DATA[0];
        uint8_t y      = radio.DATA[1];
        uint8_t button = radio.DATA[2];

        handleDataPacket(x, y, button);
        lastReceiveTime = millis();
    }

    if (millis() - lastReceiveTime >= CONNECTION_TIMEOUT) {
        if (millis() - lastNoConnectionPrint >= 1000) {
            Serial.println("no connection"); 
            lastNoConnectionPrint = millis();
        }
    }
}
\end{lstlisting}

% ---- RemoteTransiver.h ------------------------------------------------------
\begin{lstlisting}[style=appendix, caption={Robot --- RemoteTransiver.h}, label={lst:robot-trans-h}]
#ifndef RFM69_H
#define RFM69_H

#include <Arduino.h>

void transiver_init();
void transiver_receive();
void readPiSerial();

#endif
\end{lstlisting}

% ---- Robot platformio.ini ---------------------------------------------------
\begin{lstlisting}[style=appendix-ini, caption={Robot --- platformio.ini}, label={lst:robot-pio}]
[env:adafruit_feather_m0]
platform = atmelsam
board = adafruit_feather_m0
framework = arduino
monitor_speed = 115200
build_flags = -D USB_SERIAL
lib_deps = 
	adafruit/Adafruit SSD1306@^2.5.16
	lowpowerlab/RFM69@^1.6.0
	lowpowerlab/SPIFlash@^101.1.3
\end{lstlisting}

% -----------------------------------------------------------------------------
\section{Test and Measurement Scripts (Python)}
% -----------------------------------------------------------------------------

The following Python scripts were used during the project to characterise
the radio link's latency, packet-loss, and signal-strength (RSSI)
performance.  All measurement scripts communicate with the robot-side
Feather~M0 over a serial port and use \texttt{matplotlib} to produce plots.

% ---- test_feather.py --------------------------------------------------------
\begin{lstlisting}[style=appendix-py, caption={Feather serial-communication test}, label={lst:test-feather}]
#!/usr/bin/env python3
"""
Feather M0 Communication Test Script
============================
Usage: python test_feather.py

This script tests serial communication with the Feather M0 (RobotSide).

Setup:
  - Windows: Uses COM8 (change below if different)
  - Linux: Uses /dev/ttyACM1 (change below if different)

Run:
  python test_feather.py

What it does:
  1. Sends "name:TESTBOT,bat:12.34\n" to Feather
  2. Waits for response
  3. Reports what was received

Output examples:
  - "Response: OK" = Feather received and responded
  - "No response" = Feather not receiving data
"""

import serial
import time
import sys

WINDOWS_PORT = 'COM8'
LINUX_PORT = '/dev/ttyACM1'

if sys.platform.startswith('win'):
    SERIAL_PORT = WINDOWS_PORT
elif sys.platform.startswith('linux'):
    SERIAL_PORT = LINUX_PORT
else:
    SERIAL_PORT = '/dev/ttyUSB0'  # Mac

BAUD_RATE = 115200
TIMEOUT = 2

def test_feather():
    print("=== Feather M0 Communication Test ===")
    print(f"Platform: {sys.platform}")
    print(f"Port: {SERIAL_PORT}")
    print(f"Baud: {BAUD_RATE}")
    print("-" * 40)

    try:
        print("Opening serial port...")
        ser = serial.Serial(SERIAL_PORT, BAUD_RATE, timeout=TIMEOUT)
        time.sleep(2)
        print("Port opened!")

        print("\nSending test data...")
        test_msg = "name:TESTBOT,bat:12.34\n"
        ser.write(test_msg.encode())
        print(f"Sent: {test_msg.strip()}")

        print("\nReading response...")
        received_any = False
        
        for i in range(3):
            time.sleep(0.5)
            if ser.in_waiting > 0:
                data = ser.read(ser.in_waiting)
                print(f"Response {i+1}: {data.decode('utf-8', errors='replace').strip()}")
                received_any = True
            else:
                print(f"Attempt {i+1}: No data available")

        if not received_any:
            print("\nNo response received from Feather!")
            print("This means Feather is not responding to the serial data.")

        ser.close()
        print("\n=== Test Complete ===")

    except serial.SerialException as e:
        print(f"\nERROR: {e}")
        print("\nTroubleshooting:")
        print("  - Make sure Feather is connected")
        print("  - Check correct COM port (Windows: COMx, Linux: /dev/ttyACMx)")
        print("  - Close other programs using the port")
        print("  - Try a different USB port")

if __name__ == "__main__":
    test_feather()
\end{lstlisting}

% ---- BitRateLatencyTest.py --------------------------------------------------
\begin{lstlisting}[style=appendix-py, caption={Latency vs.\ bitrate measurement}, label={lst:bitrate-latency}]
import serial
import time
import matplotlib.pyplot as plt
import csv
import re
from datetime import datetime

SERIAL_PORT = "COM8"
BAUD_RATE = 115200

BITRATE_LABEL = "1.2 kbps"   # Change this manually for each test
MEASURE_TIME = 10

timestamp = datetime.now().strftime("%Y%m%d_%H%M%S")

csv_filename = f"latency_{BITRATE_LABEL.replace(' ', '_')}_{timestamp}.csv"
plot_filename = f"latency_{BITRATE_LABEL.replace(' ', '_')}_{timestamp}.png"

ser = serial.Serial(SERIAL_PORT, BAUD_RATE, timeout=1)
time.sleep(2)

rtt_values = []
oneway_values = []

pattern = r"RTT:\s+(\d+)\s+us\s+\|\s+One-way latency:\s+([\d.]+)\s+ms"

print(f"\nRecording latency data for {BITRATE_LABEL}...")

start = time.time()

while time.time() - start < MEASURE_TIME:
    line = ser.readline().decode(errors='ignore').strip()
    match = re.search(pattern, line)

    if match:
        rtt = int(match.group(1))
        latency = float(match.group(2))
        rtt_values.append(rtt)
        oneway_values.append(latency)
        print(f"RTT: {rtt} us | Latency: {latency} ms")

ser.close()

print(f"\nCollected {len(rtt_values)} samples")

with open(csv_filename, mode='w', newline='') as file:
    writer = csv.writer(file)
    writer.writerow(["Sample", "Bitrate", "RTT_us", "OneWayLatency_ms"])
    for i in range(len(rtt_values)):
        writer.writerow([i + 1, BITRATE_LABEL, rtt_values[i], oneway_values[i]])

print(f"CSV saved as: {csv_filename}")

avg_latency = sum(oneway_values) / len(oneway_values)

samples = list(range(1, len(oneway_values) + 1))

plt.figure(figsize=(8,5))
plt.plot(samples, oneway_values, marker='o', label="One-way Latency")
plt.axhline(avg_latency, linestyle='--', label=f"Average = {avg_latency:.3f} ms")
plt.xlabel("Sample Number")
plt.ylabel("Latency (ms)")
plt.title(f"Latency Test - {BITRATE_LABEL}")
plt.grid(True)
plt.legend()
plt.tight_layout()
plt.savefig(plot_filename, dpi=300)
print(f"Plot saved as: {plot_filename}")
plt.show()
\end{lstlisting}

% ---- PackageSizeLatencyTest.py ----------------------------------------------
\begin{lstlisting}[style=appendix-py, caption={Latency vs.\ packet-size measurement}, label={lst:pkt-size-latency}]
import serial
import time
import matplotlib.pyplot as plt
import csv
import re
from datetime import datetime

SERIAL_PORT = "COM8"
BAUD_RATE = 115200

PACKET_SIZE = 5     # Change for each test
BIT_RATE = 115.2
MEASURE_TIME = 10

timestamp = datetime.now().strftime("%Y%m%d_%H%M%S")

csv_filename = f"packet_size_{PACKET_SIZE}B_{BIT_RATE}kb{timestamp}.csv"
plot_filename = f"packet_size_{PACKET_SIZE}B_{BIT_RATE}kb{timestamp}.png"

ser = serial.Serial(SERIAL_PORT, BAUD_RATE, timeout=1)
time.sleep(2)

rtt_values = []
latency_values = []

pattern = r"Size:\s+(\d+)\s+bytes\s+\|\s+RTT:\s+(\d+)\s+us\s+\|\s+One-way latency:\s+([\d.]+)\s+ms"

print(f"\nRecording latency data for packet size = {PACKET_SIZE} bytes")

start = time.time()

while time.time() - start < MEASURE_TIME:
    line = ser.readline().decode(errors='ignore').strip()
    match = re.search(pattern, line)

    if match:
        size = int(match.group(1))
        rtt = int(match.group(2))
        latency = float(match.group(3))
        rtt_values.append(rtt)
        latency_values.append(latency)
        print(f"Size: {size} B | RTT: {rtt} us | Latency: {latency} ms")

ser.close()

print(f"\nCollected {len(latency_values)} samples")

with open(csv_filename, mode='w', newline='') as file:
    writer = csv.writer(file)
    writer.writerow(["Sample", "PacketSize_Bytes", "RTT_us", "OneWayLatency_ms"])
    for i in range(len(latency_values)):
        writer.writerow([i + 1, PACKET_SIZE, rtt_values[i], latency_values[i]])

print(f"CSV saved as: {csv_filename}")

avg_latency = sum(latency_values) / len(latency_values)

samples = list(range(1, len(latency_values) + 1))

plt.figure(figsize=(8,5))
plt.plot(samples, latency_values, marker='o', label="One-way Latency")
plt.axhline(avg_latency, linestyle='--', label=f"Average = {avg_latency:.3f} ms")
plt.xlabel("Sample Number")
plt.ylabel("Latency (ms)")
plt.title(f"Latency vs Time - Packet Size {PACKET_SIZE} Bytes")
plt.grid(True)
plt.legend()
plt.tight_layout()
plt.savefig(plot_filename, dpi=300)
print(f"Plot saved as: {plot_filename}")
plt.show()
\end{lstlisting}

% ---- packagelossTest.py -----------------------------------------------------
\begin{lstlisting}[style=appendix-py, caption={Packet loss vs.\ distance}, label={lst:packet-loss}]
import serial
import time
import matplotlib.pyplot as plt
import csv
from datetime import datetime

SERIAL_PORT = "COM8"
BAUD_RATE = 115200

DISTANCES = list([20, 40, 50, 70])
EXPECTED_PACKETS = 100

timestamp = datetime.now().strftime("%Y%m%d_%H%M%S")

csv_filename = f"packet_loss_test_{timestamp}.csv"
plot_filename = f"packet_loss_plot_{timestamp}.png"

ser = serial.Serial(SERIAL_PORT, BAUD_RATE, timeout=1)
time.sleep(2)

results = {}

for dist in DISTANCES:
    print(f"\nMove robot to {dist} meters")
    input("Press ENTER when ready...")
    print("Waiting for packets...")

    received_packets = set()
    start_time = time.time()

    while time.time() - start_time < 10:
        line = ser.readline().decode(errors='ignore').strip()

        if "Received packet:" in line:
            try:
                packet_num = int(line.split(":")[1])
                received_packets.add(packet_num)
                print(f"{dist}m -> Packet {packet_num}")
            except:
                pass

    received_count = len(received_packets)
    lost_count = EXPECTED_PACKETS - received_count
    loss_percent = (lost_count / EXPECTED_PACKETS) * 100

    results[dist] = {
        "received": received_count,
        "lost": lost_count,
        "loss_percent": loss_percent
    }

    print("\n========== RESULT ==========")
    print(f"Distance: {dist} m")
    print(f"Received: {received_count}")
    print(f"Lost: {lost_count}")
    print(f"Packet Loss: {loss_percent:.1f}%")
    print("============================")

ser.close()

with open(csv_filename, mode='w', newline='') as file:
    writer = csv.writer(file)
    writer.writerow(["Distance_m", "Received_Packets", "Lost_Packets", "Packet_Loss_Percent"])
    for dist in DISTANCES:
        writer.writerow([dist, results[dist]["received"], results[dist]["lost"], results[dist]["loss_percent"]])

print(f"\nCSV saved as: {csv_filename}")

distances = [dist for dist in DISTANCES]
losses = [results[dist]["loss_percent"] for dist in DISTANCES]

plt.figure(figsize=(8, 5))
plt.plot(distances, losses, marker='o')
plt.xlabel("Distance (m)")
plt.ylabel("Packet Loss (%)")
plt.title("Packet Loss vs Distance")
plt.grid(True)
plt.tight_layout()
plt.savefig(plot_filename, dpi=300)
print(f"Plot saved as: {plot_filename}")
plt.show()
\end{lstlisting}

% ---- OrrientationsTest.py ---------------------------------------------------
\begin{lstlisting}[style=appendix-py, caption={RSSI vs.\ antenna orientation}, label={lst:orientation}]
import serial
import time
import matplotlib.pyplot as plt
import csv
from datetime import datetime

SERIAL_PORT = "COM8"
BAUD_RATE = 115200

ORIENTATIONS = ["Front", "Back", "Left", "Right"]
MEASURE_TIME = 10
ANTENNA = "Horizontal_16.5cm"

timestamp = datetime.now().strftime("%Y%m%d_%H%M%S")

csv_filename = f"rssi_orientation_data_{ANTENNA}_{timestamp}.csv"
plot_filename = f"rssi_orientation_plot_{ANTENNA}_{timestamp}.png"

ser = serial.Serial(SERIAL_PORT, BAUD_RATE, timeout=1)
time.sleep(2)

results = {}

def parse_line(line):
    try:
        parts = line.split(",")
        x = int(parts[0].split(":")[1])
        rssi = int(parts[1].split(":")[1])
        return x, rssi
    except:
        return None, None

for orientation in ORIENTATIONS:
    print(f"\nRotate robot to: {orientation}")
    print("Push joystick to start recording...")

    while True:
        line = ser.readline().decode().strip()
        x, rssi = parse_line(line)
        if x is not None and x > 240:
            print(f"Recording {orientation} orientation...")
            break

    rssi_values = []
    start = time.time()

    while time.time() - start < MEASURE_TIME:
        line = ser.readline().decode().strip()
        x, rssi = parse_line(line)
        if rssi is not None:
            rssi_values.append(rssi)

    results[orientation] = rssi_values
    print(f"Collected {len(rssi_values)} samples")

ser.close()

with open(csv_filename, mode='w', newline='') as file:
    writer = csv.writer(file)
    writer.writerow(["Orientation", "Sample_Number", "RSSI_dBm"])
    for orientation in ORIENTATIONS:
        for i, rssi in enumerate(results[orientation]):
            writer.writerow([orientation, i + 1, rssi])

print(f"\nCSV data saved as: {csv_filename}")

orientations = []
means = []
mins = []
maxs = []

for orientation in ORIENTATIONS:
    vals = results[orientation]
    if len(vals) == 0:
        continue
    orientations.append(orientation)
    means.append(sum(vals) / len(vals))
    mins.append(min(vals))
    maxs.append(max(vals))

plt.figure(figsize=(8, 5))
plt.plot(orientations, means, marker='o', label="Average RSSI")
plt.fill_between(range(len(orientations)), mins, maxs, alpha=0.2, label="Variation")
plt.xlabel("Orientation")
plt.ylabel("RSSI (dBm)")
plt.title("RSSI vs Orientation")
plt.grid(True)
plt.legend()
plt.tight_layout()
plt.savefig(plot_filename, dpi=300)
print(f"Plot saved as: {plot_filename}")
plt.show()
\end{lstlisting}

% ---- RSSItest.py ------------------------------------------------------------
\begin{lstlisting}[style=appendix-py, caption={RSSI vs.\ distance and antenna length}, label={lst:rssi-test}]
import serial
import time
import matplotlib.pyplot as plt
import csv
from datetime import datetime

SERIAL_PORT = "COM8"
BAUD_RATE = 115200

DISTANCES = [1, 5, 10, 20, 40]
MEASURE_TIME = 10
ANTENNA = "Horizontal_16.5cm"

timestamp = datetime.now().strftime("%Y%m%d_%H%M%S")

csv_filename = f"rssi_data_{ANTENNA}_{timestamp}.csv"
plot_filename = f"rssi_plot_{ANTENNA}_{timestamp}.png"

ser = serial.Serial(SERIAL_PORT, BAUD_RATE, timeout=1)
time.sleep(2)

results = {}

def parse_line(line):
    try:
        parts = line.split(",")
        x = int(parts[0].split(":")[1])
        rssi = int(parts[1].split(":")[1])
        return x, rssi
    except:
        return None, None

for dist in DISTANCES:
    print(f"\nMove to {dist} meters. Push joystick to start...")

    while True:
        line = ser.readline().decode().strip()
        x, rssi = parse_line(line)
        if x is not None and x > 240:
            print(f"Recording at {dist} m...")
            break

    rssi_values = []
    start = time.time()

    while time.time() - start < MEASURE_TIME:
        line = ser.readline().decode().strip()
        x, rssi = parse_line(line)
        if rssi is not None:
            rssi_values.append(rssi)

    results[dist] = rssi_values
    print(f"Collected {len(rssi_values)} samples")

ser.close()

with open(csv_filename, mode='w', newline='') as file:
    writer = csv.writer(file)
    writer.writerow(["Distance_m", "Sample_Number", "RSSI_dBm"])
    for dist in DISTANCES:
        for i, rssi in enumerate(results[dist]):
            writer.writerow([dist, i + 1, rssi])

print(f"\nCSV data saved as: {csv_filename}")

distances = []
means = []
mins = []
maxs = []

for d in DISTANCES:
    vals = results[d]
    if len(vals) == 0:
        continue
    distances.append(d)
    means.append(sum(vals) / len(vals))
    mins.append(min(vals))
    maxs.append(max(vals))

plt.figure()
plt.plot(distances, means, marker='o', label="Average RSSI")
plt.fill_between(distances, mins, maxs, alpha=0.2, label="Variation")
plt.xlabel("Distance (m)")
plt.ylabel("RSSI (dBm)")
plt.title("RSSI vs Distance")
plt.grid(True)
plt.legend()
plt.tight_layout()
plt.savefig(plot_filename, dpi=300)
print(f"Plot saved as: {plot_filename}")
plt.show()
\end{lstlisting}

% ---- plot_rssi_antenna.py ---------------------------------------------------
\begin{lstlisting}[style=appendix-py, caption={Multi-file RSSI plotting (antenna comparison)}, label={lst:rssi-plot}]
"""
RSSI vs Distance -- comparison across antenna lengths.

Filename format:  rssi_data_*_<LENGTH>cm_<timestamp>.csv
Columns: Distance_m, Sample_Number, RSSI_dBm
"""

import os
import re
import glob
import pandas as pd
import numpy as np
import matplotlib.pyplot as plt
import matplotlib.ticker as ticker
from pathlib import Path
from collections import defaultdict

DATA_DIR = os.path.dirname(os.path.abspath(__file__))
OUTPUT_DIR = DATA_DIR

LENGTH_ORDER = [16.5, 16.7, 16.9, 17.1, 17.3, 17.5, 17.8]

PALETTE = ["#1f77b4", "#ff7f0e", "#2ca02c", "#d62728",
           "#9467bd", "#8c564b", "#e377c2"]

def parse_length(path: str):
    m = re.search(r'_([\d]+[_.][\d]+)cm[_.]', Path(path).name)
    if m:
        return float(m.group(1).replace("_", "."))
    m2 = re.search(r'_(\d+)cm[_.]', Path(path).name)
    if m2:
        return float(m2.group(1))
    return None

def load_all(data_dir: str) -> pd.DataFrame:
    pattern = os.path.join(data_dir, "rssi_data_*.csv")
    files = glob.glob(pattern)
    if not files:
        raise FileNotFoundError(f"No files matching 'rssi_data_*.csv' in {data_dir}")

    by_length = defaultdict(list)
    for f in sorted(files):
        length = parse_length(f)
        if length is None:
            continue
        df = pd.read_csv(f)
        df.columns = [c.strip() for c in df.columns]
        df = df.rename(columns={c: c.lower().replace(" ", "_") for c in df.columns})
        df = df.rename(columns={"distance_m": "distance", "rssi_dbm": "rssi"})
        df["length_cm"] = length
        by_length[length].append(df)

    frames = []
    for length, dfs in sorted(by_length.items()):
        frames.append(pd.concat(dfs, ignore_index=True))
    return pd.concat(frames, ignore_index=True)

def build_stats(df: pd.DataFrame) -> pd.DataFrame:
    g = df.groupby(["length_cm", "distance"])["rssi"]
    stats = pd.DataFrame({
        "mean": g.mean(), "median": g.median(), "std": g.std(),
        "p5": g.quantile(0.05), "p95": g.quantile(0.95), "n": g.count(),
    }).reset_index()
    return stats

def plot_rssi_vs_distance(stats: pd.DataFrame, out_dir: str):
    lengths = sorted(stats["length_cm"].unique(),
                     key=lambda x: LENGTH_ORDER.index(x) if x in LENGTH_ORDER else x)
    fig, ax = plt.subplots(figsize=(9, 6))
    for i, length in enumerate(lengths):
        sub = stats[stats["length_cm"] == length].sort_values("distance")
        color = PALETTE[i % len(PALETTE)]
        marker = ["o", "s", "^", "D", "v", "P", "X"][i % 7]
        ax.plot(sub["distance"], sub["median"],
                color=color, marker=marker, linewidth=2, markersize=7,
                label=f"{length} cm", zorder=3)
        ax.fill_between(sub["distance"], sub["p5"], sub["p95"],
                        alpha=0.08, color=color)
        ax.fill_between(sub["distance"],
                        sub["median"] - sub["std"], sub["median"] + sub["std"],
                        alpha=0.15, color=color)
    ax.set_xscale("log")
    ax.set_xlabel("Distance (m)", fontsize=11)
    ax.set_ylabel("RSSI (dBm)", fontsize=11)
    ax.set_title("RSSI vs Distance -- Antenna Length Comparison", fontsize=12)
    distances = sorted(stats["distance"].unique())
    ax.set_xticks(distances)
    ax.set_xticklabels([f"{d} m" for d in distances])
    ax.xaxis.set_minor_formatter(ticker.NullFormatter())
    ax.legend(title="Antenna length", loc="lower left")
    ax.grid(True, which="both", linestyle="--", alpha=0.4)
    fig.tight_layout()
    fig.savefig(os.path.join(out_dir, "plot_rssi_vs_distance.png"), dpi=150)
    plt.close(fig)

def plot_bar_per_distance(stats: pd.DataFrame, out_dir: str):
    lengths = sorted(stats["length_cm"].unique(),
                     key=lambda x: LENGTH_ORDER.index(x) if x in LENGTH_ORDER else x)
    distances = sorted(stats["distance"].unique())
    n_len = len(lengths)
    x = np.arange(len(distances))
    width = 0.7 / n_len
    offsets = np.linspace(-(n_len - 1) / 2, (n_len - 1) / 2, n_len) * width
    fig, ax = plt.subplots(figsize=(11, 6))
    for i, length in enumerate(lengths):
        sub = stats[stats["length_cm"] == length].sort_values("distance")
        vals = [sub[sub["distance"] == d]["median"].values[0]
                if len(sub[sub["distance"] == d]) else np.nan for d in distances]
        errs = [sub[sub["distance"] == d]["std"].values[0]
                if len(sub[sub["distance"] == d]) else 0 for d in distances]
        ax.bar(x + offsets[i], vals, width, yerr=errs,
               color=PALETTE[i], alpha=0.8, label=f"{length} cm", capsize=3)
    ax.set_xticks(x)
    ax.set_xticklabels([f"{d} m" for d in distances])
    ax.set_ylabel("Median RSSI (dBm)")
    ax.set_title("Median RSSI per Distance -- Antenna Length Comparison")
    ax.legend(loc="lower left")
    ax.grid(axis="y", linestyle="--", alpha=0.4)
    fig.tight_layout()
    fig.savefig(os.path.join(out_dir, "plot_rssi_bar_per_distance.png"), dpi=150)
    plt.close(fig)

def plot_boxplot_far(df: pd.DataFrame, out_dir: str):
    max_dist = df["distance"].max()
    sub = df[df["distance"] == max_dist]
    lengths = sorted(sub["length_cm"].unique(),
                     key=lambda x: LENGTH_ORDER.index(x) if x in LENGTH_ORDER else x)
    groups = [sub[sub["length_cm"] == l]["rssi"].values for l in lengths]
    labels = [f"{l} cm" for l in lengths]
    colors = PALETTE[:len(lengths)]
    fig, ax = plt.subplots(figsize=(9, 5))
    bp = ax.boxplot(groups, tick_labels=labels, patch_artist=True,
                    medianprops={"color": "black", "linewidth": 2})
    for patch, color in zip(bp["boxes"], colors):
        patch.set_facecolor(color)
        patch.set_alpha(0.75)
    ax.set_title(f"RSSI Distribution at {max_dist} m")
    ax.grid(axis="y", linestyle="--", alpha=0.4)
    fig.tight_layout()
    fig.savefig(os.path.join(out_dir, f"plot_rssi_boxplot_{max_dist}m.png"), dpi=150)
    plt.close(fig)

if __name__ == "__main__":
    df = load_all(DATA_DIR)
    stats = build_stats(df)
    os.makedirs(OUTPUT_DIR, exist_ok=True)
    plot_rssi_vs_distance(stats, OUTPUT_DIR)
    plot_bar_per_distance(stats, OUTPUT_DIR)
    plot_boxplot_far(df, OUTPUT_DIR)
\end{lstlisting}

% ---- plot_latency_vs_bitrate.py ---------------------------------------------
\begin{lstlisting}[style=appendix-py, caption={Latency vs.\ bitrate plotting}, label={lst:latency-bitrate-plot}]
"""
Latency vs Bitrate -- fixed packet size (5 B).

Loads all CSVs matching: latency_*kbps*.csv
Columns: Sample, Bitrate (string), RTT_us, OneWayLatency_ms
"""

import os, re, glob
import pandas as pd
import numpy as np
import matplotlib.pyplot as plt
import matplotlib.ticker as ticker
from pathlib import Path

DATA_DIR = os.path.dirname(os.path.abspath(__file__))
OUTPUT_DIR = DATA_DIR
PACKET_SIZE_B = 5
BITRATE_ORDER = [1.2, 9.6, 57.6, 115.2, 250.0]

def load_all_csvs(data_dir: str) -> pd.DataFrame:
    pattern = os.path.join(data_dir, "latency_*kbps*.csv")
    files = glob.glob(pattern)
    if not files:
        raise FileNotFoundError(f"No CSVs matching 'latency_*kbps*.csv' in {data_dir}")
    frames = []
    for f in sorted(files):
        df = pd.read_csv(f)
        df.columns = [c.strip() for c in df.columns]
        col_map = {c: ["sample", "bitrate_str", "RTT_us", "OWL_ms"][i]
                   for i, c in enumerate(df.columns)}
        df = df.rename(columns=col_map)
        frames.append(df)
    combined = pd.concat(frames, ignore_index=True)
    combined["bitrate_kbps"] = combined["bitrate_str"].str.extract(r"([\d.]+)").astype(float)
    return combined

def remove_outliers(df: pd.DataFrame) -> pd.DataFrame:
    clean = []
    for br, g in df.groupby("bitrate_kbps"):
        med = g["RTT_us"].median()
        clean.append(g[g["RTT_us"] < 5 * med])
    return pd.concat(clean, ignore_index=True)

def build_stats(df: pd.DataFrame) -> pd.DataFrame:
    g = df.groupby("bitrate_kbps")["RTT_us"]
    stats = pd.DataFrame({
        "n": g.count(), "mean_ms": g.mean()/1000, "median_ms": g.median()/1000,
        "std_ms": g.std()/1000, "p5_ms": g.quantile(0.05)/1000,
        "p95_ms": g.quantile(0.95)/1000,
    }).reset_index()
    order = {b: i for i, b in enumerate(BITRATE_ORDER)}
    stats["_order"] = stats["bitrate_kbps"].map(order).fillna(999)
    return stats.sort_values("_order").drop(columns="_order").reset_index(drop=True)

def plot_rtt_line(stats: pd.DataFrame, out_dir: str):
    fig, ax = plt.subplots(figsize=(8, 5))
    ax.plot(stats["bitrate_kbps"], stats["median_ms"],
            color="#2196F3", marker="o", linewidth=2, markersize=7, zorder=3,
            label="Median RTT")
    ax.fill_between(stats["bitrate_kbps"], stats["p5_ms"], stats["p95_ms"],
                    alpha=0.15, color="#2196F3", label="5th-95th percentile")
    ax.fill_between(stats["bitrate_kbps"],
                    stats["median_ms"] - stats["std_ms"],
                    stats["median_ms"] + stats["std_ms"],
                    alpha=0.25, color="#2196F3", label=r"$\pm1\sigma$")
    for _, row in stats.iterrows():
        ax.annotate(f"{row['median_ms']:.1f} ms",
                    xy=(row["bitrate_kbps"], row["median_ms"]),
                    xytext=(0, 10), textcoords="offset points", ha="center", fontsize=8)
    ax.set_xscale("log")
    ax.set_xlabel("Bitrate (kbps)")
    ax.set_ylabel("RTT (ms)")
    ax.set_title(f"Round-Trip Time vs Bitrate  |  Packet size = {PACKET_SIZE_B} B")
    ax.set_xticks(stats["bitrate_kbps"])
    ax.set_xticklabels([f"{b:g}" for b in stats["bitrate_kbps"]])
    ax.xaxis.set_minor_formatter(ticker.NullFormatter())
    ax.legend()
    ax.grid(True, which="both", linestyle="--", alpha=0.4)
    fig.tight_layout()
    fig.savefig(os.path.join(out_dir, "plot_rtt_vs_bitrate.png"), dpi=150)
    plt.close(fig)

def plot_owl_line(df: pd.DataFrame, out_dir: str):
    owl = df.groupby("bitrate_kbps")["OWL_ms"]
    owl_stats = pd.DataFrame({
        "median_ms": owl.median(), "std_ms": owl.std(),
        "p5_ms": owl.quantile(0.05), "p95_ms": owl.quantile(0.95),
    }).reset_index()
    order = {b: i for i, b in enumerate(BITRATE_ORDER)}
    owl_stats["_order"] = owl_stats["bitrate_kbps"].map(order).fillna(999)
    owl_stats = owl_stats.sort_values("_order").drop(columns="_order").reset_index(drop=True)
    fig, ax = plt.subplots(figsize=(8, 5))
    ax.plot(owl_stats["bitrate_kbps"], owl_stats["median_ms"],
            color="#4CAF50", marker="s", linewidth=2, markersize=7, zorder=3,
            label="Median one-way latency")
    ax.fill_between(owl_stats["bitrate_kbps"],
                    owl_stats["median_ms"] - owl_stats["std_ms"],
                    owl_stats["median_ms"] + owl_stats["std_ms"],
                    alpha=0.25, color="#4CAF50", label=r"$\pm1\sigma$")
    for _, row in owl_stats.iterrows():
        ax.annotate(f"{row['median_ms']:.1f} ms",
                    xy=(row["bitrate_kbps"], row["median_ms"]),
                    xytext=(0, 10), textcoords="offset points", ha="center", fontsize=8)
    ax.set_xscale("log")
    ax.set_xlabel("Bitrate (kbps)")
    ax.set_ylabel("One-Way Latency (ms)")
    ax.set_title(f"One-Way Latency vs Bitrate  |  Packet size = {PACKET_SIZE_B} B")
    ax.set_xticks(owl_stats["bitrate_kbps"])
    ax.set_xticklabels([f"{b:g}" for b in owl_stats["bitrate_kbps"]])
    ax.xaxis.set_minor_formatter(ticker.NullFormatter())
    ax.legend()
    ax.grid(True, which="both", linestyle="--", alpha=0.4)
    fig.tight_layout()
    fig.savefig(os.path.join(out_dir, "plot_owl_vs_bitrate.png"), dpi=150)
    plt.close(fig)

def plot_boxplot(df: pd.DataFrame, out_dir: str):
    stats = build_stats(df)
    bitrates = stats["bitrate_kbps"].tolist()
    groups = [df[df["bitrate_kbps"] == br]["RTT_us"].values / 1000 for br in bitrates]
    labels = [f"{b:g} kbps" for b in bitrates]
    fig, ax = plt.subplots(figsize=(9, 5))
    bp = ax.boxplot(groups, tick_labels=labels, patch_artist=True,
                    medianprops={"color": "black", "linewidth": 2})
    colors = ["#2196F3", "#4CAF50", "#FF9800", "#9C27B0", "#F44336"]
    for patch, color in zip(bp["boxes"], colors):
        patch.set_facecolor(color); patch.set_alpha(0.7)
    ax.set_title(f"RTT Distribution per Bitrate  |  Packet size = {PACKET_SIZE_B} B")
    ax.grid(axis="y", linestyle="--", alpha=0.4)
    fig.tight_layout()
    fig.savefig(os.path.join(out_dir, "plot_boxplot_bitrate.png"), dpi=150)
    plt.close(fig)

def plot_bar_comparison(df: pd.DataFrame, out_dir: str):
    stats = build_stats(df)
    bitrates = stats["bitrate_kbps"].tolist()
    x = np.arange(len(bitrates))
    owl_med = [df[df["bitrate_kbps"] == br]["OWL_ms"].median() for br in bitrates]
    owl_std = [df[df["bitrate_kbps"] == br]["OWL_ms"].std() for br in bitrates]
    fig, ax = plt.subplots(figsize=(9, 5))
    width = 0.35
    ax.bar(x - width/2, stats["median_ms"], width, yerr=stats["std_ms"],
           capsize=4, color="#2196F3", alpha=0.8, label="RTT")
    ax.bar(x + width/2, owl_med, width, yerr=owl_std,
           capsize=4, color="#4CAF50", alpha=0.8, label="One-Way Latency")
    ax.set_xticks(x)
    ax.set_xticklabels([f"{b:g} kbps" for b in bitrates])
    ax.set_ylabel("Latency (ms)")
    ax.set_title(f"RTT vs One-Way Latency per Bitrate  |  Packet size = {PACKET_SIZE_B} B")
    ax.legend()
    ax.grid(axis="y", linestyle="--", alpha=0.4)
    fig.tight_layout()
    fig.savefig(os.path.join(out_dir, "plot_bar_rtt_owl.png"), dpi=150)
    plt.close(fig)

if __name__ == "__main__":
    df = load_all_csvs(DATA_DIR)
    df = remove_outliers(df)
    os.makedirs(OUTPUT_DIR, exist_ok=True)
    plot_rtt_line(build_stats(df), OUTPUT_DIR)
    plot_owl_line(df, OUTPUT_DIR)
    plot_boxplot(df, OUTPUT_DIR)
    plot_bar_comparison(df, OUTPUT_DIR)
\end{lstlisting}

% ---- plot_latency.py --------------------------------------------------------
\begin{lstlisting}[style=appendix-py, caption={Latency vs.\ packet-size plotting}, label={lst:latency-pktsize-plot}]
"""
Latency comparison across packet sizes and bitrates.

Filename format: packet_size_<SIZE>B_<BITRATE>kb<timestamp>.csv
Columns: test_num, packet_size_B, RTT_us, one_way_us
"""

import os, re, glob
import pandas as pd
import numpy as np
import matplotlib.pyplot as plt
import matplotlib.ticker as ticker
from pathlib import Path

DATA_DIR = os.path.dirname(os.path.abspath(__file__))
OUTPUT_DIR = DATA_DIR

PACKET_SIZES = [5, 16, 32, 48, 64]
BITRATES = [9.6, 57.6, 115.2]

COLORS = {5: "#2196F3", 16: "#4CAF50", 32: "#FF9800", 48: "#9C27B0", 64: "#F44336"}
MARKERS = {9.6: "o", 57.6: "s", 115.2: "^"}
LINESTYLES = {9.6: "-", 57.6: "--", 115.2: ":"}

def parse_filename(path: str):
    name = Path(path).stem
    m = re.search(r'packet_size_(\d+)B_([\d._]+)kb', name)
    if m:
        return int(m.group(1)), float(m.group(2).replace("_", "."))
    return None

def load_all_csvs(data_dir: str):
    pattern = os.path.join(data_dir, "packet_size_*.csv")
    files = glob.glob(pattern)
    if not files:
        raise FileNotFoundError(f"No CSVs matching 'packet_size_*.csv' in {data_dir}")
    data = {}
    for f in sorted(files):
        key = parse_filename(f)
        if key is None:
            continue
        df = pd.read_csv(f)
        df.columns = ["test_num", "packet_size_B", "RTT_us", "one_way_us"]
        median_rtt = df["RTT_us"].median()
        df = df[df["RTT_us"] < 5 * median_rtt]
        data[key] = df
    return data

def summary_table(data: dict):
    rows = []
    for (ps, br), df in sorted(data.items()):
        rows.append({
            "Packet size (B)": ps, "Bitrate (kbps)": br,
            "RTT mean (ms)": df["RTT_us"].mean()/1000,
            "RTT median (ms)": df["RTT_us"].median()/1000,
            "RTT std (ms)": df["RTT_us"].std()/1000,
            "OWL median (ms)": df["one_way_us"].median()/1000,
            "OWL std (ms)": df["one_way_us"].std()/1000,
        })
    return pd.DataFrame(rows)

def plot_rtt_vs_packet_size(stats: pd.DataFrame, out_dir: str):
    fig, ax = plt.subplots(figsize=(8, 5))
    for br in sorted(stats["Bitrate (kbps)"].unique()):
        sub = stats[stats["Bitrate (kbps)"] == br].sort_values("Packet size (B)")
        ax.plot(sub["Packet size (B)"], sub["RTT median (ms)"],
                marker=MARKERS[br], linestyle=LINESTYLES[br],
                linewidth=1.8, markersize=7, label=f"{br} kbps")
        ax.fill_between(sub["Packet size (B)"],
                        sub["RTT median (ms)"] - sub["RTT std (ms)"],
                        sub["RTT median (ms)"] + sub["RTT std (ms)"], alpha=0.10)
    ax.set_xlabel("Packet Size (bytes)")
    ax.set_ylabel("Median RTT (ms)")
    ax.set_title("Round-Trip Time vs Packet Size")
    ax.set_xticks(PACKET_SIZES)
    ax.legend(title="Bitrate")
    ax.grid(True, linestyle="--", alpha=0.4)
    fig.tight_layout()
    fig.savefig(os.path.join(out_dir, "plot_rtt_vs_packet_size.png"), dpi=150)
    plt.close(fig)

def plot_owl_vs_packet_size(stats: pd.DataFrame, out_dir: str):
    fig, ax = plt.subplots(figsize=(8, 5))
    for br in sorted(stats["Bitrate (kbps)"].unique()):
        sub = stats[stats["Bitrate (kbps)"] == br].sort_values("Packet size (B)")
        ax.plot(sub["Packet size (B)"], sub["OWL median (ms)"],
                marker=MARKERS[br], linestyle=LINESTYLES[br],
                linewidth=1.8, markersize=7, label=f"{br} kbps")
        ax.fill_between(sub["Packet size (B)"],
                        sub["OWL median (ms)"] - sub["OWL std (ms)"],
                        sub["OWL median (ms)"] + sub["OWL std (ms)"], alpha=0.10)
    ax.set_xlabel("Packet Size (bytes)")
    ax.set_ylabel("Median One-Way Latency (ms)")
    ax.set_title("One-Way Latency vs Packet Size")
    ax.set_xticks(PACKET_SIZES)
    ax.legend(title="Bitrate")
    ax.grid(True, linestyle="--", alpha=0.4)
    fig.tight_layout()
    fig.savefig(os.path.join(out_dir, "plot_owl_vs_packet_size.png"), dpi=150)
    plt.close(fig)

def plot_grouped_bar(stats: pd.DataFrame, out_dir: str):
    bitrates = sorted(stats["Bitrate (kbps)"].unique())
    packet_sizes = sorted(stats["Packet size (B)"].unique())
    n_br, n_ps = len(bitrates), len(packet_sizes)
    x = np.arange(n_ps)
    width = 0.7 / n_br
    offsets = np.linspace(-(n_br-1)/2, (n_br-1)/2, n_br) * width
    fig, axes = plt.subplots(1, 2, figsize=(13, 5))
    for ax, metric, ylabel in [
        (axes[0], "RTT median (ms)", "Median RTT (ms)"),
        (axes[1], "OWL median (ms)", "Median One-Way Latency (ms)"),
    ]:
        for i, br in enumerate(bitrates):
            sub = stats[stats["Bitrate (kbps)"] == br].sort_values("Packet size (B)")
            vals = [sub[sub["Packet size (B)"] == ps][metric].values[0]
                    if len(sub[sub["Packet size (B)"] == ps]) else 0 for ps in packet_sizes]
            std_col = metric.replace("median", "std")
            errs = [sub[sub["Packet size (B)"] == ps][std_col].values[0]
                    if len(sub[sub["Packet size (B)"] == ps]) else 0 for ps in packet_sizes]
            ax.bar(x + offsets[i], vals, width, yerr=errs, label=f"{br} kbps", capsize=3)
        ax.set_xlabel("Packet Size (bytes)")
        ax.set_ylabel(ylabel)
        ax.set_xticks(x)
        ax.set_xticklabels([f"{ps} B" for ps in packet_sizes])
        ax.legend(title="Bitrate")
        ax.grid(axis="y", linestyle="--", alpha=0.4)
    fig.suptitle("Latency Comparison -- All Packet Sizes and Bitrates")
    fig.tight_layout()
    fig.savefig(os.path.join(out_dir, "plot_grouped_bar.png"), dpi=150)
    plt.close(fig)

def plot_boxplot(data: dict, out_dir: str):
    bitrates = sorted({k[1] for k in data})
    packet_sizes = sorted({k[0] for k in data})
    fig, axes = plt.subplots(1, len(bitrates), figsize=(5*len(bitrates), 5), sharey=True)
    if len(bitrates) == 1:
        axes = [axes]
    for ax, br in zip(axes, bitrates):
        groups = [data[(ps, br)]["RTT_us"].values/1000 for ps in packet_sizes if (ps, br) in data]
        labels = [f"{ps} B" for ps in packet_sizes if (ps, br) in data]
        bp = ax.boxplot(groups, labels=labels, patch_artist=True,
                        medianprops={"color": "black", "linewidth": 2})
        for patch, ps in zip(bp["boxes"], [ps for ps in packet_sizes if (ps, br) in data]):
            patch.set_facecolor(COLORS.get(ps, "#888")); patch.set_alpha(0.75)
        ax.set_title(f"{br} kbps")
        ax.grid(axis="y", linestyle="--", alpha=0.4)
    axes[0].set_ylabel("RTT (ms)")
    fig.suptitle("RTT Distribution -- Grouped by Bitrate")
    fig.tight_layout()
    fig.savefig(os.path.join(out_dir, "plot_boxplot_rtt.png"), dpi=150)
    plt.close(fig)

def plot_heatmap(stats: pd.DataFrame, out_dir: str):
    import matplotlib.colors as mcolors
    for metric, title, fname in [
        ("RTT median (ms)", "Median RTT (ms)", "plot_heatmap_rtt.png"),
        ("OWL median (ms)", "Median One-Way Latency (ms)", "plot_heatmap_owl.png"),
    ]:
        pivot = stats.pivot(index="Packet size (B)", columns="Bitrate (kbps)", values=metric)
        fig, ax = plt.subplots(figsize=(6, 4))
        im = ax.imshow(pivot.values, aspect="auto", cmap="YlOrRd",
                       norm=mcolors.Normalize(vmin=pivot.values.min(), vmax=pivot.values.max()))
        ax.set_xticks(range(len(pivot.columns)))
        ax.set_xticklabels([f"{c} kbps" for c in pivot.columns])
        ax.set_yticks(range(len(pivot.index)))
        ax.set_yticklabels([f"{r} B" for r in pivot.index])
        ax.set_title(f"Heatmap -- {title}")
        for i in range(len(pivot.index)):
            for j in range(len(pivot.columns)):
                val = pivot.values[i, j]
                if not np.isnan(val):
                    ax.text(j, i, f"{val:.2f}", ha="center", va="center", fontsize=9,
                            color="white" if val > pivot.values.max()*0.65 else "black")
        plt.colorbar(im, ax=ax, label=title)
        fig.tight_layout()
        fig.savefig(os.path.join(out_dir, fname), dpi=150)
        plt.close(fig)

if __name__ == "__main__":
    data = load_all_csvs(DATA_DIR)
    stats = summary_table(data)
    os.makedirs(OUTPUT_DIR, exist_ok=True)
    plot_rtt_vs_packet_size(stats, OUTPUT_DIR)
    plot_owl_vs_packet_size(stats, OUTPUT_DIR)
    plot_grouped_bar(stats, OUTPUT_DIR)
    plot_boxplot(data, OUTPUT_DIR)
    plot_heatmap(stats, OUTPUT_DIR)
\end{lstlisting}

\endinput

% in your report where you want the appendix to appear.
%
% If you already have these packages in your preamble, comment out the
% \usepackage lines below.
% =============================================================================

% ---- Packages (uncomment any NOT already in your preamble) ------------------
% \usepackage{xcolor}
% \usepackage{listings}
% \usepackage{caption}

% ---- Listing style -----------------------------------------------------------
\definecolor{codebg}{rgb}{0.95,0.95,0.97}
\definecolor{codeframe}{rgb}{0.80,0.80,0.85}
\definecolor{stringgreen}{rgb}{0.15,0.55,0.15}
\definecolor{keywordblue}{rgb}{0.10,0.20,0.70}
\definecolor{commentgray}{rgb}{0.35,0.35,0.35}

\lstdefinestyle{appendix}{
  language        = C++,
  basicstyle      = \footnotesize\ttfamily,
  backgroundcolor = \color{codebg},
  frame           = single,
  rulecolor       = \color{codeframe},
  numbers         = left,
  numberstyle     = \tiny\color{gray},
  stepnumber      = 1,
  numbersep       = 8pt,
  tabsize         = 2,
  showspaces      = false,
  showstringspaces = false,
  breaklines      = true,
  breakatwhitespace = true,
  keywordstyle    = \color{keywordblue}\bfseries,
  commentstyle    = \color{commentgray}\itshape,
  stringstyle     = \color{stringgreen},
  captionpos      = t,
  xleftmargin     = 2em,
  framexleftmargin = 1.5em,
}

\lstdefinestyle{appendix-py}{
  style           = appendix,
  language        = Python,
}

\lstdefinestyle{appendix-ini}{
  style           = appendix,
  language        = {[LaTeX]TeX},
}

% =============================================================================
\chapter{Source Code Listings}  \label{app:code}
% =============================================================================

% -----------------------------------------------------------------------------
\section{Remote Controller Firmware (C++)}
% -----------------------------------------------------------------------------

The handheld remote controller firmware runs on an Adafruit Feather~M0
with an RFM69~433\,MHz transceiver, an SSD1306~128$\times$32~OLED display,
and an analog dual-axis joystick.  PlatformIO is used as the build system.

% ---- main.cpp ---------------------------------------------------------------
\begin{lstlisting}[style=appendix, caption={Remote --- main.cpp}, label={lst:remote-main}]
#include <Arduino.h>
#include "OledDis.h"
#include "RadioTrans.h"
#include "Joystick.h"

#define BUTTON_PIN 10
#define VBATPIN A7
#define PWR_CONTROL 5

volatile uint8_t buttonState = 0;

void handleButton() {
  buttonState = !buttonState;
}

void setup() {
  Serial.begin(115200);
  delay(2000);

  pinMode(PWR_CONTROL, OUTPUT);
  digitalWrite(PWR_CONTROL, HIGH);


  attachInterrupt(digitalPinToInterrupt(BUTTON_PIN), handleButton, FALLING);

  Display_init();

  transiver_init();   

  Serial.println("System ready");
}

void loop() {
  uint8_t x, y;
  int8_t RSSI = 0;
  uint8_t battery = 0;
  char robotName[15] = "No connection";

  readJoystick(&x, &y);


  // Remote battery mesurements
  float measuredvbat = analogRead(VBATPIN);
  measuredvbat *= 2;
  measuredvbat *= 3.3;
  measuredvbat /= 1024.0;



  sendData(x, y, buttonState, &battery, robotName, &RSSI);
  Display_draw(robotName, measuredvbat, battery, RSSI);

  delay(100);
}
\end{lstlisting}

% ---- Joystick.cpp -----------------------------------------------------------
\begin{lstlisting}[style=appendix, caption={Remote --- Joystick.cpp}, label={lst:remote-joystick}]
#include <Arduino.h>


void readJoystick(uint8_t *x, uint8_t *y) {
    *x = -analogRead(A0) >> 2;
    *y = -analogRead(A1) >> 2;
}
\end{lstlisting}

% ---- OledDis.cpp ------------------------------------------------------------
\begin{lstlisting}[style=appendix, caption={Remote --- OledDis.cpp}, label={lst:remote-oled}]
#include <Arduino.h>
#include "OledDis.h"
#include <Wire.h>
#include <Adafruit_SSD1306.h>

#define SCREEN_WIDTH 128
#define SCREEN_HEIGHT 32

Adafruit_SSD1306 display(
  SCREEN_WIDTH,
  SCREEN_HEIGHT,
  &Wire,     
  -1
);

void Display_init() {
  Wire.begin();
  delay(100);
  if (!display.begin(SSD1306_SWITCHCAPVCC, 0x3C)) {
    Serial.println("OLED not found");
    while (1);
  }
  Serial.println("OLED found");

  display.clearDisplay();
  display.setRotation(2);
  display.ssd1306_command(SSD1306_DISPLAYON);
  display.ssd1306_command(SSD1306_NORMALDISPLAY);
  display.display();
}

void Display_test() {
  display.clearDisplay();

  display.drawRect(0, 0, 128, 32, SSD1306_WHITE);

  display.drawLine(0, 0, 127, 31, SSD1306_WHITE);

  display.display();
}

void Display_clear() {
  display.clearDisplay();
  display.display();
}

void Display_draw(char *name, float remoteBattery, uint8_t robotBattery, int8_t RSSI) {
  display.clearDisplay();
  display.setTextSize(1);
  display.setTextColor(SSD1306_WHITE);
  display.setCursor(0, 0);

  display.print("Name:");
  display.println(name);

  display.print("R:");
  display.print(remoteBattery, 1);
  display.print("V");

  display.print("B:");
  display.print(robotBattery / 10.0, 1);
  display.println("V");

  display.print("RSSI:");
  display.println(RSSI);

  display.display();
}
\end{lstlisting}

% ---- RadioTrans.cpp ---------------------------------------------------------
\begin{lstlisting}[style=appendix, caption={Remote --- RadioTrans.cpp}, label={lst:remote-radio}]
#include "RadioTrans.h"
#include <SPI.h>
#include <RFM69.h>
#include <RFM69registers.h>

#define RFM69_CS    0
#define RFM69_RST   1
#define RFM69_INT   A4
#define RFM69_DIO2  A3
#define RFM69_DIO3  A5

#define NETWORKID   42
#define NODEID      2
#define TONODEID    1
#define FREQUENCY   RF69_433MHZ
#define LATENCY_PACKET 0x42
#define PACKET_TEST 0x55

RFM69 radio(RFM69_CS, RFM69_INT);

void transiver_init()
{
    pinMode(RFM69_RST, OUTPUT);
    digitalWrite(RFM69_RST, HIGH);
    delay(50);
    digitalWrite(RFM69_RST, LOW);
    delay(50);

    SPI.begin();

    if (!radio.initialize(FREQUENCY, NODEID, NETWORKID))
    {
        Serial.println("RFM69 init FAILED");
        while (1);
    }

    radio.setHighPower();

    Serial.println("RFM69 init OK");
}

bool sendData(uint8_t x, uint8_t y, uint8_t button, uint8_t *battery, char *robotName, int8_t *RSSI)
{
    uint8_t packet[3];
    packet[0] = x;
    packet[1] = y;
    packet[2] = button;

    if (!radio.sendWithRetry(TONODEID, packet, sizeof(packet), 1, 50))
    {
        Serial.println("No ACK received");
        return false;
    }

    *battery = radio.DATA[0];

    for (int i = 0; i < 14; i++)
    {
        robotName[i] = (char)radio.DATA[i + 1];

        if (robotName[i] == '\0')
            break;
    }

    robotName[14] = '\0';

    *RSSI = radio.RSSI;

    return true;
}
\end{lstlisting}

% ---- Joystick.h -------------------------------------------------------------
\begin{lstlisting}[style=appendix, caption={Remote --- Joystick.h}, label={lst:remote-joystick-h}]
#include <Arduino.h>

void readJoystick(uint8_t *x, uint8_t *y);
\end{lstlisting}

% ---- OledDis.h --------------------------------------------------------------
\begin{lstlisting}[style=appendix, caption={Remote --- OledDis.h}, label={lst:remote-oled-h}]
#pragma once
#include <Arduino.h>

void Display_init();
void Display_test();
void Display_clear();
void Display_draw(char *name, float remoteBattery, uint8_t robotBattery, int8_t RSSI);
\end{lstlisting}

% ---- RadioTrans.h -----------------------------------------------------------
\begin{lstlisting}[style=appendix, caption={Remote --- RadioTrans.h}, label={lst:remote-radio-h}]
#ifndef RFM69_H
#define RFM69_H

#include <Arduino.h>

void transiver_init();

bool sendData(uint8_t x, uint8_t y, uint8_t button, uint8_t *battery, char *robotName, int8_t *RSSI);


#endif
\end{lstlisting}

% ---- platformio.ini ---------------------------------------------------------
\begin{lstlisting}[style=appendix-ini, caption={Remote --- platformio.ini}, label={lst:remote-pio}]
[env:adafruit_feather_m0]
platform = atmelsam
board = adafruit_feather_m0
framework = arduino
monitor_speed = 115200
build_flags = -D USB_SERIAL
lib_deps = 
	adafruit/Adafruit SSD1306@^2.5.16
	lowpowerlab/RFM69@^1.6.0
	lowpowerlab/SPIFlash@^101.1.3
\end{lstlisting}

% -----------------------------------------------------------------------------
\section{Robot-Side Firmware (C++)}
% -----------------------------------------------------------------------------

The robot-side firmware runs on a second Adafruit Feather~M0.  It receives
joystick data over the RFM69 link, reads battery voltage and a human-readable
name from a Raspberry Pi over the serial port, and echoes these back to the
remote inside the radio's ACK payload.

% ---- Robot main.cpp ---------------------------------------------------------
\begin{lstlisting}[style=appendix, caption={Robot --- main.cpp}, label={lst:robot-main}]
#include <Arduino.h>
#include <SPI.h>
#include <RFM69.h>
#include "RemoteTransiver.h"

void setup() {
  Serial.begin(115200);
  while (!Serial) delay(10);

  transiver_init();
}

void loop() {
  transiver_receive();
}
\end{lstlisting}

% ---- RemoteTransiver.cpp ----------------------------------------------------
\begin{lstlisting}[style=appendix, caption={Robot --- RemoteTransiver.cpp}, label={lst:robot-trans}]
#include <SPI.h>
#include <RFM69.h>
#include "RemoteTransiver.h"
#include <RFM69registers.h>

#define RFM69_CS   A5
#define RFM69_INT  6
#define RFM69_RST  A4
#define ESTOP      A3
#define LED_PIN    13


#define NETWORKID  42
#define NODEID     1
#define TONODEID   2
#define FREQUENCY  RF69_433MHZ
#define LATENCY_PACKET 0x42
#define PACKET_TEST 0x55

RFM69 radio(RFM69_CS, RFM69_INT);

float batteryVoltage = 0;       
char robotName[16] = "none/con";     
bool piDataReceived = false;     

unsigned long lastReceiveTime = 0;
unsigned long lastNoConnectionPrint = 0;
const unsigned long CONNECTION_TIMEOUT = 1000;


void transiver_init()
{
    pinMode(LED_PIN, OUTPUT);
    pinMode(RFM69_RST, OUTPUT);
    digitalWrite(RFM69_RST, HIGH);
    delay(100);
    digitalWrite(RFM69_RST, LOW);
    delay(100);

    pinMode(ESTOP, OUTPUT);

    if (!radio.initialize(FREQUENCY, NODEID, NETWORKID))
    {
        Serial.println("RFM69 init FAILED");
        while (1);
    }

    radio.setHighPower();

    Serial.println("RFM69 init OK");
    lastReceiveTime = millis();
}

void readPiSerial()
{
    static char line[64];
    static uint8_t idx = 0;

    while (Serial.available())
    {
        char c = Serial.read();
        if (c == '\r') continue;

        if (c == '\n')
        {
            line[idx] = '\0';
            idx = 0;

            // Find "name:" 
            char* namePtr = strstr(line, "name:");
            char* batPtr  = strstr(line, ",bat:");

            if (namePtr && batPtr)
            {
                namePtr += 5; // skip "name:"

                // Copy name up to the comma
                uint8_t i = 0;
                while (namePtr[i] != ',' && namePtr[i] != '\0' && i < 15)
                {
                    robotName[i] = namePtr[i];
                    i++;
                }
                robotName[i] = '\0';

                // Parse float manually
                batPtr += 5; // skip ",bat:"
                batteryVoltage = atof(batPtr);
                piDataReceived = true;


            }

        }
        else
        {
            if (idx < sizeof(line) - 1)
                line[idx++] = c;
            else
                idx = 0;
        }
    }
}


void handleDataPacket(uint8_t x, uint8_t y, uint8_t button)
{

    Serial.print("x:");
    Serial.print(x);
    Serial.print(", y:");
    Serial.print(y);
    Serial.print(", btn:");
    Serial.println(button);

    if (button == 1)
        digitalWrite(ESTOP, HIGH);
    else
        digitalWrite(ESTOP, LOW);
}


void transiver_receive()
{

    readPiSerial();

    if (radio.receiveDone())
    {
        if (radio.ACKRequested())
        {

            uint8_t ackPayload[16];

            // Battery voltage, e.g. 24.3V -> 243
            ackPayload[0] = (uint8_t)(batteryVoltage * 10.0);

            // Copy robot name into remaining bytes
            uint8_t i = 0;
            for (; i < 14 && robotName[i] != '\0'; i++)
            {
                ackPayload[i + 1] = robotName[i];
            }

            // Zero-fill rest
            for (; i < 14; i++)
            {
                ackPayload[i + 1] = 0;
            }

            radio.sendACK(ackPayload, 15);
        }

        uint8_t x      = radio.DATA[0];
        uint8_t y      = radio.DATA[1];
        uint8_t button = radio.DATA[2];

        handleDataPacket(x, y, button);
        lastReceiveTime = millis();
    }

    if (millis() - lastReceiveTime >= CONNECTION_TIMEOUT) {
        if (millis() - lastNoConnectionPrint >= 1000) {
            Serial.println("no connection"); 
            lastNoConnectionPrint = millis();
        }
    }
}
\end{lstlisting}

% ---- RemoteTransiver.h ------------------------------------------------------
\begin{lstlisting}[style=appendix, caption={Robot --- RemoteTransiver.h}, label={lst:robot-trans-h}]
#ifndef RFM69_H
#define RFM69_H

#include <Arduino.h>

void transiver_init();
void transiver_receive();
void readPiSerial();

#endif
\end{lstlisting}

% ---- Robot platformio.ini ---------------------------------------------------
\begin{lstlisting}[style=appendix-ini, caption={Robot --- platformio.ini}, label={lst:robot-pio}]
[env:adafruit_feather_m0]
platform = atmelsam
board = adafruit_feather_m0
framework = arduino
monitor_speed = 115200
build_flags = -D USB_SERIAL
lib_deps = 
	adafruit/Adafruit SSD1306@^2.5.16
	lowpowerlab/RFM69@^1.6.0
	lowpowerlab/SPIFlash@^101.1.3
\end{lstlisting}

% -----------------------------------------------------------------------------
\section{Test and Measurement Scripts (Python)}
% -----------------------------------------------------------------------------

The following Python scripts were used during the project to characterise
the radio link's latency, packet-loss, and signal-strength (RSSI)
performance.  All measurement scripts communicate with the robot-side
Feather~M0 over a serial port and use \texttt{matplotlib} to produce plots.

% ---- test_feather.py --------------------------------------------------------
\begin{lstlisting}[style=appendix-py, caption={Feather serial-communication test}, label={lst:test-feather}]
#!/usr/bin/env python3
"""
Feather M0 Communication Test Script
============================
Usage: python test_feather.py

This script tests serial communication with the Feather M0 (RobotSide).

Setup:
  - Windows: Uses COM8 (change below if different)
  - Linux: Uses /dev/ttyACM1 (change below if different)

Run:
  python test_feather.py

What it does:
  1. Sends "name:TESTBOT,bat:12.34\n" to Feather
  2. Waits for response
  3. Reports what was received

Output examples:
  - "Response: OK" = Feather received and responded
  - "No response" = Feather not receiving data
"""

import serial
import time
import sys

WINDOWS_PORT = 'COM8'
LINUX_PORT = '/dev/ttyACM1'

if sys.platform.startswith('win'):
    SERIAL_PORT = WINDOWS_PORT
elif sys.platform.startswith('linux'):
    SERIAL_PORT = LINUX_PORT
else:
    SERIAL_PORT = '/dev/ttyUSB0'  # Mac

BAUD_RATE = 115200
TIMEOUT = 2

def test_feather():
    print("=== Feather M0 Communication Test ===")
    print(f"Platform: {sys.platform}")
    print(f"Port: {SERIAL_PORT}")
    print(f"Baud: {BAUD_RATE}")
    print("-" * 40)

    try:
        print("Opening serial port...")
        ser = serial.Serial(SERIAL_PORT, BAUD_RATE, timeout=TIMEOUT)
        time.sleep(2)
        print("Port opened!")

        print("\nSending test data...")
        test_msg = "name:TESTBOT,bat:12.34\n"
        ser.write(test_msg.encode())
        print(f"Sent: {test_msg.strip()}")

        print("\nReading response...")
        received_any = False
        
        for i in range(3):
            time.sleep(0.5)
            if ser.in_waiting > 0:
                data = ser.read(ser.in_waiting)
                print(f"Response {i+1}: {data.decode('utf-8', errors='replace').strip()}")
                received_any = True
            else:
                print(f"Attempt {i+1}: No data available")

        if not received_any:
            print("\nNo response received from Feather!")
            print("This means Feather is not responding to the serial data.")

        ser.close()
        print("\n=== Test Complete ===")

    except serial.SerialException as e:
        print(f"\nERROR: {e}")
        print("\nTroubleshooting:")
        print("  - Make sure Feather is connected")
        print("  - Check correct COM port (Windows: COMx, Linux: /dev/ttyACMx)")
        print("  - Close other programs using the port")
        print("  - Try a different USB port")

if __name__ == "__main__":
    test_feather()
\end{lstlisting}

% ---- BitRateLatencyTest.py --------------------------------------------------
\begin{lstlisting}[style=appendix-py, caption={Latency vs.\ bitrate measurement}, label={lst:bitrate-latency}]
import serial
import time
import matplotlib.pyplot as plt
import csv
import re
from datetime import datetime

SERIAL_PORT = "COM8"
BAUD_RATE = 115200

BITRATE_LABEL = "1.2 kbps"   # Change this manually for each test
MEASURE_TIME = 10

timestamp = datetime.now().strftime("%Y%m%d_%H%M%S")

csv_filename = f"latency_{BITRATE_LABEL.replace(' ', '_')}_{timestamp}.csv"
plot_filename = f"latency_{BITRATE_LABEL.replace(' ', '_')}_{timestamp}.png"

ser = serial.Serial(SERIAL_PORT, BAUD_RATE, timeout=1)
time.sleep(2)

rtt_values = []
oneway_values = []

pattern = r"RTT:\s+(\d+)\s+us\s+\|\s+One-way latency:\s+([\d.]+)\s+ms"

print(f"\nRecording latency data for {BITRATE_LABEL}...")

start = time.time()

while time.time() - start < MEASURE_TIME:
    line = ser.readline().decode(errors='ignore').strip()
    match = re.search(pattern, line)

    if match:
        rtt = int(match.group(1))
        latency = float(match.group(2))
        rtt_values.append(rtt)
        oneway_values.append(latency)
        print(f"RTT: {rtt} us | Latency: {latency} ms")

ser.close()

print(f"\nCollected {len(rtt_values)} samples")

with open(csv_filename, mode='w', newline='') as file:
    writer = csv.writer(file)
    writer.writerow(["Sample", "Bitrate", "RTT_us", "OneWayLatency_ms"])
    for i in range(len(rtt_values)):
        writer.writerow([i + 1, BITRATE_LABEL, rtt_values[i], oneway_values[i]])

print(f"CSV saved as: {csv_filename}")

avg_latency = sum(oneway_values) / len(oneway_values)

samples = list(range(1, len(oneway_values) + 1))

plt.figure(figsize=(8,5))
plt.plot(samples, oneway_values, marker='o', label="One-way Latency")
plt.axhline(avg_latency, linestyle='--', label=f"Average = {avg_latency:.3f} ms")
plt.xlabel("Sample Number")
plt.ylabel("Latency (ms)")
plt.title(f"Latency Test - {BITRATE_LABEL}")
plt.grid(True)
plt.legend()
plt.tight_layout()
plt.savefig(plot_filename, dpi=300)
print(f"Plot saved as: {plot_filename}")
plt.show()
\end{lstlisting}

% ---- PackageSizeLatencyTest.py ----------------------------------------------
\begin{lstlisting}[style=appendix-py, caption={Latency vs.\ packet-size measurement}, label={lst:pkt-size-latency}]
import serial
import time
import matplotlib.pyplot as plt
import csv
import re
from datetime import datetime

SERIAL_PORT = "COM8"
BAUD_RATE = 115200

PACKET_SIZE = 5     # Change for each test
BIT_RATE = 115.2
MEASURE_TIME = 10

timestamp = datetime.now().strftime("%Y%m%d_%H%M%S")

csv_filename = f"packet_size_{PACKET_SIZE}B_{BIT_RATE}kb{timestamp}.csv"
plot_filename = f"packet_size_{PACKET_SIZE}B_{BIT_RATE}kb{timestamp}.png"

ser = serial.Serial(SERIAL_PORT, BAUD_RATE, timeout=1)
time.sleep(2)

rtt_values = []
latency_values = []

pattern = r"Size:\s+(\d+)\s+bytes\s+\|\s+RTT:\s+(\d+)\s+us\s+\|\s+One-way latency:\s+([\d.]+)\s+ms"

print(f"\nRecording latency data for packet size = {PACKET_SIZE} bytes")

start = time.time()

while time.time() - start < MEASURE_TIME:
    line = ser.readline().decode(errors='ignore').strip()
    match = re.search(pattern, line)

    if match:
        size = int(match.group(1))
        rtt = int(match.group(2))
        latency = float(match.group(3))
        rtt_values.append(rtt)
        latency_values.append(latency)
        print(f"Size: {size} B | RTT: {rtt} us | Latency: {latency} ms")

ser.close()

print(f"\nCollected {len(latency_values)} samples")

with open(csv_filename, mode='w', newline='') as file:
    writer = csv.writer(file)
    writer.writerow(["Sample", "PacketSize_Bytes", "RTT_us", "OneWayLatency_ms"])
    for i in range(len(latency_values)):
        writer.writerow([i + 1, PACKET_SIZE, rtt_values[i], latency_values[i]])

print(f"CSV saved as: {csv_filename}")

avg_latency = sum(latency_values) / len(latency_values)

samples = list(range(1, len(latency_values) + 1))

plt.figure(figsize=(8,5))
plt.plot(samples, latency_values, marker='o', label="One-way Latency")
plt.axhline(avg_latency, linestyle='--', label=f"Average = {avg_latency:.3f} ms")
plt.xlabel("Sample Number")
plt.ylabel("Latency (ms)")
plt.title(f"Latency vs Time - Packet Size {PACKET_SIZE} Bytes")
plt.grid(True)
plt.legend()
plt.tight_layout()
plt.savefig(plot_filename, dpi=300)
print(f"Plot saved as: {plot_filename}")
plt.show()
\end{lstlisting}

% ---- packagelossTest.py -----------------------------------------------------
\begin{lstlisting}[style=appendix-py, caption={Packet loss vs.\ distance}, label={lst:packet-loss}]
import serial
import time
import matplotlib.pyplot as plt
import csv
from datetime import datetime

SERIAL_PORT = "COM8"
BAUD_RATE = 115200

DISTANCES = list([20, 40, 50, 70])
EXPECTED_PACKETS = 100

timestamp = datetime.now().strftime("%Y%m%d_%H%M%S")

csv_filename = f"packet_loss_test_{timestamp}.csv"
plot_filename = f"packet_loss_plot_{timestamp}.png"

ser = serial.Serial(SERIAL_PORT, BAUD_RATE, timeout=1)
time.sleep(2)

results = {}

for dist in DISTANCES:
    print(f"\nMove robot to {dist} meters")
    input("Press ENTER when ready...")
    print("Waiting for packets...")

    received_packets = set()
    start_time = time.time()

    while time.time() - start_time < 10:
        line = ser.readline().decode(errors='ignore').strip()

        if "Received packet:" in line:
            try:
                packet_num = int(line.split(":")[1])
                received_packets.add(packet_num)
                print(f"{dist}m -> Packet {packet_num}")
            except:
                pass

    received_count = len(received_packets)
    lost_count = EXPECTED_PACKETS - received_count
    loss_percent = (lost_count / EXPECTED_PACKETS) * 100

    results[dist] = {
        "received": received_count,
        "lost": lost_count,
        "loss_percent": loss_percent
    }

    print("\n========== RESULT ==========")
    print(f"Distance: {dist} m")
    print(f"Received: {received_count}")
    print(f"Lost: {lost_count}")
    print(f"Packet Loss: {loss_percent:.1f}%")
    print("============================")

ser.close()

with open(csv_filename, mode='w', newline='') as file:
    writer = csv.writer(file)
    writer.writerow(["Distance_m", "Received_Packets", "Lost_Packets", "Packet_Loss_Percent"])
    for dist in DISTANCES:
        writer.writerow([dist, results[dist]["received"], results[dist]["lost"], results[dist]["loss_percent"]])

print(f"\nCSV saved as: {csv_filename}")

distances = [dist for dist in DISTANCES]
losses = [results[dist]["loss_percent"] for dist in DISTANCES]

plt.figure(figsize=(8, 5))
plt.plot(distances, losses, marker='o')
plt.xlabel("Distance (m)")
plt.ylabel("Packet Loss (%)")
plt.title("Packet Loss vs Distance")
plt.grid(True)
plt.tight_layout()
plt.savefig(plot_filename, dpi=300)
print(f"Plot saved as: {plot_filename}")
plt.show()
\end{lstlisting}

% ---- OrrientationsTest.py ---------------------------------------------------
\begin{lstlisting}[style=appendix-py, caption={RSSI vs.\ antenna orientation}, label={lst:orientation}]
import serial
import time
import matplotlib.pyplot as plt
import csv
from datetime import datetime

SERIAL_PORT = "COM8"
BAUD_RATE = 115200

ORIENTATIONS = ["Front", "Back", "Left", "Right"]
MEASURE_TIME = 10
ANTENNA = "Horizontal_16.5cm"

timestamp = datetime.now().strftime("%Y%m%d_%H%M%S")

csv_filename = f"rssi_orientation_data_{ANTENNA}_{timestamp}.csv"
plot_filename = f"rssi_orientation_plot_{ANTENNA}_{timestamp}.png"

ser = serial.Serial(SERIAL_PORT, BAUD_RATE, timeout=1)
time.sleep(2)

results = {}

def parse_line(line):
    try:
        parts = line.split(",")
        x = int(parts[0].split(":")[1])
        rssi = int(parts[1].split(":")[1])
        return x, rssi
    except:
        return None, None

for orientation in ORIENTATIONS:
    print(f"\nRotate robot to: {orientation}")
    print("Push joystick to start recording...")

    while True:
        line = ser.readline().decode().strip()
        x, rssi = parse_line(line)
        if x is not None and x > 240:
            print(f"Recording {orientation} orientation...")
            break

    rssi_values = []
    start = time.time()

    while time.time() - start < MEASURE_TIME:
        line = ser.readline().decode().strip()
        x, rssi = parse_line(line)
        if rssi is not None:
            rssi_values.append(rssi)

    results[orientation] = rssi_values
    print(f"Collected {len(rssi_values)} samples")

ser.close()

with open(csv_filename, mode='w', newline='') as file:
    writer = csv.writer(file)
    writer.writerow(["Orientation", "Sample_Number", "RSSI_dBm"])
    for orientation in ORIENTATIONS:
        for i, rssi in enumerate(results[orientation]):
            writer.writerow([orientation, i + 1, rssi])

print(f"\nCSV data saved as: {csv_filename}")

orientations = []
means = []
mins = []
maxs = []

for orientation in ORIENTATIONS:
    vals = results[orientation]
    if len(vals) == 0:
        continue
    orientations.append(orientation)
    means.append(sum(vals) / len(vals))
    mins.append(min(vals))
    maxs.append(max(vals))

plt.figure(figsize=(8, 5))
plt.plot(orientations, means, marker='o', label="Average RSSI")
plt.fill_between(range(len(orientations)), mins, maxs, alpha=0.2, label="Variation")
plt.xlabel("Orientation")
plt.ylabel("RSSI (dBm)")
plt.title("RSSI vs Orientation")
plt.grid(True)
plt.legend()
plt.tight_layout()
plt.savefig(plot_filename, dpi=300)
print(f"Plot saved as: {plot_filename}")
plt.show()
\end{lstlisting}

% ---- RSSItest.py ------------------------------------------------------------
\begin{lstlisting}[style=appendix-py, caption={RSSI vs.\ distance and antenna length}, label={lst:rssi-test}]
import serial
import time
import matplotlib.pyplot as plt
import csv
from datetime import datetime

SERIAL_PORT = "COM8"
BAUD_RATE = 115200

DISTANCES = [1, 5, 10, 20, 40]
MEASURE_TIME = 10
ANTENNA = "Horizontal_16.5cm"

timestamp = datetime.now().strftime("%Y%m%d_%H%M%S")

csv_filename = f"rssi_data_{ANTENNA}_{timestamp}.csv"
plot_filename = f"rssi_plot_{ANTENNA}_{timestamp}.png"

ser = serial.Serial(SERIAL_PORT, BAUD_RATE, timeout=1)
time.sleep(2)

results = {}

def parse_line(line):
    try:
        parts = line.split(",")
        x = int(parts[0].split(":")[1])
        rssi = int(parts[1].split(":")[1])
        return x, rssi
    except:
        return None, None

for dist in DISTANCES:
    print(f"\nMove to {dist} meters. Push joystick to start...")

    while True:
        line = ser.readline().decode().strip()
        x, rssi = parse_line(line)
        if x is not None and x > 240:
            print(f"Recording at {dist} m...")
            break

    rssi_values = []
    start = time.time()

    while time.time() - start < MEASURE_TIME:
        line = ser.readline().decode().strip()
        x, rssi = parse_line(line)
        if rssi is not None:
            rssi_values.append(rssi)

    results[dist] = rssi_values
    print(f"Collected {len(rssi_values)} samples")

ser.close()

with open(csv_filename, mode='w', newline='') as file:
    writer = csv.writer(file)
    writer.writerow(["Distance_m", "Sample_Number", "RSSI_dBm"])
    for dist in DISTANCES:
        for i, rssi in enumerate(results[dist]):
            writer.writerow([dist, i + 1, rssi])

print(f"\nCSV data saved as: {csv_filename}")

distances = []
means = []
mins = []
maxs = []

for d in DISTANCES:
    vals = results[d]
    if len(vals) == 0:
        continue
    distances.append(d)
    means.append(sum(vals) / len(vals))
    mins.append(min(vals))
    maxs.append(max(vals))

plt.figure()
plt.plot(distances, means, marker='o', label="Average RSSI")
plt.fill_between(distances, mins, maxs, alpha=0.2, label="Variation")
plt.xlabel("Distance (m)")
plt.ylabel("RSSI (dBm)")
plt.title("RSSI vs Distance")
plt.grid(True)
plt.legend()
plt.tight_layout()
plt.savefig(plot_filename, dpi=300)
print(f"Plot saved as: {plot_filename}")
plt.show()
\end{lstlisting}

% ---- plot_rssi_antenna.py ---------------------------------------------------
\begin{lstlisting}[style=appendix-py, caption={Multi-file RSSI plotting (antenna comparison)}, label={lst:rssi-plot}]
"""
RSSI vs Distance -- comparison across antenna lengths.

Filename format:  rssi_data_*_<LENGTH>cm_<timestamp>.csv
Columns: Distance_m, Sample_Number, RSSI_dBm
"""

import os
import re
import glob
import pandas as pd
import numpy as np
import matplotlib.pyplot as plt
import matplotlib.ticker as ticker
from pathlib import Path
from collections import defaultdict

DATA_DIR = os.path.dirname(os.path.abspath(__file__))
OUTPUT_DIR = DATA_DIR

LENGTH_ORDER = [16.5, 16.7, 16.9, 17.1, 17.3, 17.5, 17.8]

PALETTE = ["#1f77b4", "#ff7f0e", "#2ca02c", "#d62728",
           "#9467bd", "#8c564b", "#e377c2"]

def parse_length(path: str):
    m = re.search(r'_([\d]+[_.][\d]+)cm[_.]', Path(path).name)
    if m:
        return float(m.group(1).replace("_", "."))
    m2 = re.search(r'_(\d+)cm[_.]', Path(path).name)
    if m2:
        return float(m2.group(1))
    return None

def load_all(data_dir: str) -> pd.DataFrame:
    pattern = os.path.join(data_dir, "rssi_data_*.csv")
    files = glob.glob(pattern)
    if not files:
        raise FileNotFoundError(f"No files matching 'rssi_data_*.csv' in {data_dir}")

    by_length = defaultdict(list)
    for f in sorted(files):
        length = parse_length(f)
        if length is None:
            continue
        df = pd.read_csv(f)
        df.columns = [c.strip() for c in df.columns]
        df = df.rename(columns={c: c.lower().replace(" ", "_") for c in df.columns})
        df = df.rename(columns={"distance_m": "distance", "rssi_dbm": "rssi"})
        df["length_cm"] = length
        by_length[length].append(df)

    frames = []
    for length, dfs in sorted(by_length.items()):
        frames.append(pd.concat(dfs, ignore_index=True))
    return pd.concat(frames, ignore_index=True)

def build_stats(df: pd.DataFrame) -> pd.DataFrame:
    g = df.groupby(["length_cm", "distance"])["rssi"]
    stats = pd.DataFrame({
        "mean": g.mean(), "median": g.median(), "std": g.std(),
        "p5": g.quantile(0.05), "p95": g.quantile(0.95), "n": g.count(),
    }).reset_index()
    return stats

def plot_rssi_vs_distance(stats: pd.DataFrame, out_dir: str):
    lengths = sorted(stats["length_cm"].unique(),
                     key=lambda x: LENGTH_ORDER.index(x) if x in LENGTH_ORDER else x)
    fig, ax = plt.subplots(figsize=(9, 6))
    for i, length in enumerate(lengths):
        sub = stats[stats["length_cm"] == length].sort_values("distance")
        color = PALETTE[i % len(PALETTE)]
        marker = ["o", "s", "^", "D", "v", "P", "X"][i % 7]
        ax.plot(sub["distance"], sub["median"],
                color=color, marker=marker, linewidth=2, markersize=7,
                label=f"{length} cm", zorder=3)
        ax.fill_between(sub["distance"], sub["p5"], sub["p95"],
                        alpha=0.08, color=color)
        ax.fill_between(sub["distance"],
                        sub["median"] - sub["std"], sub["median"] + sub["std"],
                        alpha=0.15, color=color)
    ax.set_xscale("log")
    ax.set_xlabel("Distance (m)", fontsize=11)
    ax.set_ylabel("RSSI (dBm)", fontsize=11)
    ax.set_title("RSSI vs Distance -- Antenna Length Comparison", fontsize=12)
    distances = sorted(stats["distance"].unique())
    ax.set_xticks(distances)
    ax.set_xticklabels([f"{d} m" for d in distances])
    ax.xaxis.set_minor_formatter(ticker.NullFormatter())
    ax.legend(title="Antenna length", loc="lower left")
    ax.grid(True, which="both", linestyle="--", alpha=0.4)
    fig.tight_layout()
    fig.savefig(os.path.join(out_dir, "plot_rssi_vs_distance.png"), dpi=150)
    plt.close(fig)

def plot_bar_per_distance(stats: pd.DataFrame, out_dir: str):
    lengths = sorted(stats["length_cm"].unique(),
                     key=lambda x: LENGTH_ORDER.index(x) if x in LENGTH_ORDER else x)
    distances = sorted(stats["distance"].unique())
    n_len = len(lengths)
    x = np.arange(len(distances))
    width = 0.7 / n_len
    offsets = np.linspace(-(n_len - 1) / 2, (n_len - 1) / 2, n_len) * width
    fig, ax = plt.subplots(figsize=(11, 6))
    for i, length in enumerate(lengths):
        sub = stats[stats["length_cm"] == length].sort_values("distance")
        vals = [sub[sub["distance"] == d]["median"].values[0]
                if len(sub[sub["distance"] == d]) else np.nan for d in distances]
        errs = [sub[sub["distance"] == d]["std"].values[0]
                if len(sub[sub["distance"] == d]) else 0 for d in distances]
        ax.bar(x + offsets[i], vals, width, yerr=errs,
               color=PALETTE[i], alpha=0.8, label=f"{length} cm", capsize=3)
    ax.set_xticks(x)
    ax.set_xticklabels([f"{d} m" for d in distances])
    ax.set_ylabel("Median RSSI (dBm)")
    ax.set_title("Median RSSI per Distance -- Antenna Length Comparison")
    ax.legend(loc="lower left")
    ax.grid(axis="y", linestyle="--", alpha=0.4)
    fig.tight_layout()
    fig.savefig(os.path.join(out_dir, "plot_rssi_bar_per_distance.png"), dpi=150)
    plt.close(fig)

def plot_boxplot_far(df: pd.DataFrame, out_dir: str):
    max_dist = df["distance"].max()
    sub = df[df["distance"] == max_dist]
    lengths = sorted(sub["length_cm"].unique(),
                     key=lambda x: LENGTH_ORDER.index(x) if x in LENGTH_ORDER else x)
    groups = [sub[sub["length_cm"] == l]["rssi"].values for l in lengths]
    labels = [f"{l} cm" for l in lengths]
    colors = PALETTE[:len(lengths)]
    fig, ax = plt.subplots(figsize=(9, 5))
    bp = ax.boxplot(groups, tick_labels=labels, patch_artist=True,
                    medianprops={"color": "black", "linewidth": 2})
    for patch, color in zip(bp["boxes"], colors):
        patch.set_facecolor(color)
        patch.set_alpha(0.75)
    ax.set_title(f"RSSI Distribution at {max_dist} m")
    ax.grid(axis="y", linestyle="--", alpha=0.4)
    fig.tight_layout()
    fig.savefig(os.path.join(out_dir, f"plot_rssi_boxplot_{max_dist}m.png"), dpi=150)
    plt.close(fig)

if __name__ == "__main__":
    df = load_all(DATA_DIR)
    stats = build_stats(df)
    os.makedirs(OUTPUT_DIR, exist_ok=True)
    plot_rssi_vs_distance(stats, OUTPUT_DIR)
    plot_bar_per_distance(stats, OUTPUT_DIR)
    plot_boxplot_far(df, OUTPUT_DIR)
\end{lstlisting}

% ---- plot_latency_vs_bitrate.py ---------------------------------------------
\begin{lstlisting}[style=appendix-py, caption={Latency vs.\ bitrate plotting}, label={lst:latency-bitrate-plot}]
"""
Latency vs Bitrate -- fixed packet size (5 B).

Loads all CSVs matching: latency_*kbps*.csv
Columns: Sample, Bitrate (string), RTT_us, OneWayLatency_ms
"""

import os, re, glob
import pandas as pd
import numpy as np
import matplotlib.pyplot as plt
import matplotlib.ticker as ticker
from pathlib import Path

DATA_DIR = os.path.dirname(os.path.abspath(__file__))
OUTPUT_DIR = DATA_DIR
PACKET_SIZE_B = 5
BITRATE_ORDER = [1.2, 9.6, 57.6, 115.2, 250.0]

def load_all_csvs(data_dir: str) -> pd.DataFrame:
    pattern = os.path.join(data_dir, "latency_*kbps*.csv")
    files = glob.glob(pattern)
    if not files:
        raise FileNotFoundError(f"No CSVs matching 'latency_*kbps*.csv' in {data_dir}")
    frames = []
    for f in sorted(files):
        df = pd.read_csv(f)
        df.columns = [c.strip() for c in df.columns]
        col_map = {c: ["sample", "bitrate_str", "RTT_us", "OWL_ms"][i]
                   for i, c in enumerate(df.columns)}
        df = df.rename(columns=col_map)
        frames.append(df)
    combined = pd.concat(frames, ignore_index=True)
    combined["bitrate_kbps"] = combined["bitrate_str"].str.extract(r"([\d.]+)").astype(float)
    return combined

def remove_outliers(df: pd.DataFrame) -> pd.DataFrame:
    clean = []
    for br, g in df.groupby("bitrate_kbps"):
        med = g["RTT_us"].median()
        clean.append(g[g["RTT_us"] < 5 * med])
    return pd.concat(clean, ignore_index=True)

def build_stats(df: pd.DataFrame) -> pd.DataFrame:
    g = df.groupby("bitrate_kbps")["RTT_us"]
    stats = pd.DataFrame({
        "n": g.count(), "mean_ms": g.mean()/1000, "median_ms": g.median()/1000,
        "std_ms": g.std()/1000, "p5_ms": g.quantile(0.05)/1000,
        "p95_ms": g.quantile(0.95)/1000,
    }).reset_index()
    order = {b: i for i, b in enumerate(BITRATE_ORDER)}
    stats["_order"] = stats["bitrate_kbps"].map(order).fillna(999)
    return stats.sort_values("_order").drop(columns="_order").reset_index(drop=True)

def plot_rtt_line(stats: pd.DataFrame, out_dir: str):
    fig, ax = plt.subplots(figsize=(8, 5))
    ax.plot(stats["bitrate_kbps"], stats["median_ms"],
            color="#2196F3", marker="o", linewidth=2, markersize=7, zorder=3,
            label="Median RTT")
    ax.fill_between(stats["bitrate_kbps"], stats["p5_ms"], stats["p95_ms"],
                    alpha=0.15, color="#2196F3", label="5th-95th percentile")
    ax.fill_between(stats["bitrate_kbps"],
                    stats["median_ms"] - stats["std_ms"],
                    stats["median_ms"] + stats["std_ms"],
                    alpha=0.25, color="#2196F3", label=r"$\pm1\sigma$")
    for _, row in stats.iterrows():
        ax.annotate(f"{row['median_ms']:.1f} ms",
                    xy=(row["bitrate_kbps"], row["median_ms"]),
                    xytext=(0, 10), textcoords="offset points", ha="center", fontsize=8)
    ax.set_xscale("log")
    ax.set_xlabel("Bitrate (kbps)")
    ax.set_ylabel("RTT (ms)")
    ax.set_title(f"Round-Trip Time vs Bitrate  |  Packet size = {PACKET_SIZE_B} B")
    ax.set_xticks(stats["bitrate_kbps"])
    ax.set_xticklabels([f"{b:g}" for b in stats["bitrate_kbps"]])
    ax.xaxis.set_minor_formatter(ticker.NullFormatter())
    ax.legend()
    ax.grid(True, which="both", linestyle="--", alpha=0.4)
    fig.tight_layout()
    fig.savefig(os.path.join(out_dir, "plot_rtt_vs_bitrate.png"), dpi=150)
    plt.close(fig)

def plot_owl_line(df: pd.DataFrame, out_dir: str):
    owl = df.groupby("bitrate_kbps")["OWL_ms"]
    owl_stats = pd.DataFrame({
        "median_ms": owl.median(), "std_ms": owl.std(),
        "p5_ms": owl.quantile(0.05), "p95_ms": owl.quantile(0.95),
    }).reset_index()
    order = {b: i for i, b in enumerate(BITRATE_ORDER)}
    owl_stats["_order"] = owl_stats["bitrate_kbps"].map(order).fillna(999)
    owl_stats = owl_stats.sort_values("_order").drop(columns="_order").reset_index(drop=True)
    fig, ax = plt.subplots(figsize=(8, 5))
    ax.plot(owl_stats["bitrate_kbps"], owl_stats["median_ms"],
            color="#4CAF50", marker="s", linewidth=2, markersize=7, zorder=3,
            label="Median one-way latency")
    ax.fill_between(owl_stats["bitrate_kbps"],
                    owl_stats["median_ms"] - owl_stats["std_ms"],
                    owl_stats["median_ms"] + owl_stats["std_ms"],
                    alpha=0.25, color="#4CAF50", label=r"$\pm1\sigma$")
    for _, row in owl_stats.iterrows():
        ax.annotate(f"{row['median_ms']:.1f} ms",
                    xy=(row["bitrate_kbps"], row["median_ms"]),
                    xytext=(0, 10), textcoords="offset points", ha="center", fontsize=8)
    ax.set_xscale("log")
    ax.set_xlabel("Bitrate (kbps)")
    ax.set_ylabel("One-Way Latency (ms)")
    ax.set_title(f"One-Way Latency vs Bitrate  |  Packet size = {PACKET_SIZE_B} B")
    ax.set_xticks(owl_stats["bitrate_kbps"])
    ax.set_xticklabels([f"{b:g}" for b in owl_stats["bitrate_kbps"]])
    ax.xaxis.set_minor_formatter(ticker.NullFormatter())
    ax.legend()
    ax.grid(True, which="both", linestyle="--", alpha=0.4)
    fig.tight_layout()
    fig.savefig(os.path.join(out_dir, "plot_owl_vs_bitrate.png"), dpi=150)
    plt.close(fig)

def plot_boxplot(df: pd.DataFrame, out_dir: str):
    stats = build_stats(df)
    bitrates = stats["bitrate_kbps"].tolist()
    groups = [df[df["bitrate_kbps"] == br]["RTT_us"].values / 1000 for br in bitrates]
    labels = [f"{b:g} kbps" for b in bitrates]
    fig, ax = plt.subplots(figsize=(9, 5))
    bp = ax.boxplot(groups, tick_labels=labels, patch_artist=True,
                    medianprops={"color": "black", "linewidth": 2})
    colors = ["#2196F3", "#4CAF50", "#FF9800", "#9C27B0", "#F44336"]
    for patch, color in zip(bp["boxes"], colors):
        patch.set_facecolor(color); patch.set_alpha(0.7)
    ax.set_title(f"RTT Distribution per Bitrate  |  Packet size = {PACKET_SIZE_B} B")
    ax.grid(axis="y", linestyle="--", alpha=0.4)
    fig.tight_layout()
    fig.savefig(os.path.join(out_dir, "plot_boxplot_bitrate.png"), dpi=150)
    plt.close(fig)

def plot_bar_comparison(df: pd.DataFrame, out_dir: str):
    stats = build_stats(df)
    bitrates = stats["bitrate_kbps"].tolist()
    x = np.arange(len(bitrates))
    owl_med = [df[df["bitrate_kbps"] == br]["OWL_ms"].median() for br in bitrates]
    owl_std = [df[df["bitrate_kbps"] == br]["OWL_ms"].std() for br in bitrates]
    fig, ax = plt.subplots(figsize=(9, 5))
    width = 0.35
    ax.bar(x - width/2, stats["median_ms"], width, yerr=stats["std_ms"],
           capsize=4, color="#2196F3", alpha=0.8, label="RTT")
    ax.bar(x + width/2, owl_med, width, yerr=owl_std,
           capsize=4, color="#4CAF50", alpha=0.8, label="One-Way Latency")
    ax.set_xticks(x)
    ax.set_xticklabels([f"{b:g} kbps" for b in bitrates])
    ax.set_ylabel("Latency (ms)")
    ax.set_title(f"RTT vs One-Way Latency per Bitrate  |  Packet size = {PACKET_SIZE_B} B")
    ax.legend()
    ax.grid(axis="y", linestyle="--", alpha=0.4)
    fig.tight_layout()
    fig.savefig(os.path.join(out_dir, "plot_bar_rtt_owl.png"), dpi=150)
    plt.close(fig)

if __name__ == "__main__":
    df = load_all_csvs(DATA_DIR)
    df = remove_outliers(df)
    os.makedirs(OUTPUT_DIR, exist_ok=True)
    plot_rtt_line(build_stats(df), OUTPUT_DIR)
    plot_owl_line(df, OUTPUT_DIR)
    plot_boxplot(df, OUTPUT_DIR)
    plot_bar_comparison(df, OUTPUT_DIR)
\end{lstlisting}

% ---- plot_latency.py --------------------------------------------------------
\begin{lstlisting}[style=appendix-py, caption={Latency vs.\ packet-size plotting}, label={lst:latency-pktsize-plot}]
"""
Latency comparison across packet sizes and bitrates.

Filename format: packet_size_<SIZE>B_<BITRATE>kb<timestamp>.csv
Columns: test_num, packet_size_B, RTT_us, one_way_us
"""

import os, re, glob
import pandas as pd
import numpy as np
import matplotlib.pyplot as plt
import matplotlib.ticker as ticker
from pathlib import Path

DATA_DIR = os.path.dirname(os.path.abspath(__file__))
OUTPUT_DIR = DATA_DIR

PACKET_SIZES = [5, 16, 32, 48, 64]
BITRATES = [9.6, 57.6, 115.2]

COLORS = {5: "#2196F3", 16: "#4CAF50", 32: "#FF9800", 48: "#9C27B0", 64: "#F44336"}
MARKERS = {9.6: "o", 57.6: "s", 115.2: "^"}
LINESTYLES = {9.6: "-", 57.6: "--", 115.2: ":"}

def parse_filename(path: str):
    name = Path(path).stem
    m = re.search(r'packet_size_(\d+)B_([\d._]+)kb', name)
    if m:
        return int(m.group(1)), float(m.group(2).replace("_", "."))
    return None

def load_all_csvs(data_dir: str):
    pattern = os.path.join(data_dir, "packet_size_*.csv")
    files = glob.glob(pattern)
    if not files:
        raise FileNotFoundError(f"No CSVs matching 'packet_size_*.csv' in {data_dir}")
    data = {}
    for f in sorted(files):
        key = parse_filename(f)
        if key is None:
            continue
        df = pd.read_csv(f)
        df.columns = ["test_num", "packet_size_B", "RTT_us", "one_way_us"]
        median_rtt = df["RTT_us"].median()
        df = df[df["RTT_us"] < 5 * median_rtt]
        data[key] = df
    return data

def summary_table(data: dict):
    rows = []
    for (ps, br), df in sorted(data.items()):
        rows.append({
            "Packet size (B)": ps, "Bitrate (kbps)": br,
            "RTT mean (ms)": df["RTT_us"].mean()/1000,
            "RTT median (ms)": df["RTT_us"].median()/1000,
            "RTT std (ms)": df["RTT_us"].std()/1000,
            "OWL median (ms)": df["one_way_us"].median()/1000,
            "OWL std (ms)": df["one_way_us"].std()/1000,
        })
    return pd.DataFrame(rows)

def plot_rtt_vs_packet_size(stats: pd.DataFrame, out_dir: str):
    fig, ax = plt.subplots(figsize=(8, 5))
    for br in sorted(stats["Bitrate (kbps)"].unique()):
        sub = stats[stats["Bitrate (kbps)"] == br].sort_values("Packet size (B)")
        ax.plot(sub["Packet size (B)"], sub["RTT median (ms)"],
                marker=MARKERS[br], linestyle=LINESTYLES[br],
                linewidth=1.8, markersize=7, label=f"{br} kbps")
        ax.fill_between(sub["Packet size (B)"],
                        sub["RTT median (ms)"] - sub["RTT std (ms)"],
                        sub["RTT median (ms)"] + sub["RTT std (ms)"], alpha=0.10)
    ax.set_xlabel("Packet Size (bytes)")
    ax.set_ylabel("Median RTT (ms)")
    ax.set_title("Round-Trip Time vs Packet Size")
    ax.set_xticks(PACKET_SIZES)
    ax.legend(title="Bitrate")
    ax.grid(True, linestyle="--", alpha=0.4)
    fig.tight_layout()
    fig.savefig(os.path.join(out_dir, "plot_rtt_vs_packet_size.png"), dpi=150)
    plt.close(fig)

def plot_owl_vs_packet_size(stats: pd.DataFrame, out_dir: str):
    fig, ax = plt.subplots(figsize=(8, 5))
    for br in sorted(stats["Bitrate (kbps)"].unique()):
        sub = stats[stats["Bitrate (kbps)"] == br].sort_values("Packet size (B)")
        ax.plot(sub["Packet size (B)"], sub["OWL median (ms)"],
                marker=MARKERS[br], linestyle=LINESTYLES[br],
                linewidth=1.8, markersize=7, label=f"{br} kbps")
        ax.fill_between(sub["Packet size (B)"],
                        sub["OWL median (ms)"] - sub["OWL std (ms)"],
                        sub["OWL median (ms)"] + sub["OWL std (ms)"], alpha=0.10)
    ax.set_xlabel("Packet Size (bytes)")
    ax.set_ylabel("Median One-Way Latency (ms)")
    ax.set_title("One-Way Latency vs Packet Size")
    ax.set_xticks(PACKET_SIZES)
    ax.legend(title="Bitrate")
    ax.grid(True, linestyle="--", alpha=0.4)
    fig.tight_layout()
    fig.savefig(os.path.join(out_dir, "plot_owl_vs_packet_size.png"), dpi=150)
    plt.close(fig)

def plot_grouped_bar(stats: pd.DataFrame, out_dir: str):
    bitrates = sorted(stats["Bitrate (kbps)"].unique())
    packet_sizes = sorted(stats["Packet size (B)"].unique())
    n_br, n_ps = len(bitrates), len(packet_sizes)
    x = np.arange(n_ps)
    width = 0.7 / n_br
    offsets = np.linspace(-(n_br-1)/2, (n_br-1)/2, n_br) * width
    fig, axes = plt.subplots(1, 2, figsize=(13, 5))
    for ax, metric, ylabel in [
        (axes[0], "RTT median (ms)", "Median RTT (ms)"),
        (axes[1], "OWL median (ms)", "Median One-Way Latency (ms)"),
    ]:
        for i, br in enumerate(bitrates):
            sub = stats[stats["Bitrate (kbps)"] == br].sort_values("Packet size (B)")
            vals = [sub[sub["Packet size (B)"] == ps][metric].values[0]
                    if len(sub[sub["Packet size (B)"] == ps]) else 0 for ps in packet_sizes]
            std_col = metric.replace("median", "std")
            errs = [sub[sub["Packet size (B)"] == ps][std_col].values[0]
                    if len(sub[sub["Packet size (B)"] == ps]) else 0 for ps in packet_sizes]
            ax.bar(x + offsets[i], vals, width, yerr=errs, label=f"{br} kbps", capsize=3)
        ax.set_xlabel("Packet Size (bytes)")
        ax.set_ylabel(ylabel)
        ax.set_xticks(x)
        ax.set_xticklabels([f"{ps} B" for ps in packet_sizes])
        ax.legend(title="Bitrate")
        ax.grid(axis="y", linestyle="--", alpha=0.4)
    fig.suptitle("Latency Comparison -- All Packet Sizes and Bitrates")
    fig.tight_layout()
    fig.savefig(os.path.join(out_dir, "plot_grouped_bar.png"), dpi=150)
    plt.close(fig)

def plot_boxplot(data: dict, out_dir: str):
    bitrates = sorted({k[1] for k in data})
    packet_sizes = sorted({k[0] for k in data})
    fig, axes = plt.subplots(1, len(bitrates), figsize=(5*len(bitrates), 5), sharey=True)
    if len(bitrates) == 1:
        axes = [axes]
    for ax, br in zip(axes, bitrates):
        groups = [data[(ps, br)]["RTT_us"].values/1000 for ps in packet_sizes if (ps, br) in data]
        labels = [f"{ps} B" for ps in packet_sizes if (ps, br) in data]
        bp = ax.boxplot(groups, labels=labels, patch_artist=True,
                        medianprops={"color": "black", "linewidth": 2})
        for patch, ps in zip(bp["boxes"], [ps for ps in packet_sizes if (ps, br) in data]):
            patch.set_facecolor(COLORS.get(ps, "#888")); patch.set_alpha(0.75)
        ax.set_title(f"{br} kbps")
        ax.grid(axis="y", linestyle="--", alpha=0.4)
    axes[0].set_ylabel("RTT (ms)")
    fig.suptitle("RTT Distribution -- Grouped by Bitrate")
    fig.tight_layout()
    fig.savefig(os.path.join(out_dir, "plot_boxplot_rtt.png"), dpi=150)
    plt.close(fig)

def plot_heatmap(stats: pd.DataFrame, out_dir: str):
    import matplotlib.colors as mcolors
    for metric, title, fname in [
        ("RTT median (ms)", "Median RTT (ms)", "plot_heatmap_rtt.png"),
        ("OWL median (ms)", "Median One-Way Latency (ms)", "plot_heatmap_owl.png"),
    ]:
        pivot = stats.pivot(index="Packet size (B)", columns="Bitrate (kbps)", values=metric)
        fig, ax = plt.subplots(figsize=(6, 4))
        im = ax.imshow(pivot.values, aspect="auto", cmap="YlOrRd",
                       norm=mcolors.Normalize(vmin=pivot.values.min(), vmax=pivot.values.max()))
        ax.set_xticks(range(len(pivot.columns)))
        ax.set_xticklabels([f"{c} kbps" for c in pivot.columns])
        ax.set_yticks(range(len(pivot.index)))
        ax.set_yticklabels([f"{r} B" for r in pivot.index])
        ax.set_title(f"Heatmap -- {title}")
        for i in range(len(pivot.index)):
            for j in range(len(pivot.columns)):
                val = pivot.values[i, j]
                if not np.isnan(val):
                    ax.text(j, i, f"{val:.2f}", ha="center", va="center", fontsize=9,
                            color="white" if val > pivot.values.max()*0.65 else "black")
        plt.colorbar(im, ax=ax, label=title)
        fig.tight_layout()
        fig.savefig(os.path.join(out_dir, fname), dpi=150)
        plt.close(fig)

if __name__ == "__main__":
    data = load_all_csvs(DATA_DIR)
    stats = summary_table(data)
    os.makedirs(OUTPUT_DIR, exist_ok=True)
    plot_rtt_vs_packet_size(stats, OUTPUT_DIR)
    plot_owl_vs_packet_size(stats, OUTPUT_DIR)
    plot_grouped_bar(stats, OUTPUT_DIR)
    plot_boxplot(data, OUTPUT_DIR)
    plot_heatmap(stats, OUTPUT_DIR)
\end{lstlisting}

\endinput

% in your report where you want the appendix to appear.
%
% If you already have these packages in your preamble, comment out the
% \usepackage lines below.
% =============================================================================

% ---- Packages (uncomment any NOT already in your preamble) ------------------
% \usepackage{xcolor}
% \usepackage{listings}
% \usepackage{caption}

% ---- Listing style -----------------------------------------------------------
\definecolor{codebg}{rgb}{0.95,0.95,0.97}
\definecolor{codeframe}{rgb}{0.80,0.80,0.85}
\definecolor{stringgreen}{rgb}{0.15,0.55,0.15}
\definecolor{keywordblue}{rgb}{0.10,0.20,0.70}
\definecolor{commentgray}{rgb}{0.35,0.35,0.35}

\lstdefinestyle{appendix}{
  language        = C++,
  basicstyle      = \footnotesize\ttfamily,
  backgroundcolor = \color{codebg},
  frame           = single,
  rulecolor       = \color{codeframe},
  numbers         = left,
  numberstyle     = \tiny\color{gray},
  stepnumber      = 1,
  numbersep       = 8pt,
  tabsize         = 2,
  showspaces      = false,
  showstringspaces = false,
  breaklines      = true,
  breakatwhitespace = true,
  keywordstyle    = \color{keywordblue}\bfseries,
  commentstyle    = \color{commentgray}\itshape,
  stringstyle     = \color{stringgreen},
  captionpos      = t,
  xleftmargin     = 2em,
  framexleftmargin = 1.5em,
}

\lstdefinestyle{appendix-py}{
  style           = appendix,
  language        = Python,
}

\lstdefinestyle{appendix-ini}{
  style           = appendix,
  language        = {[LaTeX]TeX},
}

% =============================================================================
\chapter{Source Code Listings}  \label{app:code}
% =============================================================================

% -----------------------------------------------------------------------------
\section{Remote Controller Firmware (C++)}
% -----------------------------------------------------------------------------

The handheld remote controller firmware runs on an Adafruit Feather~M0
with an RFM69~433\,MHz transceiver, an SSD1306~128$\times$32~OLED display,
and an analog dual-axis joystick.  PlatformIO is used as the build system.

% ---- main.cpp ---------------------------------------------------------------
\begin{lstlisting}[style=appendix, caption={Remote --- main.cpp}, label={lst:remote-main}]
#include <Arduino.h>
#include "OledDis.h"
#include "RadioTrans.h"
#include "Joystick.h"

#define BUTTON_PIN 10
#define VBATPIN A7
#define PWR_CONTROL 5

volatile uint8_t buttonState = 0;

void handleButton() {
  buttonState = !buttonState;
}

void setup() {
  Serial.begin(115200);
  delay(2000);

  pinMode(PWR_CONTROL, OUTPUT);
  digitalWrite(PWR_CONTROL, HIGH);


  attachInterrupt(digitalPinToInterrupt(BUTTON_PIN), handleButton, FALLING);

  Display_init();

  transiver_init();   

  Serial.println("System ready");
}

void loop() {
  uint8_t x, y;
  int8_t RSSI = 0;
  uint8_t battery = 0;
  char robotName[15] = "No connection";

  readJoystick(&x, &y);


  // Remote battery mesurements
  float measuredvbat = analogRead(VBATPIN);
  measuredvbat *= 2;
  measuredvbat *= 3.3;
  measuredvbat /= 1024.0;



  sendData(x, y, buttonState, &battery, robotName, &RSSI);
  Display_draw(robotName, measuredvbat, battery, RSSI);

  delay(100);
}
\end{lstlisting}

% ---- Joystick.cpp -----------------------------------------------------------
\begin{lstlisting}[style=appendix, caption={Remote --- Joystick.cpp}, label={lst:remote-joystick}]
#include <Arduino.h>


void readJoystick(uint8_t *x, uint8_t *y) {
    *x = -analogRead(A0) >> 2;
    *y = -analogRead(A1) >> 2;
}
\end{lstlisting}

% ---- OledDis.cpp ------------------------------------------------------------
\begin{lstlisting}[style=appendix, caption={Remote --- OledDis.cpp}, label={lst:remote-oled}]
#include <Arduino.h>
#include "OledDis.h"
#include <Wire.h>
#include <Adafruit_SSD1306.h>

#define SCREEN_WIDTH 128
#define SCREEN_HEIGHT 32

Adafruit_SSD1306 display(
  SCREEN_WIDTH,
  SCREEN_HEIGHT,
  &Wire,     
  -1
);

void Display_init() {
  Wire.begin();
  delay(100);
  if (!display.begin(SSD1306_SWITCHCAPVCC, 0x3C)) {
    Serial.println("OLED not found");
    while (1);
  }
  Serial.println("OLED found");

  display.clearDisplay();
  display.setRotation(2);
  display.ssd1306_command(SSD1306_DISPLAYON);
  display.ssd1306_command(SSD1306_NORMALDISPLAY);
  display.display();
}

void Display_test() {
  display.clearDisplay();

  display.drawRect(0, 0, 128, 32, SSD1306_WHITE);

  display.drawLine(0, 0, 127, 31, SSD1306_WHITE);

  display.display();
}

void Display_clear() {
  display.clearDisplay();
  display.display();
}

void Display_draw(char *name, float remoteBattery, uint8_t robotBattery, int8_t RSSI) {
  display.clearDisplay();
  display.setTextSize(1);
  display.setTextColor(SSD1306_WHITE);
  display.setCursor(0, 0);

  display.print("Name:");
  display.println(name);

  display.print("R:");
  display.print(remoteBattery, 1);
  display.print("V");

  display.print("B:");
  display.print(robotBattery / 10.0, 1);
  display.println("V");

  display.print("RSSI:");
  display.println(RSSI);

  display.display();
}
\end{lstlisting}

% ---- RadioTrans.cpp ---------------------------------------------------------
\begin{lstlisting}[style=appendix, caption={Remote --- RadioTrans.cpp}, label={lst:remote-radio}]
#include "RadioTrans.h"
#include <SPI.h>
#include <RFM69.h>
#include <RFM69registers.h>

#define RFM69_CS    0
#define RFM69_RST   1
#define RFM69_INT   A4
#define RFM69_DIO2  A3
#define RFM69_DIO3  A5

#define NETWORKID   42
#define NODEID      2
#define TONODEID    1
#define FREQUENCY   RF69_433MHZ
#define LATENCY_PACKET 0x42
#define PACKET_TEST 0x55

RFM69 radio(RFM69_CS, RFM69_INT);

void transiver_init()
{
    pinMode(RFM69_RST, OUTPUT);
    digitalWrite(RFM69_RST, HIGH);
    delay(50);
    digitalWrite(RFM69_RST, LOW);
    delay(50);

    SPI.begin();

    if (!radio.initialize(FREQUENCY, NODEID, NETWORKID))
    {
        Serial.println("RFM69 init FAILED");
        while (1);
    }

    radio.setHighPower();

    Serial.println("RFM69 init OK");
}

bool sendData(uint8_t x, uint8_t y, uint8_t button, uint8_t *battery, char *robotName, int8_t *RSSI)
{
    uint8_t packet[3];
    packet[0] = x;
    packet[1] = y;
    packet[2] = button;

    if (!radio.sendWithRetry(TONODEID, packet, sizeof(packet), 1, 50))
    {
        Serial.println("No ACK received");
        return false;
    }

    *battery = radio.DATA[0];

    for (int i = 0; i < 14; i++)
    {
        robotName[i] = (char)radio.DATA[i + 1];

        if (robotName[i] == '\0')
            break;
    }

    robotName[14] = '\0';

    *RSSI = radio.RSSI;

    return true;
}
\end{lstlisting}

% ---- Joystick.h -------------------------------------------------------------
\begin{lstlisting}[style=appendix, caption={Remote --- Joystick.h}, label={lst:remote-joystick-h}]
#include <Arduino.h>

void readJoystick(uint8_t *x, uint8_t *y);
\end{lstlisting}

% ---- OledDis.h --------------------------------------------------------------
\begin{lstlisting}[style=appendix, caption={Remote --- OledDis.h}, label={lst:remote-oled-h}]
#pragma once
#include <Arduino.h>

void Display_init();
void Display_test();
void Display_clear();
void Display_draw(char *name, float remoteBattery, uint8_t robotBattery, int8_t RSSI);
\end{lstlisting}

% ---- RadioTrans.h -----------------------------------------------------------
\begin{lstlisting}[style=appendix, caption={Remote --- RadioTrans.h}, label={lst:remote-radio-h}]
#ifndef RFM69_H
#define RFM69_H

#include <Arduino.h>

void transiver_init();

bool sendData(uint8_t x, uint8_t y, uint8_t button, uint8_t *battery, char *robotName, int8_t *RSSI);


#endif
\end{lstlisting}

% ---- platformio.ini ---------------------------------------------------------
\begin{lstlisting}[style=appendix-ini, caption={Remote --- platformio.ini}, label={lst:remote-pio}]
[env:adafruit_feather_m0]
platform = atmelsam
board = adafruit_feather_m0
framework = arduino
monitor_speed = 115200
build_flags = -D USB_SERIAL
lib_deps = 
	adafruit/Adafruit SSD1306@^2.5.16
	lowpowerlab/RFM69@^1.6.0
	lowpowerlab/SPIFlash@^101.1.3
\end{lstlisting}

% -----------------------------------------------------------------------------
\section{Robot-Side Firmware (C++)}
% -----------------------------------------------------------------------------

The robot-side firmware runs on a second Adafruit Feather~M0.  It receives
joystick data over the RFM69 link, reads battery voltage and a human-readable
name from a Raspberry Pi over the serial port, and echoes these back to the
remote inside the radio's ACK payload.

% ---- Robot main.cpp ---------------------------------------------------------
\begin{lstlisting}[style=appendix, caption={Robot --- main.cpp}, label={lst:robot-main}]
#include <Arduino.h>
#include <SPI.h>
#include <RFM69.h>
#include "RemoteTransiver.h"

void setup() {
  Serial.begin(115200);
  while (!Serial) delay(10);

  transiver_init();
}

void loop() {
  transiver_receive();
}
\end{lstlisting}

% ---- RemoteTransiver.cpp ----------------------------------------------------
\begin{lstlisting}[style=appendix, caption={Robot --- RemoteTransiver.cpp}, label={lst:robot-trans}]
#include <SPI.h>
#include <RFM69.h>
#include "RemoteTransiver.h"
#include <RFM69registers.h>

#define RFM69_CS   A5
#define RFM69_INT  6
#define RFM69_RST  A4
#define ESTOP      A3
#define LED_PIN    13


#define NETWORKID  42
#define NODEID     1
#define TONODEID   2
#define FREQUENCY  RF69_433MHZ
#define LATENCY_PACKET 0x42
#define PACKET_TEST 0x55

RFM69 radio(RFM69_CS, RFM69_INT);

float batteryVoltage = 0;       
char robotName[16] = "none/con";     
bool piDataReceived = false;     

unsigned long lastReceiveTime = 0;
unsigned long lastNoConnectionPrint = 0;
const unsigned long CONNECTION_TIMEOUT = 1000;


void transiver_init()
{
    pinMode(LED_PIN, OUTPUT);
    pinMode(RFM69_RST, OUTPUT);
    digitalWrite(RFM69_RST, HIGH);
    delay(100);
    digitalWrite(RFM69_RST, LOW);
    delay(100);

    pinMode(ESTOP, OUTPUT);

    if (!radio.initialize(FREQUENCY, NODEID, NETWORKID))
    {
        Serial.println("RFM69 init FAILED");
        while (1);
    }

    radio.setHighPower();

    Serial.println("RFM69 init OK");
    lastReceiveTime = millis();
}

void readPiSerial()
{
    static char line[64];
    static uint8_t idx = 0;

    while (Serial.available())
    {
        char c = Serial.read();
        if (c == '\r') continue;

        if (c == '\n')
        {
            line[idx] = '\0';
            idx = 0;

            // Find "name:" 
            char* namePtr = strstr(line, "name:");
            char* batPtr  = strstr(line, ",bat:");

            if (namePtr && batPtr)
            {
                namePtr += 5; // skip "name:"

                // Copy name up to the comma
                uint8_t i = 0;
                while (namePtr[i] != ',' && namePtr[i] != '\0' && i < 15)
                {
                    robotName[i] = namePtr[i];
                    i++;
                }
                robotName[i] = '\0';

                // Parse float manually
                batPtr += 5; // skip ",bat:"
                batteryVoltage = atof(batPtr);
                piDataReceived = true;


            }

        }
        else
        {
            if (idx < sizeof(line) - 1)
                line[idx++] = c;
            else
                idx = 0;
        }
    }
}


void handleDataPacket(uint8_t x, uint8_t y, uint8_t button)
{

    Serial.print("x:");
    Serial.print(x);
    Serial.print(", y:");
    Serial.print(y);
    Serial.print(", btn:");
    Serial.println(button);

    if (button == 1)
        digitalWrite(ESTOP, HIGH);
    else
        digitalWrite(ESTOP, LOW);
}


void transiver_receive()
{

    readPiSerial();

    if (radio.receiveDone())
    {
        if (radio.ACKRequested())
        {

            uint8_t ackPayload[16];

            // Battery voltage, e.g. 24.3V -> 243
            ackPayload[0] = (uint8_t)(batteryVoltage * 10.0);

            // Copy robot name into remaining bytes
            uint8_t i = 0;
            for (; i < 14 && robotName[i] != '\0'; i++)
            {
                ackPayload[i + 1] = robotName[i];
            }

            // Zero-fill rest
            for (; i < 14; i++)
            {
                ackPayload[i + 1] = 0;
            }

            radio.sendACK(ackPayload, 15);
        }

        uint8_t x      = radio.DATA[0];
        uint8_t y      = radio.DATA[1];
        uint8_t button = radio.DATA[2];

        handleDataPacket(x, y, button);
        lastReceiveTime = millis();
    }

    if (millis() - lastReceiveTime >= CONNECTION_TIMEOUT) {
        if (millis() - lastNoConnectionPrint >= 1000) {
            Serial.println("no connection"); 
            lastNoConnectionPrint = millis();
        }
    }
}
\end{lstlisting}

% ---- RemoteTransiver.h ------------------------------------------------------
\begin{lstlisting}[style=appendix, caption={Robot --- RemoteTransiver.h}, label={lst:robot-trans-h}]
#ifndef RFM69_H
#define RFM69_H

#include <Arduino.h>

void transiver_init();
void transiver_receive();
void readPiSerial();

#endif
\end{lstlisting}

% ---- Robot platformio.ini ---------------------------------------------------
\begin{lstlisting}[style=appendix-ini, caption={Robot --- platformio.ini}, label={lst:robot-pio}]
[env:adafruit_feather_m0]
platform = atmelsam
board = adafruit_feather_m0
framework = arduino
monitor_speed = 115200
build_flags = -D USB_SERIAL
lib_deps = 
	adafruit/Adafruit SSD1306@^2.5.16
	lowpowerlab/RFM69@^1.6.0
	lowpowerlab/SPIFlash@^101.1.3
\end{lstlisting}

% -----------------------------------------------------------------------------
\section{Test and Measurement Scripts (Python)}
% -----------------------------------------------------------------------------

The following Python scripts were used during the project to characterise
the radio link's latency, packet-loss, and signal-strength (RSSI)
performance.  All measurement scripts communicate with the robot-side
Feather~M0 over a serial port and use \texttt{matplotlib} to produce plots.

% ---- test_feather.py --------------------------------------------------------
\begin{lstlisting}[style=appendix-py, caption={Feather serial-communication test}, label={lst:test-feather}]
#!/usr/bin/env python3
"""
Feather M0 Communication Test Script
============================
Usage: python test_feather.py

This script tests serial communication with the Feather M0 (RobotSide).

Setup:
  - Windows: Uses COM8 (change below if different)
  - Linux: Uses /dev/ttyACM1 (change below if different)

Run:
  python test_feather.py

What it does:
  1. Sends "name:TESTBOT,bat:12.34\n" to Feather
  2. Waits for response
  3. Reports what was received

Output examples:
  - "Response: OK" = Feather received and responded
  - "No response" = Feather not receiving data
"""

import serial
import time
import sys

WINDOWS_PORT = 'COM8'
LINUX_PORT = '/dev/ttyACM1'

if sys.platform.startswith('win'):
    SERIAL_PORT = WINDOWS_PORT
elif sys.platform.startswith('linux'):
    SERIAL_PORT = LINUX_PORT
else:
    SERIAL_PORT = '/dev/ttyUSB0'  # Mac

BAUD_RATE = 115200
TIMEOUT = 2

def test_feather():
    print("=== Feather M0 Communication Test ===")
    print(f"Platform: {sys.platform}")
    print(f"Port: {SERIAL_PORT}")
    print(f"Baud: {BAUD_RATE}")
    print("-" * 40)

    try:
        print("Opening serial port...")
        ser = serial.Serial(SERIAL_PORT, BAUD_RATE, timeout=TIMEOUT)
        time.sleep(2)
        print("Port opened!")

        print("\nSending test data...")
        test_msg = "name:TESTBOT,bat:12.34\n"
        ser.write(test_msg.encode())
        print(f"Sent: {test_msg.strip()}")

        print("\nReading response...")
        received_any = False
        
        for i in range(3):
            time.sleep(0.5)
            if ser.in_waiting > 0:
                data = ser.read(ser.in_waiting)
                print(f"Response {i+1}: {data.decode('utf-8', errors='replace').strip()}")
                received_any = True
            else:
                print(f"Attempt {i+1}: No data available")

        if not received_any:
            print("\nNo response received from Feather!")
            print("This means Feather is not responding to the serial data.")

        ser.close()
        print("\n=== Test Complete ===")

    except serial.SerialException as e:
        print(f"\nERROR: {e}")
        print("\nTroubleshooting:")
        print("  - Make sure Feather is connected")
        print("  - Check correct COM port (Windows: COMx, Linux: /dev/ttyACMx)")
        print("  - Close other programs using the port")
        print("  - Try a different USB port")

if __name__ == "__main__":
    test_feather()
\end{lstlisting}

% ---- BitRateLatencyTest.py --------------------------------------------------
\begin{lstlisting}[style=appendix-py, caption={Latency vs.\ bitrate measurement}, label={lst:bitrate-latency}]
import serial
import time
import matplotlib.pyplot as plt
import csv
import re
from datetime import datetime

SERIAL_PORT = "COM8"
BAUD_RATE = 115200

BITRATE_LABEL = "1.2 kbps"   # Change this manually for each test
MEASURE_TIME = 10

timestamp = datetime.now().strftime("%Y%m%d_%H%M%S")

csv_filename = f"latency_{BITRATE_LABEL.replace(' ', '_')}_{timestamp}.csv"
plot_filename = f"latency_{BITRATE_LABEL.replace(' ', '_')}_{timestamp}.png"

ser = serial.Serial(SERIAL_PORT, BAUD_RATE, timeout=1)
time.sleep(2)

rtt_values = []
oneway_values = []

pattern = r"RTT:\s+(\d+)\s+us\s+\|\s+One-way latency:\s+([\d.]+)\s+ms"

print(f"\nRecording latency data for {BITRATE_LABEL}...")

start = time.time()

while time.time() - start < MEASURE_TIME:
    line = ser.readline().decode(errors='ignore').strip()
    match = re.search(pattern, line)

    if match:
        rtt = int(match.group(1))
        latency = float(match.group(2))
        rtt_values.append(rtt)
        oneway_values.append(latency)
        print(f"RTT: {rtt} us | Latency: {latency} ms")

ser.close()

print(f"\nCollected {len(rtt_values)} samples")

with open(csv_filename, mode='w', newline='') as file:
    writer = csv.writer(file)
    writer.writerow(["Sample", "Bitrate", "RTT_us", "OneWayLatency_ms"])
    for i in range(len(rtt_values)):
        writer.writerow([i + 1, BITRATE_LABEL, rtt_values[i], oneway_values[i]])

print(f"CSV saved as: {csv_filename}")

avg_latency = sum(oneway_values) / len(oneway_values)

samples = list(range(1, len(oneway_values) + 1))

plt.figure(figsize=(8,5))
plt.plot(samples, oneway_values, marker='o', label="One-way Latency")
plt.axhline(avg_latency, linestyle='--', label=f"Average = {avg_latency:.3f} ms")
plt.xlabel("Sample Number")
plt.ylabel("Latency (ms)")
plt.title(f"Latency Test - {BITRATE_LABEL}")
plt.grid(True)
plt.legend()
plt.tight_layout()
plt.savefig(plot_filename, dpi=300)
print(f"Plot saved as: {plot_filename}")
plt.show()
\end{lstlisting}

% ---- PackageSizeLatencyTest.py ----------------------------------------------
\begin{lstlisting}[style=appendix-py, caption={Latency vs.\ packet-size measurement}, label={lst:pkt-size-latency}]
import serial
import time
import matplotlib.pyplot as plt
import csv
import re
from datetime import datetime

SERIAL_PORT = "COM8"
BAUD_RATE = 115200

PACKET_SIZE = 5     # Change for each test
BIT_RATE = 115.2
MEASURE_TIME = 10

timestamp = datetime.now().strftime("%Y%m%d_%H%M%S")

csv_filename = f"packet_size_{PACKET_SIZE}B_{BIT_RATE}kb{timestamp}.csv"
plot_filename = f"packet_size_{PACKET_SIZE}B_{BIT_RATE}kb{timestamp}.png"

ser = serial.Serial(SERIAL_PORT, BAUD_RATE, timeout=1)
time.sleep(2)

rtt_values = []
latency_values = []

pattern = r"Size:\s+(\d+)\s+bytes\s+\|\s+RTT:\s+(\d+)\s+us\s+\|\s+One-way latency:\s+([\d.]+)\s+ms"

print(f"\nRecording latency data for packet size = {PACKET_SIZE} bytes")

start = time.time()

while time.time() - start < MEASURE_TIME:
    line = ser.readline().decode(errors='ignore').strip()
    match = re.search(pattern, line)

    if match:
        size = int(match.group(1))
        rtt = int(match.group(2))
        latency = float(match.group(3))
        rtt_values.append(rtt)
        latency_values.append(latency)
        print(f"Size: {size} B | RTT: {rtt} us | Latency: {latency} ms")

ser.close()

print(f"\nCollected {len(latency_values)} samples")

with open(csv_filename, mode='w', newline='') as file:
    writer = csv.writer(file)
    writer.writerow(["Sample", "PacketSize_Bytes", "RTT_us", "OneWayLatency_ms"])
    for i in range(len(latency_values)):
        writer.writerow([i + 1, PACKET_SIZE, rtt_values[i], latency_values[i]])

print(f"CSV saved as: {csv_filename}")

avg_latency = sum(latency_values) / len(latency_values)

samples = list(range(1, len(latency_values) + 1))

plt.figure(figsize=(8,5))
plt.plot(samples, latency_values, marker='o', label="One-way Latency")
plt.axhline(avg_latency, linestyle='--', label=f"Average = {avg_latency:.3f} ms")
plt.xlabel("Sample Number")
plt.ylabel("Latency (ms)")
plt.title(f"Latency vs Time - Packet Size {PACKET_SIZE} Bytes")
plt.grid(True)
plt.legend()
plt.tight_layout()
plt.savefig(plot_filename, dpi=300)
print(f"Plot saved as: {plot_filename}")
plt.show()
\end{lstlisting}

% ---- packagelossTest.py -----------------------------------------------------
\begin{lstlisting}[style=appendix-py, caption={Packet loss vs.\ distance}, label={lst:packet-loss}]
import serial
import time
import matplotlib.pyplot as plt
import csv
from datetime import datetime

SERIAL_PORT = "COM8"
BAUD_RATE = 115200

DISTANCES = list([20, 40, 50, 70])
EXPECTED_PACKETS = 100

timestamp = datetime.now().strftime("%Y%m%d_%H%M%S")

csv_filename = f"packet_loss_test_{timestamp}.csv"
plot_filename = f"packet_loss_plot_{timestamp}.png"

ser = serial.Serial(SERIAL_PORT, BAUD_RATE, timeout=1)
time.sleep(2)

results = {}

for dist in DISTANCES:
    print(f"\nMove robot to {dist} meters")
    input("Press ENTER when ready...")
    print("Waiting for packets...")

    received_packets = set()
    start_time = time.time()

    while time.time() - start_time < 10:
        line = ser.readline().decode(errors='ignore').strip()

        if "Received packet:" in line:
            try:
                packet_num = int(line.split(":")[1])
                received_packets.add(packet_num)
                print(f"{dist}m -> Packet {packet_num}")
            except:
                pass

    received_count = len(received_packets)
    lost_count = EXPECTED_PACKETS - received_count
    loss_percent = (lost_count / EXPECTED_PACKETS) * 100

    results[dist] = {
        "received": received_count,
        "lost": lost_count,
        "loss_percent": loss_percent
    }

    print("\n========== RESULT ==========")
    print(f"Distance: {dist} m")
    print(f"Received: {received_count}")
    print(f"Lost: {lost_count}")
    print(f"Packet Loss: {loss_percent:.1f}%")
    print("============================")

ser.close()

with open(csv_filename, mode='w', newline='') as file:
    writer = csv.writer(file)
    writer.writerow(["Distance_m", "Received_Packets", "Lost_Packets", "Packet_Loss_Percent"])
    for dist in DISTANCES:
        writer.writerow([dist, results[dist]["received"], results[dist]["lost"], results[dist]["loss_percent"]])

print(f"\nCSV saved as: {csv_filename}")

distances = [dist for dist in DISTANCES]
losses = [results[dist]["loss_percent"] for dist in DISTANCES]

plt.figure(figsize=(8, 5))
plt.plot(distances, losses, marker='o')
plt.xlabel("Distance (m)")
plt.ylabel("Packet Loss (%)")
plt.title("Packet Loss vs Distance")
plt.grid(True)
plt.tight_layout()
plt.savefig(plot_filename, dpi=300)
print(f"Plot saved as: {plot_filename}")
plt.show()
\end{lstlisting}

% ---- OrrientationsTest.py ---------------------------------------------------
\begin{lstlisting}[style=appendix-py, caption={RSSI vs.\ antenna orientation}, label={lst:orientation}]
import serial
import time
import matplotlib.pyplot as plt
import csv
from datetime import datetime

SERIAL_PORT = "COM8"
BAUD_RATE = 115200

ORIENTATIONS = ["Front", "Back", "Left", "Right"]
MEASURE_TIME = 10
ANTENNA = "Horizontal_16.5cm"

timestamp = datetime.now().strftime("%Y%m%d_%H%M%S")

csv_filename = f"rssi_orientation_data_{ANTENNA}_{timestamp}.csv"
plot_filename = f"rssi_orientation_plot_{ANTENNA}_{timestamp}.png"

ser = serial.Serial(SERIAL_PORT, BAUD_RATE, timeout=1)
time.sleep(2)

results = {}

def parse_line(line):
    try:
        parts = line.split(",")
        x = int(parts[0].split(":")[1])
        rssi = int(parts[1].split(":")[1])
        return x, rssi
    except:
        return None, None

for orientation in ORIENTATIONS:
    print(f"\nRotate robot to: {orientation}")
    print("Push joystick to start recording...")

    while True:
        line = ser.readline().decode().strip()
        x, rssi = parse_line(line)
        if x is not None and x > 240:
            print(f"Recording {orientation} orientation...")
            break

    rssi_values = []
    start = time.time()

    while time.time() - start < MEASURE_TIME:
        line = ser.readline().decode().strip()
        x, rssi = parse_line(line)
        if rssi is not None:
            rssi_values.append(rssi)

    results[orientation] = rssi_values
    print(f"Collected {len(rssi_values)} samples")

ser.close()

with open(csv_filename, mode='w', newline='') as file:
    writer = csv.writer(file)
    writer.writerow(["Orientation", "Sample_Number", "RSSI_dBm"])
    for orientation in ORIENTATIONS:
        for i, rssi in enumerate(results[orientation]):
            writer.writerow([orientation, i + 1, rssi])

print(f"\nCSV data saved as: {csv_filename}")

orientations = []
means = []
mins = []
maxs = []

for orientation in ORIENTATIONS:
    vals = results[orientation]
    if len(vals) == 0:
        continue
    orientations.append(orientation)
    means.append(sum(vals) / len(vals))
    mins.append(min(vals))
    maxs.append(max(vals))

plt.figure(figsize=(8, 5))
plt.plot(orientations, means, marker='o', label="Average RSSI")
plt.fill_between(range(len(orientations)), mins, maxs, alpha=0.2, label="Variation")
plt.xlabel("Orientation")
plt.ylabel("RSSI (dBm)")
plt.title("RSSI vs Orientation")
plt.grid(True)
plt.legend()
plt.tight_layout()
plt.savefig(plot_filename, dpi=300)
print(f"Plot saved as: {plot_filename}")
plt.show()
\end{lstlisting}

% ---- RSSItest.py ------------------------------------------------------------
\begin{lstlisting}[style=appendix-py, caption={RSSI vs.\ distance and antenna length}, label={lst:rssi-test}]
import serial
import time
import matplotlib.pyplot as plt
import csv
from datetime import datetime

SERIAL_PORT = "COM8"
BAUD_RATE = 115200

DISTANCES = [1, 5, 10, 20, 40]
MEASURE_TIME = 10
ANTENNA = "Horizontal_16.5cm"

timestamp = datetime.now().strftime("%Y%m%d_%H%M%S")

csv_filename = f"rssi_data_{ANTENNA}_{timestamp}.csv"
plot_filename = f"rssi_plot_{ANTENNA}_{timestamp}.png"

ser = serial.Serial(SERIAL_PORT, BAUD_RATE, timeout=1)
time.sleep(2)

results = {}

def parse_line(line):
    try:
        parts = line.split(",")
        x = int(parts[0].split(":")[1])
        rssi = int(parts[1].split(":")[1])
        return x, rssi
    except:
        return None, None

for dist in DISTANCES:
    print(f"\nMove to {dist} meters. Push joystick to start...")

    while True:
        line = ser.readline().decode().strip()
        x, rssi = parse_line(line)
        if x is not None and x > 240:
            print(f"Recording at {dist} m...")
            break

    rssi_values = []
    start = time.time()

    while time.time() - start < MEASURE_TIME:
        line = ser.readline().decode().strip()
        x, rssi = parse_line(line)
        if rssi is not None:
            rssi_values.append(rssi)

    results[dist] = rssi_values
    print(f"Collected {len(rssi_values)} samples")

ser.close()

with open(csv_filename, mode='w', newline='') as file:
    writer = csv.writer(file)
    writer.writerow(["Distance_m", "Sample_Number", "RSSI_dBm"])
    for dist in DISTANCES:
        for i, rssi in enumerate(results[dist]):
            writer.writerow([dist, i + 1, rssi])

print(f"\nCSV data saved as: {csv_filename}")

distances = []
means = []
mins = []
maxs = []

for d in DISTANCES:
    vals = results[d]
    if len(vals) == 0:
        continue
    distances.append(d)
    means.append(sum(vals) / len(vals))
    mins.append(min(vals))
    maxs.append(max(vals))

plt.figure()
plt.plot(distances, means, marker='o', label="Average RSSI")
plt.fill_between(distances, mins, maxs, alpha=0.2, label="Variation")
plt.xlabel("Distance (m)")
plt.ylabel("RSSI (dBm)")
plt.title("RSSI vs Distance")
plt.grid(True)
plt.legend()
plt.tight_layout()
plt.savefig(plot_filename, dpi=300)
print(f"Plot saved as: {plot_filename}")
plt.show()
\end{lstlisting}

% ---- plot_rssi_antenna.py ---------------------------------------------------
\begin{lstlisting}[style=appendix-py, caption={Multi-file RSSI plotting (antenna comparison)}, label={lst:rssi-plot}]
"""
RSSI vs Distance -- comparison across antenna lengths.

Filename format:  rssi_data_*_<LENGTH>cm_<timestamp>.csv
Columns: Distance_m, Sample_Number, RSSI_dBm
"""

import os
import re
import glob
import pandas as pd
import numpy as np
import matplotlib.pyplot as plt
import matplotlib.ticker as ticker
from pathlib import Path
from collections import defaultdict

DATA_DIR = os.path.dirname(os.path.abspath(__file__))
OUTPUT_DIR = DATA_DIR

LENGTH_ORDER = [16.5, 16.7, 16.9, 17.1, 17.3, 17.5, 17.8]

PALETTE = ["#1f77b4", "#ff7f0e", "#2ca02c", "#d62728",
           "#9467bd", "#8c564b", "#e377c2"]

def parse_length(path: str):
    m = re.search(r'_([\d]+[_.][\d]+)cm[_.]', Path(path).name)
    if m:
        return float(m.group(1).replace("_", "."))
    m2 = re.search(r'_(\d+)cm[_.]', Path(path).name)
    if m2:
        return float(m2.group(1))
    return None

def load_all(data_dir: str) -> pd.DataFrame:
    pattern = os.path.join(data_dir, "rssi_data_*.csv")
    files = glob.glob(pattern)
    if not files:
        raise FileNotFoundError(f"No files matching 'rssi_data_*.csv' in {data_dir}")

    by_length = defaultdict(list)
    for f in sorted(files):
        length = parse_length(f)
        if length is None:
            continue
        df = pd.read_csv(f)
        df.columns = [c.strip() for c in df.columns]
        df = df.rename(columns={c: c.lower().replace(" ", "_") for c in df.columns})
        df = df.rename(columns={"distance_m": "distance", "rssi_dbm": "rssi"})
        df["length_cm"] = length
        by_length[length].append(df)

    frames = []
    for length, dfs in sorted(by_length.items()):
        frames.append(pd.concat(dfs, ignore_index=True))
    return pd.concat(frames, ignore_index=True)

def build_stats(df: pd.DataFrame) -> pd.DataFrame:
    g = df.groupby(["length_cm", "distance"])["rssi"]
    stats = pd.DataFrame({
        "mean": g.mean(), "median": g.median(), "std": g.std(),
        "p5": g.quantile(0.05), "p95": g.quantile(0.95), "n": g.count(),
    }).reset_index()
    return stats

def plot_rssi_vs_distance(stats: pd.DataFrame, out_dir: str):
    lengths = sorted(stats["length_cm"].unique(),
                     key=lambda x: LENGTH_ORDER.index(x) if x in LENGTH_ORDER else x)
    fig, ax = plt.subplots(figsize=(9, 6))
    for i, length in enumerate(lengths):
        sub = stats[stats["length_cm"] == length].sort_values("distance")
        color = PALETTE[i % len(PALETTE)]
        marker = ["o", "s", "^", "D", "v", "P", "X"][i % 7]
        ax.plot(sub["distance"], sub["median"],
                color=color, marker=marker, linewidth=2, markersize=7,
                label=f"{length} cm", zorder=3)
        ax.fill_between(sub["distance"], sub["p5"], sub["p95"],
                        alpha=0.08, color=color)
        ax.fill_between(sub["distance"],
                        sub["median"] - sub["std"], sub["median"] + sub["std"],
                        alpha=0.15, color=color)
    ax.set_xscale("log")
    ax.set_xlabel("Distance (m)", fontsize=11)
    ax.set_ylabel("RSSI (dBm)", fontsize=11)
    ax.set_title("RSSI vs Distance -- Antenna Length Comparison", fontsize=12)
    distances = sorted(stats["distance"].unique())
    ax.set_xticks(distances)
    ax.set_xticklabels([f"{d} m" for d in distances])
    ax.xaxis.set_minor_formatter(ticker.NullFormatter())
    ax.legend(title="Antenna length", loc="lower left")
    ax.grid(True, which="both", linestyle="--", alpha=0.4)
    fig.tight_layout()
    fig.savefig(os.path.join(out_dir, "plot_rssi_vs_distance.png"), dpi=150)
    plt.close(fig)

def plot_bar_per_distance(stats: pd.DataFrame, out_dir: str):
    lengths = sorted(stats["length_cm"].unique(),
                     key=lambda x: LENGTH_ORDER.index(x) if x in LENGTH_ORDER else x)
    distances = sorted(stats["distance"].unique())
    n_len = len(lengths)
    x = np.arange(len(distances))
    width = 0.7 / n_len
    offsets = np.linspace(-(n_len - 1) / 2, (n_len - 1) / 2, n_len) * width
    fig, ax = plt.subplots(figsize=(11, 6))
    for i, length in enumerate(lengths):
        sub = stats[stats["length_cm"] == length].sort_values("distance")
        vals = [sub[sub["distance"] == d]["median"].values[0]
                if len(sub[sub["distance"] == d]) else np.nan for d in distances]
        errs = [sub[sub["distance"] == d]["std"].values[0]
                if len(sub[sub["distance"] == d]) else 0 for d in distances]
        ax.bar(x + offsets[i], vals, width, yerr=errs,
               color=PALETTE[i], alpha=0.8, label=f"{length} cm", capsize=3)
    ax.set_xticks(x)
    ax.set_xticklabels([f"{d} m" for d in distances])
    ax.set_ylabel("Median RSSI (dBm)")
    ax.set_title("Median RSSI per Distance -- Antenna Length Comparison")
    ax.legend(loc="lower left")
    ax.grid(axis="y", linestyle="--", alpha=0.4)
    fig.tight_layout()
    fig.savefig(os.path.join(out_dir, "plot_rssi_bar_per_distance.png"), dpi=150)
    plt.close(fig)

def plot_boxplot_far(df: pd.DataFrame, out_dir: str):
    max_dist = df["distance"].max()
    sub = df[df["distance"] == max_dist]
    lengths = sorted(sub["length_cm"].unique(),
                     key=lambda x: LENGTH_ORDER.index(x) if x in LENGTH_ORDER else x)
    groups = [sub[sub["length_cm"] == l]["rssi"].values for l in lengths]
    labels = [f"{l} cm" for l in lengths]
    colors = PALETTE[:len(lengths)]
    fig, ax = plt.subplots(figsize=(9, 5))
    bp = ax.boxplot(groups, tick_labels=labels, patch_artist=True,
                    medianprops={"color": "black", "linewidth": 2})
    for patch, color in zip(bp["boxes"], colors):
        patch.set_facecolor(color)
        patch.set_alpha(0.75)
    ax.set_title(f"RSSI Distribution at {max_dist} m")
    ax.grid(axis="y", linestyle="--", alpha=0.4)
    fig.tight_layout()
    fig.savefig(os.path.join(out_dir, f"plot_rssi_boxplot_{max_dist}m.png"), dpi=150)
    plt.close(fig)

if __name__ == "__main__":
    df = load_all(DATA_DIR)
    stats = build_stats(df)
    os.makedirs(OUTPUT_DIR, exist_ok=True)
    plot_rssi_vs_distance(stats, OUTPUT_DIR)
    plot_bar_per_distance(stats, OUTPUT_DIR)
    plot_boxplot_far(df, OUTPUT_DIR)
\end{lstlisting}

% ---- plot_latency_vs_bitrate.py ---------------------------------------------
\begin{lstlisting}[style=appendix-py, caption={Latency vs.\ bitrate plotting}, label={lst:latency-bitrate-plot}]
"""
Latency vs Bitrate -- fixed packet size (5 B).

Loads all CSVs matching: latency_*kbps*.csv
Columns: Sample, Bitrate (string), RTT_us, OneWayLatency_ms
"""

import os, re, glob
import pandas as pd
import numpy as np
import matplotlib.pyplot as plt
import matplotlib.ticker as ticker
from pathlib import Path

DATA_DIR = os.path.dirname(os.path.abspath(__file__))
OUTPUT_DIR = DATA_DIR
PACKET_SIZE_B = 5
BITRATE_ORDER = [1.2, 9.6, 57.6, 115.2, 250.0]

def load_all_csvs(data_dir: str) -> pd.DataFrame:
    pattern = os.path.join(data_dir, "latency_*kbps*.csv")
    files = glob.glob(pattern)
    if not files:
        raise FileNotFoundError(f"No CSVs matching 'latency_*kbps*.csv' in {data_dir}")
    frames = []
    for f in sorted(files):
        df = pd.read_csv(f)
        df.columns = [c.strip() for c in df.columns]
        col_map = {c: ["sample", "bitrate_str", "RTT_us", "OWL_ms"][i]
                   for i, c in enumerate(df.columns)}
        df = df.rename(columns=col_map)
        frames.append(df)
    combined = pd.concat(frames, ignore_index=True)
    combined["bitrate_kbps"] = combined["bitrate_str"].str.extract(r"([\d.]+)").astype(float)
    return combined

def remove_outliers(df: pd.DataFrame) -> pd.DataFrame:
    clean = []
    for br, g in df.groupby("bitrate_kbps"):
        med = g["RTT_us"].median()
        clean.append(g[g["RTT_us"] < 5 * med])
    return pd.concat(clean, ignore_index=True)

def build_stats(df: pd.DataFrame) -> pd.DataFrame:
    g = df.groupby("bitrate_kbps")["RTT_us"]
    stats = pd.DataFrame({
        "n": g.count(), "mean_ms": g.mean()/1000, "median_ms": g.median()/1000,
        "std_ms": g.std()/1000, "p5_ms": g.quantile(0.05)/1000,
        "p95_ms": g.quantile(0.95)/1000,
    }).reset_index()
    order = {b: i for i, b in enumerate(BITRATE_ORDER)}
    stats["_order"] = stats["bitrate_kbps"].map(order).fillna(999)
    return stats.sort_values("_order").drop(columns="_order").reset_index(drop=True)

def plot_rtt_line(stats: pd.DataFrame, out_dir: str):
    fig, ax = plt.subplots(figsize=(8, 5))
    ax.plot(stats["bitrate_kbps"], stats["median_ms"],
            color="#2196F3", marker="o", linewidth=2, markersize=7, zorder=3,
            label="Median RTT")
    ax.fill_between(stats["bitrate_kbps"], stats["p5_ms"], stats["p95_ms"],
                    alpha=0.15, color="#2196F3", label="5th-95th percentile")
    ax.fill_between(stats["bitrate_kbps"],
                    stats["median_ms"] - stats["std_ms"],
                    stats["median_ms"] + stats["std_ms"],
                    alpha=0.25, color="#2196F3", label=r"$\pm1\sigma$")
    for _, row in stats.iterrows():
        ax.annotate(f"{row['median_ms']:.1f} ms",
                    xy=(row["bitrate_kbps"], row["median_ms"]),
                    xytext=(0, 10), textcoords="offset points", ha="center", fontsize=8)
    ax.set_xscale("log")
    ax.set_xlabel("Bitrate (kbps)")
    ax.set_ylabel("RTT (ms)")
    ax.set_title(f"Round-Trip Time vs Bitrate  |  Packet size = {PACKET_SIZE_B} B")
    ax.set_xticks(stats["bitrate_kbps"])
    ax.set_xticklabels([f"{b:g}" for b in stats["bitrate_kbps"]])
    ax.xaxis.set_minor_formatter(ticker.NullFormatter())
    ax.legend()
    ax.grid(True, which="both", linestyle="--", alpha=0.4)
    fig.tight_layout()
    fig.savefig(os.path.join(out_dir, "plot_rtt_vs_bitrate.png"), dpi=150)
    plt.close(fig)

def plot_owl_line(df: pd.DataFrame, out_dir: str):
    owl = df.groupby("bitrate_kbps")["OWL_ms"]
    owl_stats = pd.DataFrame({
        "median_ms": owl.median(), "std_ms": owl.std(),
        "p5_ms": owl.quantile(0.05), "p95_ms": owl.quantile(0.95),
    }).reset_index()
    order = {b: i for i, b in enumerate(BITRATE_ORDER)}
    owl_stats["_order"] = owl_stats["bitrate_kbps"].map(order).fillna(999)
    owl_stats = owl_stats.sort_values("_order").drop(columns="_order").reset_index(drop=True)
    fig, ax = plt.subplots(figsize=(8, 5))
    ax.plot(owl_stats["bitrate_kbps"], owl_stats["median_ms"],
            color="#4CAF50", marker="s", linewidth=2, markersize=7, zorder=3,
            label="Median one-way latency")
    ax.fill_between(owl_stats["bitrate_kbps"],
                    owl_stats["median_ms"] - owl_stats["std_ms"],
                    owl_stats["median_ms"] + owl_stats["std_ms"],
                    alpha=0.25, color="#4CAF50", label=r"$\pm1\sigma$")
    for _, row in owl_stats.iterrows():
        ax.annotate(f"{row['median_ms']:.1f} ms",
                    xy=(row["bitrate_kbps"], row["median_ms"]),
                    xytext=(0, 10), textcoords="offset points", ha="center", fontsize=8)
    ax.set_xscale("log")
    ax.set_xlabel("Bitrate (kbps)")
    ax.set_ylabel("One-Way Latency (ms)")
    ax.set_title(f"One-Way Latency vs Bitrate  |  Packet size = {PACKET_SIZE_B} B")
    ax.set_xticks(owl_stats["bitrate_kbps"])
    ax.set_xticklabels([f"{b:g}" for b in owl_stats["bitrate_kbps"]])
    ax.xaxis.set_minor_formatter(ticker.NullFormatter())
    ax.legend()
    ax.grid(True, which="both", linestyle="--", alpha=0.4)
    fig.tight_layout()
    fig.savefig(os.path.join(out_dir, "plot_owl_vs_bitrate.png"), dpi=150)
    plt.close(fig)

def plot_boxplot(df: pd.DataFrame, out_dir: str):
    stats = build_stats(df)
    bitrates = stats["bitrate_kbps"].tolist()
    groups = [df[df["bitrate_kbps"] == br]["RTT_us"].values / 1000 for br in bitrates]
    labels = [f"{b:g} kbps" for b in bitrates]
    fig, ax = plt.subplots(figsize=(9, 5))
    bp = ax.boxplot(groups, tick_labels=labels, patch_artist=True,
                    medianprops={"color": "black", "linewidth": 2})
    colors = ["#2196F3", "#4CAF50", "#FF9800", "#9C27B0", "#F44336"]
    for patch, color in zip(bp["boxes"], colors):
        patch.set_facecolor(color); patch.set_alpha(0.7)
    ax.set_title(f"RTT Distribution per Bitrate  |  Packet size = {PACKET_SIZE_B} B")
    ax.grid(axis="y", linestyle="--", alpha=0.4)
    fig.tight_layout()
    fig.savefig(os.path.join(out_dir, "plot_boxplot_bitrate.png"), dpi=150)
    plt.close(fig)

def plot_bar_comparison(df: pd.DataFrame, out_dir: str):
    stats = build_stats(df)
    bitrates = stats["bitrate_kbps"].tolist()
    x = np.arange(len(bitrates))
    owl_med = [df[df["bitrate_kbps"] == br]["OWL_ms"].median() for br in bitrates]
    owl_std = [df[df["bitrate_kbps"] == br]["OWL_ms"].std() for br in bitrates]
    fig, ax = plt.subplots(figsize=(9, 5))
    width = 0.35
    ax.bar(x - width/2, stats["median_ms"], width, yerr=stats["std_ms"],
           capsize=4, color="#2196F3", alpha=0.8, label="RTT")
    ax.bar(x + width/2, owl_med, width, yerr=owl_std,
           capsize=4, color="#4CAF50", alpha=0.8, label="One-Way Latency")
    ax.set_xticks(x)
    ax.set_xticklabels([f"{b:g} kbps" for b in bitrates])
    ax.set_ylabel("Latency (ms)")
    ax.set_title(f"RTT vs One-Way Latency per Bitrate  |  Packet size = {PACKET_SIZE_B} B")
    ax.legend()
    ax.grid(axis="y", linestyle="--", alpha=0.4)
    fig.tight_layout()
    fig.savefig(os.path.join(out_dir, "plot_bar_rtt_owl.png"), dpi=150)
    plt.close(fig)

if __name__ == "__main__":
    df = load_all_csvs(DATA_DIR)
    df = remove_outliers(df)
    os.makedirs(OUTPUT_DIR, exist_ok=True)
    plot_rtt_line(build_stats(df), OUTPUT_DIR)
    plot_owl_line(df, OUTPUT_DIR)
    plot_boxplot(df, OUTPUT_DIR)
    plot_bar_comparison(df, OUTPUT_DIR)
\end{lstlisting}

% ---- plot_latency.py --------------------------------------------------------
\begin{lstlisting}[style=appendix-py, caption={Latency vs.\ packet-size plotting}, label={lst:latency-pktsize-plot}]
"""
Latency comparison across packet sizes and bitrates.

Filename format: packet_size_<SIZE>B_<BITRATE>kb<timestamp>.csv
Columns: test_num, packet_size_B, RTT_us, one_way_us
"""

import os, re, glob
import pandas as pd
import numpy as np
import matplotlib.pyplot as plt
import matplotlib.ticker as ticker
from pathlib import Path

DATA_DIR = os.path.dirname(os.path.abspath(__file__))
OUTPUT_DIR = DATA_DIR

PACKET_SIZES = [5, 16, 32, 48, 64]
BITRATES = [9.6, 57.6, 115.2]

COLORS = {5: "#2196F3", 16: "#4CAF50", 32: "#FF9800", 48: "#9C27B0", 64: "#F44336"}
MARKERS = {9.6: "o", 57.6: "s", 115.2: "^"}
LINESTYLES = {9.6: "-", 57.6: "--", 115.2: ":"}

def parse_filename(path: str):
    name = Path(path).stem
    m = re.search(r'packet_size_(\d+)B_([\d._]+)kb', name)
    if m:
        return int(m.group(1)), float(m.group(2).replace("_", "."))
    return None

def load_all_csvs(data_dir: str):
    pattern = os.path.join(data_dir, "packet_size_*.csv")
    files = glob.glob(pattern)
    if not files:
        raise FileNotFoundError(f"No CSVs matching 'packet_size_*.csv' in {data_dir}")
    data = {}
    for f in sorted(files):
        key = parse_filename(f)
        if key is None:
            continue
        df = pd.read_csv(f)
        df.columns = ["test_num", "packet_size_B", "RTT_us", "one_way_us"]
        median_rtt = df["RTT_us"].median()
        df = df[df["RTT_us"] < 5 * median_rtt]
        data[key] = df
    return data

def summary_table(data: dict):
    rows = []
    for (ps, br), df in sorted(data.items()):
        rows.append({
            "Packet size (B)": ps, "Bitrate (kbps)": br,
            "RTT mean (ms)": df["RTT_us"].mean()/1000,
            "RTT median (ms)": df["RTT_us"].median()/1000,
            "RTT std (ms)": df["RTT_us"].std()/1000,
            "OWL median (ms)": df["one_way_us"].median()/1000,
            "OWL std (ms)": df["one_way_us"].std()/1000,
        })
    return pd.DataFrame(rows)

def plot_rtt_vs_packet_size(stats: pd.DataFrame, out_dir: str):
    fig, ax = plt.subplots(figsize=(8, 5))
    for br in sorted(stats["Bitrate (kbps)"].unique()):
        sub = stats[stats["Bitrate (kbps)"] == br].sort_values("Packet size (B)")
        ax.plot(sub["Packet size (B)"], sub["RTT median (ms)"],
                marker=MARKERS[br], linestyle=LINESTYLES[br],
                linewidth=1.8, markersize=7, label=f"{br} kbps")
        ax.fill_between(sub["Packet size (B)"],
                        sub["RTT median (ms)"] - sub["RTT std (ms)"],
                        sub["RTT median (ms)"] + sub["RTT std (ms)"], alpha=0.10)
    ax.set_xlabel("Packet Size (bytes)")
    ax.set_ylabel("Median RTT (ms)")
    ax.set_title("Round-Trip Time vs Packet Size")
    ax.set_xticks(PACKET_SIZES)
    ax.legend(title="Bitrate")
    ax.grid(True, linestyle="--", alpha=0.4)
    fig.tight_layout()
    fig.savefig(os.path.join(out_dir, "plot_rtt_vs_packet_size.png"), dpi=150)
    plt.close(fig)

def plot_owl_vs_packet_size(stats: pd.DataFrame, out_dir: str):
    fig, ax = plt.subplots(figsize=(8, 5))
    for br in sorted(stats["Bitrate (kbps)"].unique()):
        sub = stats[stats["Bitrate (kbps)"] == br].sort_values("Packet size (B)")
        ax.plot(sub["Packet size (B)"], sub["OWL median (ms)"],
                marker=MARKERS[br], linestyle=LINESTYLES[br],
                linewidth=1.8, markersize=7, label=f"{br} kbps")
        ax.fill_between(sub["Packet size (B)"],
                        sub["OWL median (ms)"] - sub["OWL std (ms)"],
                        sub["OWL median (ms)"] + sub["OWL std (ms)"], alpha=0.10)
    ax.set_xlabel("Packet Size (bytes)")
    ax.set_ylabel("Median One-Way Latency (ms)")
    ax.set_title("One-Way Latency vs Packet Size")
    ax.set_xticks(PACKET_SIZES)
    ax.legend(title="Bitrate")
    ax.grid(True, linestyle="--", alpha=0.4)
    fig.tight_layout()
    fig.savefig(os.path.join(out_dir, "plot_owl_vs_packet_size.png"), dpi=150)
    plt.close(fig)

def plot_grouped_bar(stats: pd.DataFrame, out_dir: str):
    bitrates = sorted(stats["Bitrate (kbps)"].unique())
    packet_sizes = sorted(stats["Packet size (B)"].unique())
    n_br, n_ps = len(bitrates), len(packet_sizes)
    x = np.arange(n_ps)
    width = 0.7 / n_br
    offsets = np.linspace(-(n_br-1)/2, (n_br-1)/2, n_br) * width
    fig, axes = plt.subplots(1, 2, figsize=(13, 5))
    for ax, metric, ylabel in [
        (axes[0], "RTT median (ms)", "Median RTT (ms)"),
        (axes[1], "OWL median (ms)", "Median One-Way Latency (ms)"),
    ]:
        for i, br in enumerate(bitrates):
            sub = stats[stats["Bitrate (kbps)"] == br].sort_values("Packet size (B)")
            vals = [sub[sub["Packet size (B)"] == ps][metric].values[0]
                    if len(sub[sub["Packet size (B)"] == ps]) else 0 for ps in packet_sizes]
            std_col = metric.replace("median", "std")
            errs = [sub[sub["Packet size (B)"] == ps][std_col].values[0]
                    if len(sub[sub["Packet size (B)"] == ps]) else 0 for ps in packet_sizes]
            ax.bar(x + offsets[i], vals, width, yerr=errs, label=f"{br} kbps", capsize=3)
        ax.set_xlabel("Packet Size (bytes)")
        ax.set_ylabel(ylabel)
        ax.set_xticks(x)
        ax.set_xticklabels([f"{ps} B" for ps in packet_sizes])
        ax.legend(title="Bitrate")
        ax.grid(axis="y", linestyle="--", alpha=0.4)
    fig.suptitle("Latency Comparison -- All Packet Sizes and Bitrates")
    fig.tight_layout()
    fig.savefig(os.path.join(out_dir, "plot_grouped_bar.png"), dpi=150)
    plt.close(fig)

def plot_boxplot(data: dict, out_dir: str):
    bitrates = sorted({k[1] for k in data})
    packet_sizes = sorted({k[0] for k in data})
    fig, axes = plt.subplots(1, len(bitrates), figsize=(5*len(bitrates), 5), sharey=True)
    if len(bitrates) == 1:
        axes = [axes]
    for ax, br in zip(axes, bitrates):
        groups = [data[(ps, br)]["RTT_us"].values/1000 for ps in packet_sizes if (ps, br) in data]
        labels = [f"{ps} B" for ps in packet_sizes if (ps, br) in data]
        bp = ax.boxplot(groups, labels=labels, patch_artist=True,
                        medianprops={"color": "black", "linewidth": 2})
        for patch, ps in zip(bp["boxes"], [ps for ps in packet_sizes if (ps, br) in data]):
            patch.set_facecolor(COLORS.get(ps, "#888")); patch.set_alpha(0.75)
        ax.set_title(f"{br} kbps")
        ax.grid(axis="y", linestyle="--", alpha=0.4)
    axes[0].set_ylabel("RTT (ms)")
    fig.suptitle("RTT Distribution -- Grouped by Bitrate")
    fig.tight_layout()
    fig.savefig(os.path.join(out_dir, "plot_boxplot_rtt.png"), dpi=150)
    plt.close(fig)

def plot_heatmap(stats: pd.DataFrame, out_dir: str):
    import matplotlib.colors as mcolors
    for metric, title, fname in [
        ("RTT median (ms)", "Median RTT (ms)", "plot_heatmap_rtt.png"),
        ("OWL median (ms)", "Median One-Way Latency (ms)", "plot_heatmap_owl.png"),
    ]:
        pivot = stats.pivot(index="Packet size (B)", columns="Bitrate (kbps)", values=metric)
        fig, ax = plt.subplots(figsize=(6, 4))
        im = ax.imshow(pivot.values, aspect="auto", cmap="YlOrRd",
                       norm=mcolors.Normalize(vmin=pivot.values.min(), vmax=pivot.values.max()))
        ax.set_xticks(range(len(pivot.columns)))
        ax.set_xticklabels([f"{c} kbps" for c in pivot.columns])
        ax.set_yticks(range(len(pivot.index)))
        ax.set_yticklabels([f"{r} B" for r in pivot.index])
        ax.set_title(f"Heatmap -- {title}")
        for i in range(len(pivot.index)):
            for j in range(len(pivot.columns)):
                val = pivot.values[i, j]
                if not np.isnan(val):
                    ax.text(j, i, f"{val:.2f}", ha="center", va="center", fontsize=9,
                            color="white" if val > pivot.values.max()*0.65 else "black")
        plt.colorbar(im, ax=ax, label=title)
        fig.tight_layout()
        fig.savefig(os.path.join(out_dir, fname), dpi=150)
        plt.close(fig)

if __name__ == "__main__":
    data = load_all_csvs(DATA_DIR)
    stats = summary_table(data)
    os.makedirs(OUTPUT_DIR, exist_ok=True)
    plot_rtt_vs_packet_size(stats, OUTPUT_DIR)
    plot_owl_vs_packet_size(stats, OUTPUT_DIR)
    plot_grouped_bar(stats, OUTPUT_DIR)
    plot_boxplot(data, OUTPUT_DIR)
    plot_heatmap(stats, OUTPUT_DIR)
\end{lstlisting}

\endinput
